% Equations de droites — Mathématiques 3ème — Cours signé OnBuch+
% Programme officiel MINESEC. Compiler avec : tectonic N10-equations-droites.tex
\newcommand{\DOCMATIERE}{Mathématiques}
\newcommand{\DOCNIVEAU}{3ème}
\newcommand{\DOCMODULE}{Mathématiques – classe de 3ème}
\newcommand{\DOCLECON}{N10}
\newcommand{\DOCTITRE}{Equations de droites}
% OnBuch+ — préambule « cours signés » ENRICHI (compilé avec tectonic / XeLaTeX)
% Système visuel « Pop » d'OnBuch+ : fond crème (jamais blanc), bordures encre
% épaisses, ombres « dures » décalées, titres Archivo Black en capitales.
% Version MATHÉMATIQUES : graphes (pgfplots/tikz), boîtes pédagogiques
% (définition, propriété, méthode, attention, le savais-tu, exercice résolu,
% exercice, TP, bilan), page de garde et signature OnBuch+ ancrée en bas.
%
% Macros attendues dans main.tex AVANT \input{preamble} :
%   \DOCMATIERE  \DOCNIVEAU  \DOCMODULE  \DOCLECON (numéro)  \DOCTITRE
\documentclass[11pt]{article}

\usepackage{fontspec}
\usepackage[french]{babel}
\usepackage[a4paper,margin=2cm,top=3.4cm,bottom=2.8cm,headheight=1.4cm,footskip=1.3cm]{geometry}
\usepackage{amsmath,amssymb,amsfonts}
\usepackage{mathtools} % \xmapsto, \coloneqq... souvent utilisés par les modèles
\usepackage{cancel} % simplifications barrées
% \abs{z} (module, valeur absolue) et \norm{u} (norme), utilisés par les modèles
\providecommand{\abs}{}\let\abs\relax\DeclarePairedDelimiter{\abs}{\lvert}{\rvert}
\providecommand{\norm}{}\let\norm\relax\DeclarePairedDelimiter{\norm}{\lVert}{\rVert}
\usepackage[table]{xcolor} % [table] : fournit aussi \rowcolors (alternance de lignes)
\usepackage{pagecolor}
\usepackage{tikz}
\usetikzlibrary{arrows.meta,positioning,calc,shapes.geometric,shapes.misc,shapes.arrows,decorations.pathmorphing,decorations.pathreplacing,decorations.markings,patterns,shadows,mindmap,trees,fit,backgrounds,matrix,intersections,angles,quotes}
\usepackage{pgfplots}
\usepackage{pgf-pie} % diagrammes circulaires \pie (statistiques)
\pgfplotsset{compat=1.18}
\usepgfplotslibrary{fillbetween,groupplots} % aires entre courbes (calcul intégral)
\usepackage{enumitem}
\usepackage{fancyhdr}
% [most] charge toutes les bibliothèques tcolorbox (listings, minted, raster,
% documentation...) et gonfle inutilement la mémoire TeX sur un document long
% (~300 boîtes) : on ne charge que ce qui est réellement utilisé.
\usepackage[skins,breakable]{tcolorbox}
\usepackage{graphicx}
\usepackage{colortbl}
\usepackage{booktabs}
\usepackage{tabularx}
\usepackage{diagbox} % case d'en-tête diagonale des tableaux à double entrée
% Colonnes extensibles centrée / gauche / droite (C, L, R), souvent utilisées par les modèles
\newcolumntype{C}{>{\centering\arraybackslash}X}
\newcolumntype{L}{>{\raggedright\arraybackslash}X}
\newcolumntype{R}{>{\raggedleft\arraybackslash}X}
\usepackage{array}
\usepackage{multicol}
\usepackage{siunitx}
% parse-numbers=false : accepte aussi \SI{2,5 \times 10^{-3}}{...}
\sisetup{locale=FR,output-decimal-marker={,},group-minimum-digits=4,parse-numbers=false}
\DeclareSIUnit\ohm{\ensuremath{\symup{\Omega}}} % U+2126 absent de la police
\usepackage{qrcode}
% \degree : commande fréquemment utilisée par les modèles pour un angle ou une
% température en dehors de \SI{}{} (non fournie par les paquets chargés).
\providecommand{\degree}{^{\circ}}
% \qed (fin de démonstration), utilisé par les modèles sans amsthm
\providecommand{\qed}{\hfill$\square$}
% \barre : double barre des valeurs interdites dans un tableau de variations (tkz-tab)
\providecommand{\barre}{\ensuremath{\|}}
% environnement proof (amsthm non chargé) utilisé par les modèles
\ifdefined\proof\else\newenvironment{proof}[1][Démonstration]{\par\noindent\textit{#1.}\ }{\qed\par}\fi
\usepackage{lastpage}
\usepackage{hyperref}
\hypersetup{hidelinks}

% ============================================================
% Palette « Pop » OnBuch+ — mêmes hex que la classe P (plus_ui.dart)
% ============================================================
\definecolor{poporange}{HTML}{F59321}
\definecolor{popdark}{HTML}{E07A0C}
\definecolor{popcream}{HTML}{FFF6E8}
\definecolor{popink}{HTML}{1C1714}
\definecolor{poppurple}{HTML}{7A5AE0}
\definecolor{popgreen}{HTML}{2E9E6A}
\definecolor{poppink}{HTML}{E0457B}
\definecolor{popblue}{HTML}{2F7DE1}
\definecolor{popgold}{HTML}{C98A12}
\definecolor{popmuted}{HTML}{8A7F76}
\definecolor{popline}{HTML}{EADFD2}
\definecolor{popgray}{HTML}{8A7F76}
\colorlet{popgrey}{popgray}
% Alias tolérants (le modèle nomme parfois les couleurs différemment)
\colorlet{popwhite}{white}
\colorlet{popblack}{popink}
\colorlet{popred}{poppink}
\colorlet{popyellow}{popgold}
% Teintes claires pour fonds de tableaux / zones
\colorlet{popblueL}{popblue!10!white}
\colorlet{poporangeL}{poporange!14!white}
\colorlet{popgreenL}{popgreen!12!white}
\colorlet{poppurpleL}{poppurple!12!white}
\colorlet{poppinkL}{poppink!12!white}
\colorlet{popgoldL}{popgold!14!white}

\pagecolor{popcream}

% ============================================================
% Typographie
% ============================================================
\defaultfontfeatures{Ligatures=TeX}
\setmainfont{PlusJakartaSans-Medium.ttf}[Path=../../../fonts/,BoldFont=PlusJakartaSans-Bold.ttf]
% Maths en sans-serif (Fira Math) avec chiffres et lettres droites de Plus
% Jakarta, pour que les formules se fondent dans le texte.
\usepackage[math-style=ISO,bold-style=ISO]{unicode-math}
\setmathfont{FiraMath-Regular.otf}[Path=../../../fonts/]
\setmathfont{PlusJakartaSans-Medium.ttf}[Path=../../../fonts/,range={up/num,up/{Latin,latin},\mathup}]
\usepackage{icomma} % virgule décimale sans espace en mode math
% Caractères mathématiques Unicode absents des polices de texte (Plus Jakarta,
% Archivo Black) : rendus par leur équivalent mathématique, en texte comme en
% formule — sinon ils s'affichent en carré vide (ex. « Arithmétique dans ℤ »).
\usepackage{newunicodechar}
\newunicodechar{ℝ}{\ensuremath{\mathbb{R}}}
\newunicodechar{ℤ}{\ensuremath{\mathbb{Z}}}
\newunicodechar{ℂ}{\ensuremath{\mathbb{C}}}
\newunicodechar{ℕ}{\ensuremath{\mathbb{N}}}
\newunicodechar{ℚ}{\ensuremath{\mathbb{Q}}}
\newunicodechar{𝔻}{\ensuremath{\mathbb{D}}}
\newunicodechar{α}{\ensuremath{\alpha}}
\newunicodechar{β}{\ensuremath{\beta}}
\newunicodechar{γ}{\ensuremath{\gamma}}
\newunicodechar{δ}{\ensuremath{\delta}}
\newunicodechar{ε}{\ensuremath{\varepsilon}}
\newunicodechar{θ}{\ensuremath{\theta}}
\newunicodechar{λ}{\ensuremath{\lambda}}
\newunicodechar{μ}{\ensuremath{\mu}}
\newunicodechar{σ}{\ensuremath{\sigma}}
\newunicodechar{τ}{\ensuremath{\tau}}
\newunicodechar{φ}{\ensuremath{\varphi}}
\newunicodechar{ψ}{\ensuremath{\psi}}
\newunicodechar{ω}{\ensuremath{\omega}}
\newunicodechar{Σ}{\ensuremath{\Sigma}}
% majuscules produites par \MakeUppercase dans les titres (α-aminés -> Α)
\newunicodechar{Α}{\ensuremath{\alpha}}
\newunicodechar{Β}{\ensuremath{\beta}}
\newunicodechar{Γ}{\ensuremath{\Gamma}}
\newunicodechar{Δ}{\ensuremath{\Delta}}
\newunicodechar{Ε}{\ensuremath{\varepsilon}}
\newunicodechar{Θ}{\ensuremath{\Theta}}
\newunicodechar{Λ}{\ensuremath{\Lambda}}
\newunicodechar{Μ}{\ensuremath{\mu}}
\newunicodechar{Τ}{\ensuremath{\tau}}
\newunicodechar{Φ}{\ensuremath{\Phi}}
\newunicodechar{Ψ}{\ensuremath{\Psi}}
\newunicodechar{Ω}{\ensuremath{\Omega}}
\newunicodechar{↦}{\ensuremath{\mapsto}}
\newunicodechar{∈}{\ensuremath{\in}}
\newunicodechar{∉}{\ensuremath{\notin}}
\newunicodechar{∧}{\ensuremath{\wedge}}
\newunicodechar{⇔}{\ensuremath{\Leftrightarrow}}
\newunicodechar{⟺}{\ensuremath{\Longleftrightarrow}}
\newunicodechar{⇒}{\ensuremath{\Rightarrow}}
\newunicodechar{⟹}{\ensuremath{\Longrightarrow}}
\newunicodechar{∘}{\ensuremath{\circ}}
\newunicodechar{∪}{\ensuremath{\cup}}
\newunicodechar{∩}{\ensuremath{\cap}}
\newunicodechar{≡}{\ensuremath{\equiv}}
\newunicodechar{⊂}{\ensuremath{\subset}}
\newunicodechar{⊥}{\ensuremath{\perp}}
\newunicodechar{∃}{\ensuremath{\exists}}
\newunicodechar{∀}{\ensuremath{\forall}}
\newunicodechar{∛}{\ensuremath{\sqrt[3]{\,}}}
\newunicodechar{∜}{\ensuremath{\sqrt[4]{\,}}}
\newunicodechar{⃗}{\ensuremath{{}^{\to}}}
\newunicodechar{ⁿ}{\ifmmode^{n}\else\textsuperscript{\ensuremath{n}}\fi}
\newunicodechar{ˣ}{\ifmmode^{x}\else\textsuperscript{\ensuremath{x}}\fi}
\newunicodechar{ᵇ}{\ifmmode^{b}\else\textsuperscript{\ensuremath{b}}\fi}
\newunicodechar{ʳ}{\ifmmode^{r}\else\textsuperscript{\ensuremath{r}}\fi}
\newunicodechar{ᵘ}{\ifmmode^{u}\else\textsuperscript{\ensuremath{u}}\fi}
\newunicodechar{ʸ}{\ifmmode^{y}\else\textsuperscript{\ensuremath{y}}\fi}
\newunicodechar{ᵃ}{\ifmmode^{a}\else\textsuperscript{\ensuremath{a}}\fi}
\newunicodechar{ᵉ}{\ifmmode^{e}\else\textsuperscript{\ensuremath{e}}\fi}
\newunicodechar{ᵏ}{\ifmmode^{k}\else\textsuperscript{\ensuremath{k}}\fi}
\newunicodechar{ᵖ}{\ifmmode^{p}\else\textsuperscript{\ensuremath{p}}\fi}
\newunicodechar{ᴺ}{\ifmmode^{N}\else\textsuperscript{\ensuremath{N}}\fi}
\newunicodechar{ᶜ}{\ifmmode^{c}\else\textsuperscript{\ensuremath{c}}\fi}
\newunicodechar{ʰ}{\ifmmode^{h}\else\textsuperscript{\ensuremath{h}}\fi}
\newunicodechar{ᵠ}{\ifmmode^{q}\else\textsuperscript{\ensuremath{q}}\fi}
\newunicodechar{ₙ}{\ifmmode_{n}\else\textsubscript{\ensuremath{n}}\fi}
\newunicodechar{ₐ}{\ifmmode_{a}\else\textsubscript{\ensuremath{a}}\fi}
\newunicodechar{ᵢ}{\ifmmode_{i}\else\textsubscript{\ensuremath{i}}\fi}
\newunicodechar{ᵣ}{\ifmmode_{r}\else\textsubscript{\ensuremath{r}}\fi}
\newunicodechar{ₖ}{\ifmmode_{k}\else\textsubscript{\ensuremath{k}}\fi}
\newunicodechar{ᵦ}{\ifmmode_{\beta}\else\textsubscript{\ensuremath{\beta}}\fi}
\newunicodechar{ₓ}{\ifmmode_{x}\else\textsubscript{\ensuremath{x}}\fi}
\newfontfamily\popdisplay{ArchivoBlack-Regular.ttf}[Path=../../../fonts/]
\newfontfamily\poplogo{SpaceGrotesk-Bold.ttf}[Path=../../../fonts/]
\newfontfamily\popsemi{PlusJakartaSans-SemiBold.ttf}[Path=../../../fonts/]
\newfontfamily\popxbold{PlusJakartaSans-ExtraBold.ttf}[Path=../../../fonts/]
\newfontfamily\popxboldls{PlusJakartaSans-ExtraBold.ttf}[Path=../../../fonts/,LetterSpace=3.0]

\setlist[enumerate]{leftmargin=1.7em,itemsep=5pt,topsep=4pt}
\setlist[itemize]{leftmargin=1.7em,itemsep=4pt,topsep=4pt,label={\color{poporange}\textbullet}}
\setlength{\parindent}{0pt}
\setlength{\parskip}{5pt}
\renewcommand{\arraystretch}{1.25}

% pgfplots : style Pop par défaut
\pgfplotsset{
  every axis/.append style={axis line style={popink,line width=1pt},tick style={popink},
    grid style={popline,line width=0.5pt},label style={font=\small},tick label style={font=\footnotesize}},
  popcurve/.style={poporange,line width=1.8pt,no markers,smooth},
}
% Même style disponible dans un \draw TikZ simple (hors pgfplots)
\tikzset{popcurve/.style={poporange,line width=1.8pt}}

% ============================================================
% Pill / squiggle
% ============================================================
\newcommand{\poppill}[3]{%
  \tikz[baseline=-0.55ex]\node[draw=popink,line width=1.2pt,rounded corners=2.6pt,%
    fill=#2,inner xsep=6.5pt,inner ysep=3pt]{%
    \popxboldls\fontsize{7}{7}\selectfont\color{#3}\MakeUppercase{#1}};%
}
\newcommand{\popsquiggle}[1]{%
  \tikz[baseline=-0.4ex]{
    \draw[#1,line width=2.6pt,line cap=round]
      plot [smooth] coordinates {(0,0.09) (0.34,0.01) (0.68,0.21) (1.02,0.01) (1.36,0.10)};
  }%
}
% Mot-clé surligné (marqueur orange sous le texte)
\newcommand{\cle}[1]{{\popxbold\color{popdark}#1}}

% ============================================================
% En-tête / pied
% ============================================================
\pagestyle{fancy}
\fancyhf{}
\renewcommand{\headrulewidth}{0pt}
\renewcommand{\footrulewidth}{0pt}
\lhead{\raisebox{-0.15ex}{\poplogo\fontsize{13}{13}\selectfont\color{popink}OnBuch\color{poporange}+}}
\rhead{\poppill{\DOCMATIERE\ \textbullet\ \DOCNIVEAU\ \textbullet\ Leçon \DOCLECON}{white}{popink}}
\lfoot{\popxbold\fontsize{7.6}{7.6}\selectfont\color{popmuted}\MakeUppercase{Cours signé OnBuch+}\ \textbullet\ {\popsemi\DOCTITRE}}
\rfoot{\tikz[baseline=-0.6ex]\node[rounded corners=8.5pt,fill=popink,minimum height=17pt,minimum width=17pt,inner xsep=5pt,inner sep=0]{\popxbold\fontsize{8}{8}\selectfont\color{popcream}\thepage\,/\,\pageref*{LastPage}};}

% ============================================================
% Titres
% ============================================================
\newcommand{\chaptitre}[2]{%
  \begin{tcolorbox}[enhanced,colback=poporange,colframe=popink,boxrule=2.6pt,arc=11pt,
      drop shadow=popink,left=15pt,right=15pt,top=12pt,bottom=11pt,before skip=3pt,after skip=18pt]
    \poppill{#2}{popink}{popcream}\par\vspace{9pt}
    {\raggedright\hyphenpenalty=10000\exhyphenpenalty=10000\popdisplay\fontsize{22}{24}\selectfont\color{popcream}\MakeUppercase{#1}\par}
  \end{tcolorbox}
}
% Section numérotée : pastille orange avec le numéro + titre Archivo
\newcounter{popsec}
\newcommand{\coursec}[1]{%
  \stepcounter{popsec}\setcounter{popsub}{0}%
  \par\vspace{16pt}\noindent
  \begin{minipage}{\linewidth}
  \tikz[baseline=-0.7ex]\node[draw=popink,line width=1.6pt,fill=poporange,rounded corners=4pt,minimum size=22pt,inner sep=0]{\popdisplay\fontsize{12}{12}\selectfont\color{popcream}\thepopsec};%
  \hspace{8pt}{\popdisplay\fontsize{15}{17}\selectfont\color{popink}\MakeUppercase{#1}}\par
  \vspace{4pt}\noindent{\color{poporange}\rule{36pt}{3.6pt}}
  \end{minipage}\par\nopagebreak\vspace{8pt}
}
\newcounter{popsub}
\newcommand{\courssub}[1]{%
  \stepcounter{popsub}%
  \par\vspace{9pt}\noindent{\popxbold\fontsize{11.5}{13}\selectfont\color{popink}{\color{poporange}\thepopsec.\thepopsub}\hspace{6pt}#1}\par\nopagebreak\vspace{4pt}
}

% ============================================================
% Boîtes pédagogiques — même recette : fond blanc, bordure épaisse colorée,
% ombre dure encre, badge plein. Titre optionnel [#1].
% ============================================================
\tcbset{popbase/.style={breakable,enhanced,colback=white,boxrule=2.3pt,arc=9pt,
  shadow={3pt}{-3pt}{0pt}{popink},left=14pt,right=14pt,top=11pt,bottom=11pt,before skip=11pt,after skip=15pt,
  fontupper=\popsemi\fontsize{10.6}{14.5}\selectfont\color{popink},
  fontlower=\popsemi\fontsize{10.6}{14.5}\selectfont\color{popink}}}
% Coche dessinée (le glyphe ✓ n'existe pas dans les polices du document)
\newcommand{\popcheck}[1]{\tikz[baseline=-0.5ex]\draw[#1,line width=1.6pt,line cap=round,line join=round](-0.13,0.0)--(-0.04,-0.09)--(0.13,0.1);}
\newcommand{\popboxhead}[3]{\poppill{#1}{#2}{white}\ifx\relax#3\relax\else\hspace{6pt}{\popxbold\fontsize{10.6}{12}\selectfont\color{popink}#3}\fi\par\vspace{7pt}}

\newtcolorbox{aretenir}[1][]{popbase,colframe=popgreen,before upper={\popboxhead{À retenir}{popgreen}{#1}}}
\newtcolorbox{exemplebox}[1][]{popbase,colframe=poporange,before upper={\popboxhead{Exemple}{poporange}{#1}}}
\newtcolorbox{definition}[1][]{popbase,colframe=popblue,before upper={\popboxhead{Définition}{popblue}{#1}}}
\newtcolorbox{propriete}[1][]{popbase,colframe=poppurple,before upper={\popboxhead{Propriété}{poppurple}{#1}}}
\newtcolorbox{methode}[1][]{popbase,colframe=popgold,before upper={\popboxhead{Méthode}{popgold}{#1}}}
\newtcolorbox{attention}[1][]{popbase,colframe=poppink,before upper={\popboxhead{Attention}{poppink}{#1}}}
\newtcolorbox{experience}[1][]{popbase,colframe=popblue,colback=popblueL,before upper={\popboxhead{Expérience}{popblue}{#1}}}
% « Le savais-tu ? » : carte violette pleine (contexte, culture, Cameroun)
\newtcolorbox{savaistu}[1][]{popbase,colback=poppurple,colframe=popink,
  fontupper=\popsemi\fontsize{10.4}{14.2}\selectfont\color{white},
  before upper={\poppill{Le savais-tu\,?}{white}{poppurple}\ifx\relax#1\relax\else\hspace{6pt}{\popxbold\color{white}#1}\fi\par\vspace{7pt}}}
% Exercice résolu : énoncé en haut, \tcblower puis la solution sur fond crème
\newcounter{popexo}\newcounter{popexr}
\newtcolorbox{exoresolu}[1][]{popbase,colframe=poporange,colbacklower=popcream,
  segmentation style={popink,line width=1.2pt,dashed},
  before upper={\stepcounter{popexr}\popboxhead{Exercice résolu \thepopexr}{poporange}{#1}},
  before lower={\poppill{Solution}{popink}{popcream}\par\vspace{6pt}}}
% Exercice (non résolu en place) — la correction est dans la partie Corrigés
\newtcolorbox{exercice}[1][]{popbase,colframe=popink,
  before upper={\stepcounter{popexo}\popboxhead{Exercice \thepopexo}{popink}{#1}}}
% Corrigé
\newtcolorbox{corrige}[1][]{popbase,colframe=popgreen,colback=popgreenL,
  before upper={\popboxhead{Corrigé}{popgreen}{#1}}}
% Objectifs / prérequis / bilan
\newtcolorbox{objectifs}{popbase,colframe=popink,colback=poporangeL,
  before upper={\popboxhead{Objectifs de la leçon}{popink}{}}}
\newtcolorbox{prerequis}{popbase,colframe=popink,colback=popblueL,
  before upper={\popboxhead{Prérequis}{popblue}{}}}
\newtcolorbox{bilan}{popbase,colframe=popink,colback=white,boxrule=2.8pt,
  before upper={\popboxhead{Fiche bilan}{poporange}{L'essentiel à connaître pour l'examen}}}
% Figure encadrée avec légende
\newtcolorbox{popfigure}[1][]{enhanced,breakable=false,colback=white,colframe=popline,boxrule=1.4pt,arc=8pt,
  left=10pt,right=10pt,top=10pt,bottom=8pt,before skip=10pt,after skip=12pt,halign=center,
  lower separated=false,halign lower=center,
  fontlower=\popsemi\fontsize{9}{11.5}\selectfont\color{popmuted},
  #1}
\newcommand{\legende}[1]{\par\vspace{6pt}{\popsemi\fontsize{9}{11.5}\selectfont\color{popmuted}#1}}

% ============================================================
% Signature OnBuch+
% ============================================================
\newcommand{\poplogoinline}[1]{{\poplogo\fontsize{#1}{#1}\selectfont\color{popink}OnBuch\color{poporange}+}}
\newcommand{\signatureonbuch}{%
  \begin{tcolorbox}[enhanced,colback=popink,colframe=popink,boxrule=2.2pt,arc=10pt,
      drop shadow=poporange,left=14pt,right=14pt,top=10pt,bottom=10pt,before skip=0pt,after skip=0pt]
    \tikz[baseline=-0.7ex]\node[circle,fill=popgreen,draw=popcream,line width=1.3pt,minimum size=22pt,inner sep=0]{\popcheck{white}};%
    \hspace{9pt}\begin{minipage}[c]{0.62\linewidth}
      {\popxbold\fontsize{12}{13}\selectfont\color{popcream}Signé\ \textbullet\ {\poplogo OnBuch\color{poporange}+}}\par\vspace{2pt}
      {\raggedright\popsemi\fontsize{8}{10.5}\selectfont\color{popline}\MakeUppercase{Équipe pédagogique OnBuch+}\\\MakeUppercase{Conforme au programme officiel MINESEC}\par}
    \end{minipage}\hfill
    {\poplogo\fontsize{22}{22}\selectfont\color{popcream}OnBuch\color{poporange}+}
  \end{tcolorbox}%
}

% ============================================================
% Page de garde
% #1 titre · #2 matière · #3 niveau · #4 module · #5 n° leçon · #6 durée ·
% #7 description / objectifs (texte ou itemize)
% ============================================================
\newcommand{\pagegarde}[7]{%
  \thispagestyle{empty}
  \newgeometry{a4paper,margin=1.7cm,top=1.6cm,bottom=1.4cm}%
  \begin{tikzpicture}[remember picture,overlay]
    \fill[poporange!18] ([xshift=-3.2cm,yshift=-3.6cm]current page.north east) circle (4.6cm);
    \fill[poppurple!14] ([xshift=2.4cm,yshift=7.6cm]current page.south west) circle (3.8cm);
  \end{tikzpicture}%
  {\poplogo\fontsize{30}{30}\selectfont\color{popink}OnBuch\color{poporange}+}\hfill
  \raisebox{0.6ex}{\poppill{Cours signé \textbullet\ Premium}{popink}{popcream}}\par
  \vspace{1pt}\popsquiggle{poppurple}\par
  \vspace{8pt}
  \begin{tcolorbox}[enhanced,colback=poporange,colframe=popink,boxrule=2.8pt,arc=13pt,
      drop shadow=popink,left=18pt,right=18pt,top=12pt,bottom=13pt,after skip=10pt]
    \poppill{#2 \textbullet\ #3}{popink}{popcream}\hspace{5pt}\poppill{#4}{white}{popink}\par\vspace{8pt}
    {\popdisplay\fontsize{14}{14}\selectfont\color{popink}LEÇON #5}\par\vspace{4pt}
    {\raggedright\hyphenpenalty=10000\exhyphenpenalty=10000\popdisplay\fontsize{25}{27}\selectfont\color{popcream}\MakeUppercase{#1}\par}
    \vspace{8pt}
    {\popxbold\fontsize{9.5}{9.5}\selectfont\color{popink}\MakeUppercase{Durée conseillée : #6}}
  \end{tcolorbox}
  \begin{tcolorbox}[enhanced,colback=white,colframe=popink,boxrule=2.2pt,arc=10pt,
      drop shadow=popink,left=16pt,right=16pt,top=10pt,bottom=10pt,after skip=8pt]
    \poppill{Dans cette leçon}{poppurple}{white}\par\vspace{5pt}
    \popsemi\fontsize{9.3}{12.6}\selectfont\color{popink}#7
  \end{tcolorbox}
  \noindent\begin{minipage}[t]{0.487\linewidth}
    \begin{tcolorbox}[enhanced,colback=popgreen,colframe=popink,boxrule=2.2pt,arc=10pt,
        drop shadow=popink,left=11pt,right=11pt,top=10pt,bottom=10pt,height=3.1cm,valign=center]
      \tcbox[colback=white,colframe=popink,boxrule=1.3pt,arc=5pt,boxsep=0pt,nobeforeafter,
        left=3pt,right=3pt,top=3pt,bottom=3pt,valign=center]{\qrcode[height=1.9cm]{https://play.google.com/store/apps/details?id=com.luvvix.onbuch&hl=fr}}%
      \hspace{8pt}\begin{minipage}[c]{0.5\linewidth}
        {\popdisplay\fontsize{11}{12.5}\selectfont\color{white}\MakeUppercase{Télécharge OnBuch}}\par\vspace{4pt}
        {\raggedright\popsemi\fontsize{8.4}{11}\selectfont\color{white}Gratuit \textbullet\ Google Play\\Tuteur IA, annales, concours, résultats\par}
      \end{minipage}
    \end{tcolorbox}
  \end{minipage}\hfill
  \begin{minipage}[t]{0.487\linewidth}
    \begin{tcolorbox}[enhanced,colback=poppurple,colframe=popink,boxrule=2.2pt,arc=10pt,
        drop shadow=popink,left=11pt,right=11pt,top=10pt,bottom=10pt,height=3.1cm,valign=center]
      \tikz[baseline=-0.7ex]\node[circle,fill=white,draw=popink,line width=1.3pt,minimum size=1.6cm,inner sep=0]{\popdisplay\fontsize{24}{24}\selectfont\color{poppurple}+};%
      \hspace{8pt}\begin{minipage}[c]{0.6\linewidth}
        {\popdisplay\fontsize{11}{12.5}\selectfont\color{white}\MakeUppercase{Passe à OnBuch+}}\par\vspace{4pt}
        {\raggedright\popsemi\fontsize{8.4}{11}\selectfont\color{white}Tous les cours signés, Tuteur IA illimité, fiches PDF — onglet OnBuch+ de l'app\par}
      \end{minipage}
    \end{tcolorbox}
  \end{minipage}
  \vspace{8pt}
  \popcontenu
  \vfill
  \signatureonbuch
  \restoregeometry
  \newpage
}

% Rangée « Ce cours contient » (page de garde)
\newcommand{\poptile}[3]{% #1 couleur #2 gros texte #3 légende
  \begin{tcolorbox}[enhanced,colback=#1,colframe=popink,boxrule=2pt,arc=9pt,shadow={2.5pt}{-2.5pt}{0pt}{popink},
      left=6pt,right=6pt,top=9pt,bottom=9pt,halign=center,valign=center,height=1.75cm,width=\linewidth,nobeforeafter]
    {\popdisplay\fontsize{12.5}{13}\selectfont\color{white}\MakeUppercase{#2}}\par\vspace{3pt}
    {\centering\popsemi\fontsize{7.8}{9.5}\selectfont\color{white}#3\par}
  \end{tcolorbox}}
\newcommand{\popcontenu}{%
  \par\noindent{\popxboldls\fontsize{7.5}{7.5}\selectfont\color{popmuted}CE COURS SIGNÉ CONTIENT}\par\vspace{7pt}
  \noindent\begin{minipage}[t]{0.235\linewidth}\poptile{popblue}{Cours}{complet, illustré, pas à pas}\end{minipage}\hfill
  \begin{minipage}[t]{0.235\linewidth}\poptile{poporange}{Méthodes}{et exercices résolus}\end{minipage}\hfill
  \begin{minipage}[t]{0.235\linewidth}\poptile{poppink}{Exercices}{type examen + corrigés}\end{minipage}\hfill
  \begin{minipage}[t]{0.235\linewidth}\poptile{popgreen}{Fiche bilan}{l'essentiel à retenir}\end{minipage}\par}

% Sommaire
\newtcolorbox{sommaire}{enhanced,colback=white,colframe=popink,boxrule=2.2pt,arc=10pt,
  drop shadow=popink,left=18pt,right=18pt,top=13pt,bottom=13pt,before skip=6pt,after skip=20pt,
  before upper={\poppill{Sommaire}{poppurple}{white}\par\vspace{9pt}\popsemi\fontsize{10.6}{14.5}\selectfont\color{popink}}}

% Signature de fin de cours — poussée tout en bas de la dernière page
\newcommand{\findecours}{%
  \par\vfill\nopagebreak
  \begin{center}\popsquiggle{poporange}\end{center}\vspace{4pt}
  \signatureonbuch
}
\AtBeginDocument{\def\llbracket{\mathopen{[\mkern-3.5mu[}}\def\rrbracket{\mathclose{]\mkern-3.5mu]}}} % intervalles d'entiers (glyphe ⟦ absent de la police)
% Glyphes absents de Fira Math : redéfinis à partir de glyphes disponibles
\AtBeginDocument{%
  \def\setminus{\mathbin{\backslash}}\let\smallsetminus\setminus
  \def\vdots{\mathord{\vbox{\baselineskip3.2pt\lineskiplimit0pt\kern4pt\hbox{.}\hbox{.}\hbox{.}}}}%
  \def\ddots{\mathinner{\mkern1mu\raise6.5pt\hbox{.}\mkern2mu\raise3.6pt\hbox{.}\mkern2mu\raise0.7pt\hbox{.}\mkern1mu}}%
  \def\triangle{\mathord{\scalebox{0.9}{$\symup{\Delta}$}}}%
  \def\nabla{\mathord{\raisebox{\depth}{\scalebox{1}[-1]{$\symup{\Delta}$}}}}%
  \def\checkmark{\popcheck{popdark}}%
  \def\star{\mathbin{\ast}}\def\bigstar{\mathord{\ast}}%
  \def\diamond{\mathbin{\scalebox{0.6}{$\Diamond$}}}%
  \def\bigcup{\mathop{\mathchoice{\raisebox{-0.3em}{\scalebox{1.7}{$\cup$}}}{\scalebox{1.25}{$\cup$}}{\cup}{\cup}}\displaylimits}%
  \def\bigcap{\mathop{\mathchoice{\raisebox{-0.3em}{\scalebox{1.7}{$\cap$}}}{\scalebox{1.25}{$\cap$}}{\cap}{\cap}}\displaylimits}%
  \def\lesseqqgtr{\mathrel{\lessgtr}}\def\bigoplus{\mathop{\oplus}\displaylimits}%
}
\providecommand{\ssi}{si et seulement si} % abréviation fréquente
\AtBeginDocument{\providecommand{\wideparen}{\overparen}} % arc de cercle

%% FIGFIX-BEGIN v4 — correctifs globaux des figures (docs/onbuch-generation/tools/figfix)
\usepackage{adjustbox}\usepackage{etoolbox}
% toute figure TikZ/pgfplots plus large que la ligne est réduite à la largeur disponible
\BeforeBeginEnvironment{tikzpicture}{\begin{adjustbox}{max width=\linewidth}}
\AfterEndEnvironment{tikzpicture}{\end{adjustbox}}
% légendes pgfplots lisibles par-dessus les courbes
\pgfplotsset{every axis/.append style={legend style={fill=white,fill opacity=0.92,text opacity=1,draw=black!15,inner sep=3pt}}}
% étiquettes d'arêtes (syntaxe edge["..."]) avec fond blanc
\tikzset{every edge quotes/.append style={fill=white,inner sep=1.2pt}}
%% FIGFIX-END
\begin{document}

\pagegarde{Equations de droites}{Mathématiques}{3ème}{Mathématiques – classe de 3ème}{N10}{4 séances (fiche de progression harmonisée)}{Cette leçon apprend à caractériser une droite du plan par une équation cartésienne ou réduite, à partir de deux points, d'un vecteur directeur ou d'un coefficient directeur. Elle permet ensuite de déterminer algébriquement les positions relatives de deux droites (sécantes, parallèles, confondues) et de modéliser des situations concrètes de la vie courante.
\vspace{4pt}
\begin{multicols}{2}\small
\begin{itemize}
\item À la fin de cette leçon, je sais reconnaître et interpréter une équation cartésienne de droite ax + by + c = 0
\item À la fin de cette leçon, je sais déterminer une équation cartésienne d'une droite passant par deux points donnés
\item À la fin de cette leçon, je sais déterminer une équation cartésienne d'une droite dont on connaît un vecteur directeur et un point
\item À la fin de cette leçon, je sais déterminer l'équation réduite y = mx + p d'une droite de coefficient directeur donné
\item À la fin de cette leçon, je sais lire graphiquement le coefficient directeur et l'ordonnée à l'origine d'une droite
\item À la fin de cette leçon, je sais déterminer algébriquement les positions relatives de deux droites (sécantes, strictement parallèles, confondues)
\end{itemize}\end{multicols}}

\begin{sommaire}
\begin{multicols}{2}
\begin{enumerate}
\item Équation cartésienne d'une droite
\item Droite passant par deux points
\item Droite de vecteur directeur donné
\item Équation réduite et coefficient directeur
\item Positions relatives de deux droites
\item Applications et modélisation
\item Activité d'intégration
\item Méthodes et exercices résolus
\item Exercices
\item Corrigés des exercices
\item Fiche bilan
\end{enumerate}
\end{multicols}
\end{sommaire}

\begin{objectifs}
\begin{itemize}
\item À la fin de cette leçon, je sais reconnaître et interpréter une équation cartésienne de droite ax + by + c = 0
\item À la fin de cette leçon, je sais déterminer une équation cartésienne d'une droite passant par deux points donnés
\item À la fin de cette leçon, je sais déterminer une équation cartésienne d'une droite dont on connaît un vecteur directeur et un point
\item À la fin de cette leçon, je sais déterminer l'équation réduite y = mx + p d'une droite de coefficient directeur donné
\item À la fin de cette leçon, je sais lire graphiquement le coefficient directeur et l'ordonnée à l'origine d'une droite
\item À la fin de cette leçon, je sais déterminer algébriquement les positions relatives de deux droites (sécantes, strictement parallèles, confondues)
\item À la fin de cette leçon, je sais calculer les coordonnées du point d'intersection de deux droites sécantes
\item À la fin de cette leçon, je sais modéliser par des équations de droites des situations concrètes de la vie courante (tarifs, trajets, consommation) et rédiger une démarche justifiée
\end{itemize}
\end{objectifs}

\begin{prerequis}
\begin{itemize}
\item Repérage dans le plan : coordonnées d'un point, lecture graphique (classe de 4ème)
\item Vecteurs : coordonnées d'un vecteur, colinéarité de deux vecteurs, déterminant (début d'année de 3ème)
\item Fonctions linéaires et affines : représentation graphique de y = ax et y = ax + b (4ème)
\item Résolution d'équations du premier degré et de systèmes de deux équations à deux inconnues (4ème)
\item Milieu d'un segment et distance dans un repère orthonormé (début d'année de 3ème)
\end{itemize}
\end{prerequis}

\newpage

\coursec{Équation cartésienne d'une droite}

À Yaoundé, une entreprise de livraison de denrées alimentaires facture ses courses selon la distance parcourue : le prix affiché suit une relation linéaire. Un chauffeur de taxi-moto à Mokolo veut prévoir le coût d'une course avant de partir, et vérifier si deux itinéraires proposés par son application se croisent ou restent parallèles. Comment traduire ces trajectoires et ces tarifs par des équations de droites dans un repère ? C'est tout l'objet de cette leçon : donner à chaque droite une \cle{carte d'identité algébrique}.

Dans cette première section, nous allons découvrir ce qu'est une \cle{équation cartésienne} d'une droite, comment la reconnaître, comment vérifier qu'un point appartient à une droite, et comment lire sur cette équation la direction de la droite. Les sections suivantes apprendront ensuite à \emph{construire} ces équations à partir de données géométriques.

\courssub{Situation de découverte : une droite, une infinité de points}

Plaçons-nous dans un repère orthonormé $(O\,;\vec{i},\vec{j})$. Considérons la droite $(d)$ passant par le point $A(0\,;2)$ et dirigée par le vecteur $\vec{u}(3\,;2)$. Traçons-la : depuis $A$, on se déplace de $3$ unités vers la droite et de $2$ unités vers le haut, on obtient le point $B(3\,;4)$, et la droite $(AB)$ est tracée.

Prenons maintenant un point quelconque $M(x\,;y)$ du plan. À quelle condition $M$ se trouve-t-il \emph{sur} la droite $(d)$ ? Réponse géométrique : exactement lorsque le vecteur $\overrightarrow{AM}$ est \cle{colinéaire} au vecteur $\vec{u}$. Or tu sais calculer les coordonnées de $\overrightarrow{AM}$ :
\[
\overrightarrow{AM}\begin{pmatrix} x - 0 \\ y - 2 \end{pmatrix} = \begin{pmatrix} x \\ y - 2 \end{pmatrix}.
\]
La condition de colinéarité entre $\overrightarrow{AM}(x\,;y-2)$ et $\vec{u}(3\,;2)$ s'écrit, avec le critère de colinéarité vu dans la leçon sur les vecteurs ($xy' - x'y = 0$) :
\[
x \times 2 - (y-2) \times 3 = 0 \quad\Longleftrightarrow\quad 2x - 3y + 6 = 0.
\]
Voilà le miracle : l'appartenance d'un point à une droite, qui est une question de \emph{géométrie}, se traduit par une \emph{égalité algébrique} portant sur les coordonnées. Cette égalité $2x - 3y + 6 = 0$ est ce qu'on appelle une \textbf{équation cartésienne} de la droite $(d)$. Toute la section consiste à comprendre et manipuler cette idée.

\begin{savaistu}[Pourquoi « cartésienne » ?]
Le nom rend hommage au philosophe et mathématicien français \textbf{René Descartes} (1596--1650). Dans son ouvrage \emph{La Géométrie} (1637), il eut l'idée révolutionnaire de repérer les points du plan par des nombres (les coordonnées) et de traduire les figures par des équations. Ce pont entre l'algèbre et la géométrie, appelé \emph{géométrie analytique}, est l'un des outils les plus puissants des mathématiques : c'est lui qui permet aujourd'hui aux applications de navigation de ton téléphone de calculer des itinéraires à Yaoundé ou à Douala !
\end{savaistu}

\courssub{Définition d'une équation cartésienne}

\begin{definition}[Équation cartésienne d'une droite]
Le plan est muni d'un repère $(O\,;\vec{i},\vec{j})$.
\begin{itemize}
\item Toute droite $(d)$ du plan admet une équation de la forme
\[
ax + by + c = 0, \qquad \text{avec } (a\,;b) \neq (0\,;0),
\]
appelée \cle{équation cartésienne} de $(d)$.
\item Réciproquement, toute équation de la forme $ax + by + c = 0$ avec $(a\,;b) \neq (0\,;0)$ est l'équation d'une droite du plan.
\end{itemize}
Un point $M(x\,;y)$ appartient à la droite $(d)$ \textbf{si et seulement si} ses coordonnées vérifient l'équation : $ax + by + c = 0$.
\end{definition}

Décortiquons cette définition. L'équation $ax + by + c = 0$ fait intervenir trois nombres fixés $a$, $b$, $c$ (les \cle{coefficients}) et deux inconnues $x$ et $y$ (les coordonnées du point). La condition $(a\,;b) \neq (0\,;0)$ signifie que $a$ et $b$ ne sont \textbf{pas tous les deux nuls} : en effet, si $a = 0$ \emph{et} $b = 0$, l'équation devient $c = 0$, qui ne dépend plus du tout des coordonnées et ne peut donc pas décrire une droite.

\begin{attention}[La condition $(a\,;b) \neq (0\,;0)$ n'est pas facultative]
Beaucoup d'élèves oublient cette hypothèse. Elle est pourtant essentielle : si $a = b = 0$, l'équation $0x + 0y + c = 0$ se réduit à $c = 0$. Si $c \neq 0$, aucun point ne la vérifie (ensemble vide) ; si $c = 0$, tous les points la vérifient (le plan entier). Dans les deux cas, on n'obtient pas une droite ! Retiens : \textbf{au moins l'un des deux coefficients $a$ ou $b$ est non nul}.
\end{attention}

\courssub{Lien avec le vecteur directeur : une propriété fondamentale}

Rappelons d'abord un fait géométrique essentiel, déjà rencontré dans la leçon sur les vecteurs : \textbf{une droite est entièrement déterminée par un point et une direction}. La direction est donnée par un \cle{vecteur directeur}, c'est-à-dire un vecteur non nul parallèle à la droite. Bien sûr, une même droite possède une infinité de vecteurs directeurs : tous les vecteurs non nuls colinéaires entre eux qui la dirigent.

La question naturelle est alors : si l'on connaît l'équation $ax + by + c = 0$ d'une droite, peut-on lire directement un vecteur directeur sur les coefficients ? La réponse est oui, et c'est un résultat à connaître parfaitement.

\begin{propriete}[Vecteur directeur lu sur l'équation]
Si la droite $(d)$ a pour équation cartésienne $ax + by + c = 0$, alors le vecteur
\[
\vec{u}\begin{pmatrix} -b \\ a \end{pmatrix}
\]
est un \cle{vecteur directeur} de $(d)$. Tout vecteur non nul colinéaire à $\vec{u}$ est aussi un vecteur directeur de $(d)$.
\end{propriete}

Cette propriété se justifie très naturellement par la colinéarité, exactement comme dans notre situation de découverte. Voici le raisonnement, qui est formateur.

\begin{experience}[Pourquoi $\vec{u}(-b\,;a)$ dirige-t-il la droite ?]
Soit $(d)$ d'équation $ax + by + c = 0$, et soit $A(x_A\,;y_A)$ un point de $(d)$, donc $ax_A + by_A + c = 0$. Soit $M(x\,;y)$ un point quelconque du plan.
\begin{enumerate}
\item $M \in (d)$ signifie $ax + by + c = 0$.
\item Retranchons membre à membre les deux égalités :
\[
(ax + by + c) - (ax_A + by_A + c) = 0 \quad\Longleftrightarrow\quad a(x - x_A) + b(y - y_A) = 0.
\]
\item Or le vecteur $\overrightarrow{AM}$ a pour coordonnées $(x - x_A\,;\,y - y_A)$. Considérons le vecteur $\vec{u}(-b\,;a)$. Le critère de colinéarité appliqué à $\overrightarrow{AM}$ et $\vec{u}$ donne :
\[
(x - x_A)\times a - (y - y_A)\times(-b) = a(x-x_A) + b(y-y_A) = 0.
\]
\item Ce critère est vérifié, donc $\overrightarrow{AM}$ et $\vec{u}$ sont \textbf{colinéaires} : le point $M$ est sur la droite passant par $A$ et dirigée par $\vec{u}$. Cette droite est donc exactement $(d)$, et $\vec{u}(-b\,;a)$ en est un vecteur directeur. $\blacksquare$
\end{enumerate}
\end{experience}

\begin{exemplebox}[Lecture d'un vecteur directeur]
Soit $(d)$ la droite d'équation $2x - 3y + 6 = 0$. On a $a = 2$, $b = -3$, $c = 6$. Un vecteur directeur est donc
\[
\vec{u}\begin{pmatrix} -b \\ a \end{pmatrix} = \begin{pmatrix} -(-3) \\ 2 \end{pmatrix} = \begin{pmatrix} 3 \\ 2 \end{pmatrix}.
\]
C'est bien le vecteur $\vec{u}(3\,;2)$ de notre situation de découverte : tout est cohérent ! Le vecteur $\vec{u}'(-6\,;-4) = -2\vec{u}$ est aussi un vecteur directeur de $(d)$.
\end{exemplebox}

\begin{methode}[Retenir la formule $\vec{u}(-b\,;a)$ sans se tromper]
\begin{enumerate}
\item On échange les coefficients de $x$ et de $y$ : $(a\,;b) \to (b\,;a)$.
\item On change le signe du premier : $(-b\,;a)$.
\item Vérification mentale : la combinaison $a\times(-b) + b\times a$ vaut toujours $0$ : c'est le signe que l'équation $ax+by+c=0$ « annule » la direction $(-b\,;a)$.
\end{enumerate}
Remarque : écrire $\vec{u}(b\,;-a)$ est aussi correct, car $(b\,;-a) = -(-b\,;a)$ est colinéaire à $(-b\,;a)$. En revanche, $\vec{u}(a\,;b)$ est \textbf{faux} en général.
\end{methode}

\begin{popfigure}
\begin{tikzpicture}
\begin{axis}[
  width=0.85\linewidth,
  axis lines=middle,
  xlabel={$x$}, ylabel={$y$},
  xmin=-4.5, xmax=4.5, ymin=-3.5, ymax=5.5,
  xtick={-4,-3,-2,-1,1,2,3,4},
  ytick={-3,-2,-1,1,2,3,4,5},
  grid=both, grid style={popline!40},
  axis line style={-{Stealth}},
  label style={font=\small},
  ticklabel style={font=\scriptsize}
]
\addplot[popcurve,domain=-4.5:4.5,samples=2]{(2*x+6)/3};
\addplot[only marks,mark=*,mark size=2pt,poporange] coordinates {(0,2)};
\node[poporange,above left,font=\small] at (axis cs:0,2) {$A(0\,;2)$};
\addplot[-{Stealth},very thick,popgreen] coordinates {(0,2) (3,4)};
\node[popgreen,above right,font=\small] at (axis cs:2.4,3.6) {$\vec{u}(3\,;2)$};
\node[popblue,rotate=33,font=\small] at (axis cs:-3.2,-0.9) {$(d) : 2x - 3y + 6 = 0$};
\end{axis}
\end{tikzpicture}
\legende{La droite $(d) : 2x - 3y + 6 = 0$, son point $A(0\,;2)$ et un vecteur directeur $\vec{u}(3\,;2)$.}
\end{popfigure}

\courssub{Appartenance d'un point à une droite}

La définition donne immédiatement un critère pratique, qu'il faut savoir appliquer les yeux fermés.

\begin{methode}[Tester si un point appartient à une droite]
Pour savoir si $M(x_0\,;y_0)$ appartient à la droite $(d) : ax + by + c = 0$ :
\begin{enumerate}
\item On remplace $x$ par $x_0$ et $y$ par $y_0$ dans le membre de gauche $ax + by + c$.
\item On calcule la valeur obtenue.
\item Si le résultat vaut $0$, alors $M \in (d)$ ; sinon, $M \notin (d)$.
\end{enumerate}
\end{methode}

\begin{exoresolu}[Les points $A(1\,;4)$ et $B(-3\,;0)$ appartiennent-ils à $(d) : 2x - 3y + 6 = 0$ ?]
\textbf{Énoncé.} On considère la droite $(d)$ d'équation $2x - 3y + 6 = 0$. Les points $A(1\,;4)$ et $B(-3\,;0)$ appartiennent-ils à $(d)$ ?
\tcblower
\textbf{Solution détaillée.}
Pour $A(1\,;4)$, on remplace $x$ par $1$ et $y$ par $4$ :
\[
2\times 1 - 3\times 4 + 6 = 2 - 12 + 6 = -4.
\]
Comme $-4 \neq 0$, le point $A$ \textbf{n'appartient pas} à la droite $(d)$.
Pour $B(-3\,;0)$, on remplace $x$ par $-3$ et $y$ par $0$ :
\[
2\times(-3) - 3\times 0 + 6 = -6 + 0 + 6 = 0.
\]
Le résultat est nul, donc $B \in (d)$. D'ailleurs, $B$ est le point d'intersection de $(d)$ avec l'axe des abscisses, comme on le verra plus bas.
\end{exoresolu}

\courssub{Non-unicité de l'équation cartésienne}

Voici un point délicat mais capital. Reprenons notre droite $(d) : 2x - 3y + 6 = 0$. Multiplions tous les coefficients par $2$ : on obtient l'équation $4x - 6y + 12 = 0$. Un point vérifie la première équation si et seulement s'il vérifie la seconde (car $4x - 6y + 12 = 2(2x - 3y + 6)$, et un produit est nul quand l'un de ses facteurs est nul). Les deux équations décrivent donc \textbf{exactement la même droite}.

\begin{aretenir}[Une droite possède une infinité d'équations cartésiennes]
Si $ax + by + c = 0$ est une équation cartésienne de $(d)$, alors pour tout nombre réel non nul $k$, l'équation
\[
k(ax + by + c) = 0 \quad\text{c'est-à-dire}\quad kax + kby + kc = 0
\]
est aussi une équation cartésienne de la \textbf{même} droite $(d)$. On dit que les deux équations sont \cle{proportionnelles}.
\end{aretenir}

\begin{attention}[Erreur fréquente : croire l'équation unique]
Au BEPC, beaucoup d'élèves paniquent lorsque leur équation « ne ressemble pas » à celle du corrigé. Retiens bien : $2x - 3y + 6 = 0$, $4x - 6y + 12 = 0$, $-2x + 3y - 6 = 0$ et même $x - \num{1,5}y + 3 = 0$ désignent \textbf{toutes la même droite}. Pour comparer deux équations, vérifie si leurs coefficients sont proportionnels. En revanche, $2x - 3y + 5 = 0$ n'est \emph{pas} la même droite que $(d)$ : le coefficient constant n'est pas dans la même proportion.
\end{attention}

\courssub{Cas particuliers : droites parallèles aux axes}

Que se passe-t-il lorsque l'un des coefficients $a$ ou $b$ est nul ? (Rappel : pas les deux à la fois !)

\textbf{Premier cas : $a = 0$.} L'équation devient $by + c = 0$, soit $y = -\dfrac{c}{b}$ (division licite car $b \neq 0$). En posant $k = -\dfrac{c}{b}$, on obtient $y = k$ : c'est l'ensemble des points dont l'ordonnée vaut $k$, c'est-à-dire une droite \textbf{parallèle à l'axe des abscisses} (horizontale). Un vecteur directeur est $\vec{u}(-b\,;0)$, bien colinéaire à $\vec{i}$.

\textbf{Deuxième cas : $b = 0$.} L'équation devient $ax + c = 0$, soit $x = -\dfrac{c}{a}$ (car $a \neq 0$). En posant $k = -\dfrac{c}{a}$, on obtient $x = k$ : c'est l'ensemble des points dont l'abscisse vaut $k$, c'est-à-dire une droite \textbf{parallèle à l'axe des ordonnées} (verticale). Un vecteur directeur est $\vec{u}(0\,;a)$, colinéaire à $\vec{j}$.

\begin{aretenir}[Droites parallèles aux axes]
\begin{itemize}
\item Une équation de la forme $x = k$ (avec $k$ réel) représente une droite \textbf{parallèle à l'axe des ordonnées}, passant par le point $(k\,;0)$.
\item Une équation de la forme $y = k$ représente une droite \textbf{parallèle à l'axe des abscisses}, passant par le point $(0\,;k)$.
\item En particulier, l'axe des abscisses a pour équation $y = 0$ et l'axe des ordonnées a pour équation $x = 0$.
\end{itemize}
\end{aretenir}

\begin{popfigure}
\begin{tikzpicture}
\begin{axis}[
  width=0.8\linewidth,
  axis lines=middle,
  xlabel={$x$}, ylabel={$y$},
  xmin=-3.5, xmax=4.5, ymin=-3.5, ymax=4.5,
  xtick={-3,-2,-1,1,2,3,4},
  ytick={-3,-2,-1,1,2,3,4},
  grid=both, grid style={popline!40},
  axis line style={-{Stealth}},
  ticklabel style={font=\scriptsize}
]
\addplot[very thick,popblue] coordinates {(2,-3.5) (2,4.5)};
\node[popblue,right,font=\small] at (axis cs:2,4) {$(d_1) : x = 2$};
\addplot[very thick,poporange,domain=-3.5:4.5,samples=2]{-1};
\node[poporange,above right,font=\small] at (axis cs:3.4,-1) {$(d_2) : y = -1$};
\addplot[only marks,mark=*,mark size=2pt,popdark] coordinates {(2,-1)};
\node[popdark,below left,font=\small] at (axis cs:2,-1) {$(2\,;-1)$};
\end{axis}
\end{tikzpicture}
\legende{Les droites $x = 2$ (verticale) et $y = -1$ (horizontale) se coupent au point $(2\,;-1)$.}
\end{popfigure}

\courssub{Intersections d'une droite avec les axes du repère}

Une application directe et très utile du test d'appartenance consiste à trouver où une droite coupe les axes du repère. C'est souvent la première étape pour la tracer proprement.

\begin{methode}[Trouver les points d'intersection avec les axes]
Soit $(d) : ax + by + c = 0$.
\begin{enumerate}
\item \textbf{Intersection avec l'axe des abscisses} (équation $y = 0$) : on pose $y = 0$ dans l'équation de $(d)$ et on résout $ax + c = 0$. Si $a \neq 0$, on trouve le point $\left(-\dfrac{c}{a}\,;0\right)$.
\item \textbf{Intersection avec l'axe des ordonnées} (équation $x = 0$) : on pose $x = 0$ et on résout $by + c = 0$. Si $b \neq 0$, on trouve le point $\left(0\,;-\dfrac{c}{b}\right)$.
\end{enumerate}
Ces deux points suffisent à tracer la droite à la règle !
\end{methode}

\begin{exoresolu}[Tracer $(d) : 2x - 3y + 6 = 0$ à partir de ses intersections avec les axes]
\textbf{Énoncé.} Déterminer les coordonnées des points d'intersection de la droite $(d) : 2x - 3y + 6 = 0$ avec les axes du repère, puis expliquer comment tracer $(d)$.
\tcblower
\textbf{Solution détaillée.}
\textbf{Avec l'axe des abscisses} : on pose $y = 0$ :
\[
2x - 3\times 0 + 6 = 0 \quad\Longleftrightarrow\quad 2x = -6 \quad\Longleftrightarrow\quad x = -3.
\]
La droite coupe l'axe des abscisses en $B(-3\,;0)$.
\textbf{Avec l'axe des ordonnées} : on pose $x = 0$ :
\[
2\times 0 - 3y + 6 = 0 \quad\Longleftrightarrow\quad -3y = -6 \quad\Longleftrightarrow\quad y = 2.
\]
La droite coupe l'axe des ordonnées en $A(0\,;2)$.
\textbf{Tracé} : on place les points $A(0\,;2)$ et $B(-3\,;0)$ dans le repère, puis on trace à la règle la droite passant par ces deux points. On retrouve exactement la figure étudiée plus haut.
\end{exoresolu}

\begin{attention}[Deux pièges à éviter]
\begin{itemize}
\item Ne confonds pas les deux axes : pour l'axe des \emph{abscisses}, c'est $y$ qui vaut $0$ ; pour l'axe des \emph{ordonnées}, c'est $x$ qui vaut $0$. Beaucoup d'élèves inversent !
\item Si $c = 0$, la droite passe par l'origine $O(0\,;0)$ : les deux intersections avec les axes se confondent en un seul point. Il faut alors trouver un \emph{autre} point de la droite (en choisissant par exemple $x = 1$) pour pouvoir la tracer.
\end{itemize}
\end{attention}

\courssub{Synthèse et ouverture}

Récapitulons ce que nous avons établi dans cette section. À toute droite du plan est associée une famille d'équations de la forme $ax + by + c = 0$ avec $(a\,;b) \neq (0\,;0)$, toutes proportionnelles entre elles. Cette « carte d'identité algébrique » permet de :
\begin{itemize}
\item \textbf{tester l'appartenance} d'un point à la droite par un simple calcul ;
\item \textbf{lire un vecteur directeur} : $\vec{u}(-b\,;a)$ ;
\item \textbf{tracer la droite} grâce à ses intersections avec les axes ;
\item \textbf{reconnaître les cas particuliers} : $x = k$ (verticale) et $y = k$ (horizontale).
\end{itemize}

Revenons à notre chauffeur de taxi-moto de Mokolo : si son application lui indique que le tarif suit la relation $p = 150d + 100$ (prix en FCFA, distance en km), on peut réécrire cette relation sous la forme $150d - p + 100 = 0$ : c'est bien une équation cartésienne de droite ! Prévoir un coût revient à tester l'appartenance d'un point à cette droite. Dans les sections suivantes, nous apprendrons à \emph{fabriquer} ces équations à partir de deux points, d'un vecteur directeur ou d'un coefficient directeur, puis à étudier les positions relatives de deux droites — c'est-à-dire à déterminer si deux itinéraires se croisent ou restent parallèles.

\begin{exercice}[À toi de jouer]
On considère la droite $(d) : 3x + 2y - 12 = 0$.
\begin{enumerate}
\item Les points $M(2\,;3)$ et $N(4\,;1)$ appartiennent-ils à $(d)$ ?
\item Donner un vecteur directeur de $(d)$, puis un deuxième, colinéaire au premier.
\item Donner une autre équation cartésienne de $(d)$.
\item Déterminer les points d'intersection de $(d)$ avec les axes du repère, puis tracer $(d)$.
\end{enumerate}
\end{exercice}

\begin{corrige}[Exercice : droite $3x + 2y - 12 = 0$]
\begin{enumerate}
\item Pour $M(2\,;3)$ : $3\times 2 + 2\times 3 - 12 = 6 + 6 - 12 = 0$, donc $M \in (d)$. Pour $N(4\,;1)$ : $3\times 4 + 2\times 1 - 12 = 12 + 2 - 12 = 2 \neq 0$, donc $N \notin (d)$.
\item Avec $a = 3$ et $b = 2$, un vecteur directeur est $\vec{u}(-2\,;3)$. Un autre : $\vec{u}'(2\,;-3) = -\vec{u}$.
\item En multipliant par $2$ : $6x + 4y - 24 = 0$ est une autre équation de $(d)$.
\item Axe des abscisses ($y = 0$) : $3x - 12 = 0$ donne $x = 4$, point $(4\,;0)$. Axe des ordonnées ($x = 0$) : $2y - 12 = 0$ donne $y = 6$, point $(0\,;6)$. On trace la droite passant par $(4\,;0)$ et $(0\,;6)$.
\end{enumerate}
\end{corrige}

\coursec{Droite passant par deux points}

Dans la section précédente, nous avons découvert ce qu'est une \cle{équation cartésienne} d'une droite : une relation de la forme $ax + by + c = 0$ que vérifient les coordonnées $(x\,;y)$ de tous les points de la droite, et d'eux seulement. Nous avons vu qu'une droite possède une infinité d'équations cartésiennes, toutes proportionnelles entre elles. Mais une question essentielle reste en suspens : \textbf{comment fabriquer concrètement une telle équation ?}

La situation la plus naturelle est la suivante. On te donne deux points $A$ et $B$ du plan, repérés par leurs coordonnées. Depuis la classe de 6\textsuperscript{e}, tu sais que \emph{par deux points distincts passe une unique droite}. La droite $(AB)$ existe donc, elle est unique : il ne reste qu'à en trouver une équation. C'est exactement ce que fait un logiciel comme GeoGebra quand tu cliques sur deux points : il calcule instantanément une équation de la droite qui les relie. Dans cette section, tu vas apprendre à faire ce calcul toi-même, à la main, avec un outil que tu connais déjà : la \cle{colinéarité} de deux vecteurs et le \cle{déterminant}.

\courssub{Le problème et l'idée directrice}

\begin{definition}[Le problème posé]
Le plan est muni d'un repère $(O\,;\vec{\imath},\vec{\jmath})$. On donne deux points \textbf{distincts} $A(x_A\,;y_A)$ et $B(x_B\,;y_B)$. Déterminer une équation cartésienne de la droite $(AB)$, c'est trouver trois nombres $a$, $b$, $c$ (avec $a$ et $b$ non tous les deux nuls) tels que :
\[ M(x\,;y) \in (AB) \iff ax + by + c = 0. \]
\end{definition}

L'idée fondamentale repose sur une observation géométrique très simple. Prends un point $M$ quelconque du plan. Dire que $M$ appartient à la droite $(AB)$, c'est dire que les points $A$, $B$ et $M$ sont \textbf{alignés}. Or tu as appris en début d'année que trois points sont alignés si et seulement si les vecteurs $\overrightarrow{AM}$ et $\overrightarrow{AB}$ sont \cle{colinéaires}. Et la colinéarité se teste par un calcul : le \textbf{déterminant} des deux vecteurs doit être nul.

\begin{aretenir}[La traduction en trois langages]
\begin{center}
\begin{tabularx}{\linewidth}{|l|X|}
\hline
\rowcolor{popblueL} \textbf{Langage} & \textbf{Énoncé} \\
\hline
Géométrie & $M(x\,;y)$ appartient à la droite $(AB)$ \\
\hline
Vectoriel & les vecteurs $\overrightarrow{AM}$ et $\overrightarrow{AB}$ sont colinéaires \\
\hline
Algébrique & $\det\left(\overrightarrow{AM},\overrightarrow{AB}\right) = 0$ \\
\hline
\end{tabularx}
\end{center}
C'est ce passage du langage géométrique au langage algébrique qui va produire l'équation cherchée.
\end{aretenir}

\courssub{Étape 1 : le vecteur $\overrightarrow{AB}$, vecteur directeur de $(AB)$}

La première chose à calculer, c'est le vecteur qui \emph{dirige} la droite : le vecteur $\overrightarrow{AB}$.

\begin{propriete}[Coordonnées du vecteur $\overrightarrow{AB}$]
Si $A(x_A\,;y_A)$ et $B(x_B\,;y_B)$, alors le vecteur $\overrightarrow{AB}$ a pour coordonnées :
\[ \overrightarrow{AB}\begin{pmatrix} x_B - x_A \\ y_B - y_A \end{pmatrix}. \]
On retient la formule : \textbf{« extrémité moins origine »}. Comme $A$ et $B$ sont distincts, ce vecteur est non nul : c'est bien un \cle{vecteur directeur} de la droite $(AB)$.
\end{propriete}

\begin{attention}[L'erreur classique : inverser origine et extrémité]
L'erreur la plus fréquente consiste à calculer $(x_A - x_B\,;y_A - y_B)$ au lieu de $(x_B - x_A\,;y_B - y_A)$, ou pire, à mélanger : $(x_B - x_A\,;y_A - y_B)$. Retiens la règle : dans $\overrightarrow{AB}$, la lettre $A$ est l'\textbf{origine}, la lettre $B$ est l'\textbf{extrémité}. On calcule toujours \textbf{extrémité moins origine}, coordonnée par coordonnée : $x$ avec $x$, $y$ avec $y$. Une bonne habitude : écris les coordonnées de $B$ sur une ligne, celles de $A$ juste en dessous, puis soustrais verticalement.
\end{attention}

\courssub{Étape 2 : la condition de colinéarité}

Soit maintenant $M(x\,;y)$ un point quelconque du plan. Le vecteur $\overrightarrow{AM}$ a pour coordonnées :
\[ \overrightarrow{AM}\begin{pmatrix} x - x_A \\ y - y_A \end{pmatrix}. \]

\begin{propriete}[Critère de colinéarité par le déterminant]
Deux vecteurs $\vec{u}\begin{pmatrix} x \\ y \end{pmatrix}$ et $\vec{v}\begin{pmatrix} x' \\ y' \end{pmatrix}$ sont colinéaires si et seulement si leur déterminant est nul :
\[ \det(\vec{u},\vec{v}) = xy' - x'y = 0. \]
\end{propriete}

Appliqué à notre situation, ce critère donne la condition d'appartenance cherchée :
\[ M(x\,;y) \in (AB) \iff \det\left(\overrightarrow{AM},\overrightarrow{AB}\right) = 0 \iff (x - x_A)(y_B - y_A) - (x_B - x_A)(y - y_A) = 0. \]

\begin{attention}[Bien placer les facteurs dans le déterminant]
Pour deux vecteurs $\vec{u}(x\,;y)$ et $\vec{v}(x'\,;y')$, le déterminant vaut $xy' - x'y$ : on multiplie \emph{en croix}. Avec $\overrightarrow{AM}(x - x_A\,;y - y_A)$ et $\overrightarrow{AB}(x_B - x_A\,;y_B - y_A)$, cela donne $(x - x_A)\times(y_B - y_A) \;-\; (x_B - x_A)\times(y - y_A)$. Ne permute pas les coordonnées $x$ et $y$ entre les deux vecteurs !
\end{attention}

\courssub{Étape 3 : développer et simplifier}

Il ne reste plus qu'à développer l'égalité obtenue et à regrouper les termes pour la mettre sous la forme $ax + by + c = 0$. C'est un simple travail de calcul littéral, que tu maîtrises depuis le chapitre sur le calcul littéral : développement, réduction, puis éventuellement division de tous les coefficients par un même nombre non nul pour simplifier l'équation.

\begin{methode}[Trouver une équation de la droite $(AB)$]
\begin{enumerate}
\item \textbf{Calculer} les coordonnées du vecteur $\overrightarrow{AB}(x_B - x_A\,;y_B - y_A)$, vecteur directeur de $(AB)$.
\item \textbf{Écrire} qu'un point $M(x\,;y)$ appartient à $(AB)$ si et seulement si $\det\left(\overrightarrow{AM},\overrightarrow{AB}\right) = 0$, c'est-à-dire :
\[ (x - x_A)(y_B - y_A) - (x_B - x_A)(y - y_A) = 0. \]
\item \textbf{Développer}, réduire et ordonner pour obtenir la forme $ax + by + c = 0$. Simplifier si possible en divisant tous les coefficients par un même nombre non nul.
\item \textbf{Vérifier} : remplacer $(x\,;y)$ par les coordonnées de $A$, puis par celles de $B$ ; l'égalité doit être vraie dans les deux cas.
\end{enumerate}
\end{methode}

\courssub{Le théorème et sa démonstration guidée}

\begin{propriete}[Théorème — droite définie par deux points]
Par deux points distincts $A(x_A\,;y_A)$ et $B(x_B\,;y_B)$ du plan passe une unique droite $(AB)$. Une équation cartésienne de cette droite s'obtient en écrivant que, pour tout point $M(x\,;y)$ :
\[ M \in (AB) \iff \det\left(\overrightarrow{AM},\overrightarrow{AB}\right) = 0, \]
ce qui, après développement, donne une équation de la forme $ax + by + c = 0$ avec $(a\,;b) \neq (0\,;0)$.
\end{propriete}

\begin{experience}[Démonstration guidée du théorème]
Nous allons prouver que la condition $\det\left(\overrightarrow{AM},\overrightarrow{AB}\right) = 0$ fournit bien une équation cartésienne de droite, c'est-à-dire une équation de la forme $ax + by + c = 0$ où $a$ et $b$ ne sont pas tous les deux nuls.
\begin{enumerate}
\item \textbf{Existence et unicité de la droite.} $A$ et $B$ étant distincts, le vecteur $\overrightarrow{AB}$ est non nul. La droite $(AB)$ est l'ensemble des points $M$ tels que $\overrightarrow{AM}$ soit colinéaire à $\overrightarrow{AB}$ : c'est la définition même de la droite passant par $A$ et dirigée par $\overrightarrow{AB}$. Elle existe et elle est unique.
\item \textbf{Traduction algébrique.} Posons $\overrightarrow{AB}\begin{pmatrix} x_B - x_A \\ y_B - y_A \end{pmatrix}$ et $\overrightarrow{AM}\begin{pmatrix} x - x_A \\ y - y_A \end{pmatrix}$. La colinéarité équivaut à :
\[ (x - x_A)(y_B - y_A) - (x_B - x_A)(y - y_A) = 0. \]
\item \textbf{Développement.} En développant chaque produit :
\[ x(y_B - y_A) - x_A(y_B - y_A) - y(x_B - x_A) + y_A(x_B - x_A) = 0. \]
Regroupons les termes en $x$, en $y$ et les constantes :
\[ \underbrace{(y_B - y_A)}_{a}\,x \;+\; \underbrace{(x_A - x_B)}_{b}\,y \;+\; \underbrace{y_A(x_B - x_A) - x_A(y_B - y_A)}_{c} = 0. \]
C'est bien une équation de la forme $ax + by + c = 0$.
\item \textbf{Les coefficients $a$ et $b$ ne sont pas tous les deux nuls.} Si l'on avait à la fois $a = y_B - y_A = 0$ et $b = x_A - x_B = 0$, alors $x_A = x_B$ et $y_A = y_B$, autrement dit $A = B$, ce qui est exclu puisque les points sont distincts. Le couple $(a\,;b)$ est donc non nul : l'équation obtenue est bien une équation de droite.
\item \textbf{Conclusion.} L'ensemble des points $M$ vérifiant $\det\left(\overrightarrow{AM},\overrightarrow{AB}\right) = 0$ est exactement la droite $(AB)$, et la condition obtenue est bien une équation cartésienne de droite. $\blacksquare$
\end{enumerate}
\end{experience}

\courssub{Exemple chiffré complet}

\begin{exemplebox}[Une équation de la droite $(AB)$ avec $A(1\,;2)$ et $B(4\,;5)$]
\textbf{Étape 1.} Coordonnées du vecteur directeur (extrémité moins origine) :
\[ \overrightarrow{AB}\begin{pmatrix} 4 - 1 \\ 5 - 2 \end{pmatrix} = \begin{pmatrix} 3 \\ 3 \end{pmatrix}. \]
\textbf{Étape 2.} Soit $M(x\,;y)$. Alors $\overrightarrow{AM}\begin{pmatrix} x - 1 \\ y - 2 \end{pmatrix}$. La condition de colinéarité s'écrit :
\[ \det\left(\overrightarrow{AM},\overrightarrow{AB}\right) = 3(x - 1) - 3(y - 2) = 0. \]
\textbf{Étape 3.} Développons :
\[ 3x - 3 - 3y + 6 = 0 \iff 3x - 3y + 3 = 0. \]
En divisant tous les coefficients par $3$ :
\[ (AB) : x - y + 1 = 0. \]
\textbf{Étape 4. Vérification.} Pour $A(1\,;2)$ : $1 - 2 + 1 = 0$ \checkmark. Pour $B(4\,;5)$ : $4 - 5 + 1 = 0$ \checkmark. Les deux points vérifient l'équation : le résultat est cohérent.
\end{exemplebox}

\begin{popfigure}
\begin{tikzpicture}[scale=0.9]
\draw[popline,->] (-1,0) -- (6,0) node[below left] {$x$};
\draw[popline,->] (0,-1) -- (0,6.5) node[below left] {$y$};
\draw[popline] (0,0) node[below left] {$O$};
\foreach \x in {1,2,3,4,5} \draw[popline] (\x,0.08) -- (\x,-0.08) node[below] {\footnotesize $\x$};
\foreach \y in {1,2,3,4,5,6} \draw[popline] (0.08,\y) -- (-0.08,\y) node[left] {\footnotesize $\y$};
\draw[popblue,thick,domain=-1.2:5.5] plot (\x,{\x+1});
\node[popblue] at (5.1,6.35) {$(AB)$};
\draw[poporange,very thick,->] (1,2) -- (4,5) node[midway,above left] {$\overrightarrow{AB}$};
\draw[poppurple,very thick,->] (1,2) -- (3,4) node[midway,below right] {$\overrightarrow{AM}$};
\fill[popink] (1,2) circle (2.2pt) node[below left] {$A(1\,;2)$};
\fill[popink] (4,5) circle (2.2pt) node[above right] {$B(4\,;5)$};
\fill[poppurple] (3,4) circle (2.2pt) node[above left] {$M(x\,;y)$};
\end{tikzpicture}
\legende{La droite $(AB)$ d'équation $x - y + 1 = 0$. Un point $M$ est sur $(AB)$ si et seulement si $\overrightarrow{AM}$ et $\overrightarrow{AB}$ sont colinéaires.}
\end{popfigure}

\courssub{Cas particuliers : droites verticales et horizontales}

Que se passe-t-il lorsque les deux points ont la même abscisse ou la même ordonnée ? La méthode par le déterminant fonctionne toujours, et elle aboutit à des équations remarquablement simples.

\begin{propriete}[Droite verticale et droite horizontale]
Soit $A(x_A\,;y_A)$ et $B(x_B\,;y_B)$ deux points distincts.
\begin{itemize}
\item Si $x_A = x_B$ (même abscisse), la droite $(AB)$ est \textbf{verticale} (parallèle à l'axe des ordonnées) et a pour équation $x = x_A$.
\item Si $y_A = y_B$ (même ordonnée), la droite $(AB)$ est \textbf{horizontale} (parallèle à l'axe des abscisses) et a pour équation $y = y_A$.
\end{itemize}
\end{propriete}

\begin{exemplebox}[Vérification sur les cas particuliers]
\textbf{Cas vertical.} $A(3\,;1)$ et $B(3\,;4)$ : $\overrightarrow{AB}(0\,;3)$. La condition $\det\left(\overrightarrow{AM},\overrightarrow{AB}\right) = 0$ donne $3(x - 3) - 0(y - 1) = 0$, soit $3x - 9 = 0$, c'est-à-dire $x = 3$. C'est bien une droite verticale.

\textbf{Cas horizontal.} $A(1\,;2)$ et $B(5\,;2)$ : $\overrightarrow{AB}(4\,;0)$. La condition donne $0(x - 1) - 4(y - 2) = 0$, soit $-4y + 8 = 0$, c'est-à-dire $y = 2$. C'est bien une droite horizontale.
\end{exemplebox}

\begin{attention}[Une droite verticale n'a pas de coefficient directeur]
L'équation $x = 3$ est bien une équation cartésienne (elle s'écrit $1x + 0y - 3 = 0$), mais elle ne peut pas se mettre sous la forme $y = mx + p$ : on ne peut pas « isoler $y$ ». N'écris donc jamais qu'une droite verticale a un coefficient directeur nul : elle n'en a \emph{pas du tout}. Nous y reviendrons dans la section sur l'équation réduite.
\end{attention}

\courssub{La vérification systématique : ton filet de sécurité}

Une équation de droite se fabrique par une chaîne de calculs : coordonnées du vecteur, déterminant, développement, simplification. Une erreur de signe peut se glisser à chaque maillon. Heureusement, il existe un test de fiabilité simple et rapide : \textbf{les coordonnées de $A$ et de $B$ doivent vérifier l'équation trouvée}.

\begin{methode}[Vérifier une équation de $(AB)$]
\begin{enumerate}
\item Remplace $x$ et $y$ par les coordonnées de $A$ dans l'équation : le membre de gauche doit valoir $0$.
\item Fais de même avec les coordonnées de $B$.
\item Si l'un des deux tests échoue, l'équation est fausse : reprends le calcul du vecteur $\overrightarrow{AB}$, puis le déterminant, puis le développement.
\end{enumerate}
Ce test prend moins d'une minute et détecte la quasi-totalité des erreurs. Au BEPC, c'est un réflexe qui peut te sauver des points.
\end{methode}

\begin{exoresolu}[Déterminer et vérifier une équation de droite]
\textbf{Énoncé.} Dans un repère, on donne $E(-2\,;3)$ et $F(4\,;-1)$.
\begin{enumerate}
\item Déterminer une équation cartésienne de la droite $(EF)$.
\item Le point $G(10\,;-5)$ appartient-il à la droite $(EF)$ ?
\end{enumerate}
\tcblower
\textbf{Solution détaillée.}
\begin{enumerate}
\item \textbf{Étape 1 :} $\overrightarrow{EF}\begin{pmatrix} 4 - (-2) \\ -1 - 3 \end{pmatrix} = \begin{pmatrix} 6 \\ -4 \end{pmatrix}$.

\textbf{Étape 2 :} soit $M(x\,;y)$, alors $\overrightarrow{EM}\begin{pmatrix} x - (-2) \\ y - 3 \end{pmatrix} = \begin{pmatrix} x + 2 \\ y - 3 \end{pmatrix}$. La colinéarité donne :
\[ -4(x + 2) - 6(y - 3) = 0. \]

\textbf{Étape 3 :} développons :
\begin{align*}
-4x - 8 - 6y + 18 &= 0\\
-4x - 6y + 10 &= 0
\end{align*}
En divisant par $-2$ : $(EF) : 2x + 3y - 5 = 0$.

\textbf{Étape 4 : vérification.} Pour $E(-2\,;3)$ : $2\times(-2) + 3\times 3 - 5 = -4 + 9 - 5 = 0$ \checkmark. Pour $F(4\,;-1)$ : $2\times 4 + 3\times(-1) - 5 = 8 - 3 - 5 = 0$ \checkmark.
\item Pour $G(10\,;-5)$ : $2\times 10 + 3\times(-5) - 5 = 20 - 15 - 5 = 0$. Les coordonnées de $G$ vérifient l'équation, donc $G \in (EF)$ : les points $E$, $F$ et $G$ sont alignés.
\end{enumerate}
\end{exoresolu}

\begin{exercice}[À toi de jouer]
Dans un repère, on donne les points $K(2\,;-1)$ et $L(5\,;3)$.
\begin{enumerate}
\item Déterminer une équation cartésienne de la droite $(KL)$.
\item Vérifier que les coordonnées de $K$ et de $L$ satisfont l'équation trouvée.
\item Le point $N(8\,;7)$ appartient-il à $(KL)$ ? Et le point $P(0\,;-3)$ ?
\end{enumerate}
\end{exercice}

\begin{corrige}[Exercice : droite $(KL)$]
\begin{enumerate}
\item $\overrightarrow{KL}\begin{pmatrix} 5 - 2 \\ 3 - (-1) \end{pmatrix} = \begin{pmatrix} 3 \\ 4 \end{pmatrix}$. Pour $M(x\,;y)$, $\overrightarrow{KM}\begin{pmatrix} x - 2 \\ y - (-1) \end{pmatrix} = \begin{pmatrix} x - 2 \\ y + 1 \end{pmatrix}$. Colinéarité : $4(x - 2) - 3(y + 1) = 0$, soit $4x - 8 - 3y - 3 = 0$, d'où $(KL) : 4x - 3y - 11 = 0$.
\item Pour $K(2\,;-1)$ : $4\times 2 - 3\times(-1) - 11 = 8 + 3 - 11 = 0$ \checkmark. Pour $L(5\,;3)$ : $4\times 5 - 3\times 3 - 11 = 20 - 9 - 11 = 0$ \checkmark.
\item Pour $N(8\,;7)$ : $4\times 8 - 3\times 7 - 11 = 32 - 21 - 11 = 0$, donc $N \in (KL)$. Pour $P(0\,;-3)$ : $4\times 0 - 3\times(-3) - 11 = 0 + 9 - 11 = -2 \neq 0$, donc $P \notin (KL)$.
\end{enumerate}
\end{corrige}

\begin{savaistu}[Comment GeoGebra trace tes droites]
Quand tu places deux points $A$ et $B$ dans GeoGebra et que tu traces la droite $(AB)$, le logiciel effectue exactement le calcul de cette section : il calcule les coordonnées de $\overrightarrow{AB}$, écrit la nullité du déterminant $\det\left(\overrightarrow{AM},\overrightarrow{AB}\right)$, puis simplifie l'équation obtenue. C'est pourquoi la fenêtre « Algèbre » affiche instantanément quelque chose comme $x - y = -1$, qui est notre équation $x - y + 1 = 0$ écrite autrement. La même idée est utilisée en cartographie numérique : pour tracer la route qui relie deux villes sur une carte GPS (par exemple Douala et Yaoundé, repérées par leurs coordonnées), le logiciel calcule l'équation de la droite passant par les deux points. Les mathématiques de la classe de 3\textsuperscript{e} sont littéralement embarquées dans ton téléphone !
\end{savaistu}

\begin{aretenir}[L'essentiel de la section]
\begin{itemize}
\item Par deux points distincts $A$ et $B$ passe une unique droite $(AB)$, dirigée par $\overrightarrow{AB}(x_B - x_A\,;y_B - y_A)$.
\item Un point $M(x\,;y)$ appartient à $(AB)$ si et seulement si $\det\left(\overrightarrow{AM},\overrightarrow{AB}\right) = 0$, c'est-à-dire $(x - x_A)(y_B - y_A) - (x_B - x_A)(y - y_A) = 0$.
\item On développe et on simplifie pour obtenir la forme $ax + by + c = 0$.
\item Cas particuliers : $x_A = x_B$ donne la droite verticale $x = x_A$ ; $y_A = y_B$ donne la droite horizontale $y = y_A$.
\item On vérifie toujours le résultat en testant les coordonnées de $A$ et de $B$ dans l'équation.
\end{itemize}
\end{aretenir}

Tu sais désormais construire l'équation d'une droite à partir de deux points. Dans la section suivante, nous verrons comment construire directement l'équation d'une droite lorsqu'on connaît un point et un \cle{vecteur directeur} donné, ce qui rendra le calcul encore plus rapide.

\coursec{Droite de vecteur directeur donné}

Dans la section précédente, nous avons appris à déterminer l'équation d'une droite connue par \textbf{deux de ses points}. Mais ce n'est pas la seule façon de décrire une droite. Imagine un chauffeur de taxi-moto à Yaoundé : pour décrire son trajet en ligne droite, il suffit de donner \emph{où il part} et \emph{dans quelle direction il roule}. En mathématiques, c'est exactement la même idée : une droite est entièrement déterminée par \textbf{un point} et \textbf{une direction}, la direction étant donnée par un \cle{vecteur directeur}. C'est cette situation que nous étudions maintenant.

\courssub{Le problème posé}

\begin{definition}[Vecteur directeur d'une droite]
Soit $(d)$ une droite du plan. On appelle \cle{vecteur directeur} de $(d)$ tout vecteur \emph{non nul} $\vec{u}$ dont la direction est celle de la droite $(d)$, c'est-à-dire tel que $\vec{u} = \overrightarrow{AB}$ pour deux points $A$ et $B$ distincts de $(d)$.
\end{definition}

Une droite possède une infinité de vecteurs directeurs : ils sont tous \textbf{colinéaires} entre eux. Si $\vec{u}$ dirige $(d)$, alors $2\vec{u}$, $-3\vec{u}$, $\tfrac{1}{2}\vec{u}$ dirigent aussi $(d)$.

\textbf{Problème à résoudre.} Le plan est muni d'un repère $(O\,;\vec{\imath},\vec{\jmath})$. On donne :
\begin{itemize}
\item un point $A(x_A\,;y_A)$ ;
\item un vecteur $\vec{u}(\alpha\,;\beta)$ non nul.
\end{itemize}
On cherche une \textbf{équation cartésienne} de la droite $(d)$ passant par $A$ et dirigée par $\vec{u}$, c'est-à-dire une condition sur les coordonnées $(x\,;y)$ d'un point $M$ qui traduise exactement l'appartenance $M \in (d)$.

\courssub{Condition d'appartenance : la colinéarité}

L'idée fondamentale est la suivante : un point $M$ appartient à la droite $(d)$ si et seulement si le vecteur $\overrightarrow{AM}$ a la même direction que $\vec{u}$, autrement dit si et seulement si $\overrightarrow{AM}$ et $\vec{u}$ sont \textbf{colinéaires}. Or nous connaissons un critère algébrique de colinéarité : deux vecteurs $\vec{v}(x\,;y)$ et $\vec{v'}(x'\,;y')$ sont colinéaires si et seulement si $xy' - x'y = 0$.

\begin{propriete}[Équation d'une droite définie par un point et un vecteur directeur]
La droite $(d)$ passant par $A(x_A\,;y_A)$ et de vecteur directeur $\vec{u}(\alpha\,;\beta)$ admet pour équation :
\[ \beta\,(x - x_A) - \alpha\,(y - y_A) = 0. \]
En développant, on obtient une équation de la forme $ax + by + c = 0$ avec $a = \beta$, $b = -\alpha$ et $c = \alpha y_A - \beta x_A$.
\end{propriete}

\begin{experience}[Démonstration guidée]
Démontrons cette propriété pas à pas.
\begin{enumerate}
\item \textbf{Traduction géométrique.} Soit $M(x\,;y)$ un point du plan. Par définition d'un vecteur directeur :
\[ M \in (d) \iff \overrightarrow{AM} \text{ et } \vec{u} \text{ sont colinéaires.} \]
\item \textbf{Coordonnées de $\overrightarrow{AM}$.} On a $\overrightarrow{AM}(x - x_A\,;y - y_A)$.
\item \textbf{Critère de colinéarité.} Les vecteurs $\overrightarrow{AM}(x - x_A\,;y - y_A)$ et $\vec{u}(\alpha\,;\beta)$ sont colinéaires si et seulement si :
\[ (x - x_A)\times \beta - \alpha \times (y - y_A) = 0, \]
c'est-à-dire $\beta(x - x_A) - \alpha(y - y_A) = 0$.
\item \textbf{Mise sous forme développée.} En développant :
\[ \beta x - \beta x_A - \alpha y + \alpha y_A = 0 \iff \beta x - \alpha y + (\alpha y_A - \beta x_A) = 0, \]
qui est bien de la forme $ax + by + c = 0$ en posant $a = \beta$, $b = -\alpha$, $c = \alpha y_A - \beta x_A$. $\square$
\end{enumerate}
\end{experience}

\begin{aretenir}[Le lien fondamental]
Une droite d'équation $ax + by + c = 0$ admet pour vecteur directeur :
\[ \vec{u}(-b\,;a). \]
Réciproquement, si $\vec{u}(\alpha\,;\beta)$ dirige une droite, celle-ci a une équation de la forme $\beta x - \alpha y + c = 0$ (ou $-\beta x + \alpha y + c' = 0$, forme équivalente). Autrement dit : \textbf{les coefficients de $x$ et $y$ dans l'équation (sous la forme $\beta x - \alpha y + c = 0$) sont les coordonnées du vecteur directeur $\vec{u}(\alpha\,;\beta)$, échangées et avec un changement de signe sur la seconde}.
\end{aretenir}

\courssub{Exemple chiffré détaillé}

\begin{exemplebox}[Droite passant par $A(2\,;1)$ dirigée par $\vec{u}(-1\,;3)$]
Déterminons une équation de la droite $(d)$ passant par $A(2\,;1)$ et de vecteur directeur $\vec{u}(-1\,;3)$.

On applique la propriété avec $\alpha = -1$, $\beta = 3$, $x_A = 2$, $y_A = 1$ :
\begin{align*}
\beta(x - x_A) - \alpha(y - y_A) &= 0\\
3(x - 2) - (-1)(y - 1) &= 0\\
3(x - 2) + (y - 1) &= 0\\
3x - 6 + y - 1 &= 0\\
3x + y - 7 &= 0.
\end{align*}
La droite $(d)$ a pour équation $\boxed{3x + y - 7 = 0}$.

\textbf{Vérification.} Le point $A(2\,;1)$ doit satisfaire l'équation : $3\times 2 + 1 - 7 = 6 + 1 - 7 = 0$. C'est bien le cas. De plus, un vecteur directeur lu sur l'équation $3x + y - 7 = 0$ (avec $a=3, b=1$) est $\vec{u}(-b\,;a) = (-1\,;3)$ : on retrouve bien $\vec{u}$.
\end{exemplebox}

\begin{popfigure}
\begin{center}
\begin{tikzpicture}[scale=0.85,>=Stealth]
\clip (-1.4,-1.4) rectangle (4.4,4.4);
\draw[popline,thin] (-1.4,0) -- (4.4,0) node[below right] {$x$};
\draw[popline,thin] (0,-1.4) -- (0,4.4) node[above left] {$y$};
\foreach \x in {1,2,3,4} \draw[popline] (\x,0.06) -- (\x,-0.06) node[below,font=\small] {\x};
\foreach \y in {1,2,3,4} \draw[popline] (0.06,\y) -- (-0.06,\y) node[left,font=\small] {\y};
\node[below left,font=\small] at (0,0) {$O$};
% Droite 3x+y-7=0 => y = -3x+7. Intersections avec le cadre : x=0.86->y=4.4, x=2.8->y=-1.4
\draw[popblue,very thick,domain=0.8:2.9,samples=100] plot (\x,{-3*\x+7});
\fill[poporange] (2,1) circle (2.2pt) node[below right,popdark] {$A(2\,;1)$};
\draw[->,poporange,very thick] (2,1) -- (1,4) node[midway,left,popdark] {$\vec{u}(-1\,;3)$};
\fill[popgreen] (1,4) circle (2.2pt) node[above left,popdark] {$M(1\,;4)$};
\node[popblue,font=\small,align=center] at (3.6,1.6) {$(d)$ : $3x + y - 7 = 0$};
\end{tikzpicture}
\end{center}
\legende{La droite $(d)$ passe par $A(2\,;1)$ et est dirigée par $\vec{u}(-1\,;3)$ : son équation est $3x + y - 7 = 0$. Le point $M(1\,;4)$ appartient à $(d)$ car $3\times 1 + 4 - 7 = 0$.}
\end{popfigure}

\courssub{Lecture inverse : retrouver un vecteur directeur à partir d'une équation}

\begin{methode}[Passer de l'équation au vecteur directeur, et inversement]
\textbf{De l'équation vers le vecteur.} Si une droite a pour équation $ax + by + c = 0$ :
\begin{enumerate}
\item repérer les coefficients $a$ (devant $x$) et $b$ (devant $y$) ;
\item un vecteur directeur est $\vec{u}(-b\,;a)$.
\end{enumerate}
\textbf{Du vecteur vers l'équation.} Si $\vec{u}(\alpha\,;\beta)$ dirige la droite :
\begin{enumerate}
\item écrire l'équation sous la forme $\beta(x - x_A) - \alpha(y - y_A) = 0$ (ou $\beta x - \alpha y + c = 0$) ;
\item déterminer $c$ en écrivant qu'un point connu de la droite vérifie l'équation.
\end{enumerate}
\end{methode}

\begin{exoresolu}[Retrouver un vecteur directeur]
Soit $(d)$ la droite d'équation $5x - 2y + 8 = 0$. Donner un vecteur directeur de $(d)$, puis vérifier que le point $B(-2\,;-1)$ appartient à $(d)$.
\tcblower
\textbf{Solution.} Ici $a = 5$ et $b = -2$. Un vecteur directeur est donc :
\[ \vec{u}(-b\,;a) = (2\,;5). \]
Vérifions l'appartenance de $B(-2\,;-1)$ : $5\times(-2) - 2\times(-1) + 8 = -10 + 2 + 8 = 0$. Donc $B \in (d)$.

Remarque : tout vecteur colinéaire à $(2\,;5)$ convient aussi, par exemple $(-2\,;-5)$ ou $(4\,;10)$.
\end{exoresolu}

\begin{attention}[Erreur fréquente : coordonnées du vecteur et coefficients de l'équation]
L'erreur la plus classique est de confondre les coordonnées du vecteur directeur avec les coefficients $(a\,;b)$ de l'équation $ax + by + c = 0$. Retiens bien :
\begin{itemize}
\item le vecteur directeur est $\vec{u}(-b\,;a)$ et \textbf{non pas} $(a\,;b)$ ;
\item il y a un \textbf{échange} des deux nombres \textbf{et} un \textbf{changement de signe} sur le premier.
\end{itemize}
Par exemple, pour $3x + y - 7 = 0$, le vecteur directeur est $(-1\,;3)$, et non $(3\,;1)$. Pour te rassurer, vérifie toujours que le point donné satisfait l'équation obtenue : si ce n'est pas le cas, tu t'es trompé.
\end{attention}

\courssub{Cas particuliers : droites parallèles aux axes}

Que se passe-t-il lorsque le vecteur directeur est parallèle à l'un des axes du repère ?

\textbf{Premier cas : $\vec{u}$ parallèle à l'axe des ordonnées.} Alors $\alpha = 0$ et $\beta \neq 0$, par exemple $\vec{u}(0\,;1)$. L'équation devient :
\[ \beta(x - x_A) - 0 = 0 \iff x = x_A. \]
La droite est \textbf{verticale} : tous ses points ont la même abscisse. Son équation est $x = k$.

\textbf{Deuxième cas : $\vec{u}$ parallèle à l'axe des abscisses.} Alors $\beta = 0$ et $\alpha \neq 0$, par exemple $\vec{u}(1\,;0)$. L'équation devient :
\[ 0 - \alpha(y - y_A) = 0 \iff y = y_A. \]
La droite est \textbf{horizontale} : tous ses points ont la même ordonnée. Son équation est $y = k$.

\begin{aretenir}[À retenir sur les cas particuliers]
\begin{itemize}
\item Une droite d'équation $x = k$ est parallèle à l'axe des ordonnées (verticale) ; un vecteur directeur est $(0\,;1)$.
\item Une droite d'équation $y = k$ est parallèle à l'axe des abscisses (horizontale) ; un vecteur directeur est $(1\,;0)$.
\end{itemize}
\end{aretenir}

\courssub{Application : construire une droite dont on connaît l'équation}

Pour tracer une droite d'équation $ax + by + c = 0$ dans un repère, deux stratégies sont possibles.

\begin{methode}[Tracer une droite à partir de son équation]
\textbf{Stratégie 1 : deux points.} Choisir deux valeurs de $x$, calculer les $y$ correspondants, placer les deux points et tracer la droite.

\textbf{Stratégie 2 : un point et un vecteur directeur.}
\begin{enumerate}
\item trouver un point de la droite (par exemple en posant $x = 0$ et en calculant $y$) ;
\item lire un vecteur directeur $\vec{u}(-b\,;a)$ ;
\item placer le point, représenter $\vec{u}$ depuis ce point, puis tracer la droite.
\end{enumerate}
\end{methode}

\begin{exoresolu}[Construire une droite]
Tracer la droite $(d)$ d'équation $2x - 3y + 6 = 0$.
\tcblower
\textbf{Solution par deux points.} Si $x = 0$, alors $-3y + 6 = 0$, donc $y = 2$ : le point $P(0\,;2)$ est sur $(d)$. Si $x = 3$, alors $6 - 3y + 6 = 0$, donc $y = 4$ : le point $Q(3\,;4)$ est sur $(d)$. On place $P$ et $Q$ et on trace la droite.

\textbf{Solution par point et vecteur.} Un vecteur directeur est $\vec{u}(-b\,;a) = (3\,;2)$. Depuis $P(0\,;2)$, on se déplace de $3$ unités en abscisse et de $2$ unités en ordonnée : on tombe sur $Q(3\,;4)$. On retrouve la même droite.
\end{exoresolu}

\begin{savaistu}[Vecteur vitesse et trajectoire rectiligne]
En physique, lorsqu'un mobile se déplace \emph{en ligne droite à vitesse constante}, sa trajectoire est entièrement décrite par deux données : sa \textbf{position de départ} et son \textbf{vecteur vitesse}. Le vecteur vitesse joue exactement le rôle du vecteur directeur de la droite-trajectoire ! C'est ainsi qu'on prévoit la position d'un avion, d'un bateau ou d'un satellite : connaître le point et la direction suffit pour connaître toute la route. Les mathématiques que tu apprends en 3ème sont donc au cœur de la navigation et de la mécanique.
\end{savaistu}

\begin{exercice}[Pour s'entraîner]
\begin{enumerate}
\item Déterminer une équation de la droite passant par $C(-1\,;4)$ et de vecteur directeur $\vec{v}(2\,;-5)$.
\item La droite $(d')$ a pour équation $-x + 4y - 3 = 0$. Donner un vecteur directeur de $(d')$ et deux points de $(d')$.
\item La droite passant par $D(3\,;-2)$ et dirigée par $\vec{w}(0\,;7)$ : donner son équation sans calcul long. Justifier.
\end{enumerate}
\end{exercice}

\begin{corrige}[Exercice]
\begin{enumerate}
\item Avec $\alpha = 2$, $\beta = -5$ : $-5(x - (-1)) - 2(y - 4) = 0$, soit $-5x - 5 - 2y + 8 = 0$, d'où $-5x - 2y + 3 = 0$, ou encore $5x + 2y - 3 = 0$.
\item Un vecteur directeur est $\vec{u}(-b\,;a) = (-4\,;-1)$, ou plus simplement $(4\,;1)$. Points : pour $x = 0$, $y = \tfrac{3}{4}$ ; pour $y = 1$, $x = 1$. Donc $(0\,;\tfrac{3}{4})$ et $(1\,;1)$ conviennent.
\item Le vecteur $\vec{w}(0\,;7)$ est parallèle à l'axe des ordonnées : la droite est verticale, d'équation $x = 3$.
\end{enumerate}
\end{corrige}

\coursec{Équation réduite et coefficient directeur}

Dans la section précédente, nous avons appris à écrire l'équation cartésienne $ax + by + c = 0$ d'une droite dont on connaît un vecteur directeur. Cette forme est très pratique, mais elle a un défaut : une même droite possède une infinité d'équations cartésiennes (il suffit de tout multiplier par un même nombre non nul). Il serait donc agréable de disposer d'une écriture \emph{unique} et \emph{lisible} d'une droite, qui nous donne immédiatement ses caractéristiques. C'est exactement ce que fait l'\cle{équation réduite} : elle isole $y$ et fait apparaître deux nombres essentiels, le \cle{coefficient directeur} et l'\cle{ordonnée à l'origine}.

\courssub{De la forme cartésienne à la forme réduite}

\textbf{Intuition.} Imaginons un taxi-moto à Yaoundé qui facture une prise en charge fixe de $100$ FCFA, puis $200$ FCFA par kilomètre parcouru. Si $x$ désigne le nombre de kilomètres et $y$ le prix payé, on a tout de suite $y = 200x + 100$. Cette écriture est « réduite » : $y$ est seul d'un côté, et on lit directement les deux informations utiles. Partons maintenant d'une équation cartésienne quelconque et voyons comment obtenir cette forme.

\begin{definition}[Équation réduite d'une droite]
Toute droite $(d)$ \textbf{non verticale} (c'est-à-dire non parallèle à l'axe des ordonnées) admet une \textbf{unique} équation de la forme
\[ y = mx + p \]
appelée \cle{équation réduite} de la droite.
\begin{itemize}
\item Le nombre $m$ est appelé \cle{coefficient directeur} (ou \emph{pente}) de la droite.
\item Le nombre $p$ est appelé \cle{ordonnée à l'origine} : c'est l'ordonnée du point où la droite coupe l'axe des ordonnées, c'est-à-dire le point $(0\,; p)$.
\end{itemize}
\end{definition}

\begin{propriete}[Passage de la forme cartésienne à la forme réduite]
Soit $(d)$ une droite d'équation cartésienne $ax + by + c = 0$.
\begin{itemize}
\item Si $b \neq 0$, alors $(d)$ n'est pas verticale et son équation réduite est
\[ y = -\dfrac{a}{b}\,x - \dfrac{c}{b}, \qquad \text{donc } m = -\dfrac{a}{b} \ \text{ et } \ p = -\dfrac{c}{b}. \]
\item Si $b = 0$, alors $(d)$ est \textbf{verticale} : son équation s'écrit $x = k$ (avec $k = -\dfrac{c}{a}$) et elle \textbf{n'admet pas} d'équation réduite.
\end{itemize}
\end{propriete}

\begin{experience}[Démonstration guidée]
Partons de l'équation cartésienne $ax + by + c = 0$ avec $b \neq 0$.
\begin{enumerate}
\item On isole le terme contenant $y$ : $by = -ax - c$.
\item Comme $b \neq 0$, on peut diviser les deux membres par $b$ : $y = \dfrac{-ax - c}{b} = -\dfrac{a}{b}x - \dfrac{c}{b}$.
\item On pose $m = -\dfrac{a}{b}$ et $p = -\dfrac{c}{b}$ : on obtient bien $y = mx + p$.
\item Pour l'unicité : supposons qu'une même droite admette deux équations réduites $y = mx + p$ et $y = m'x + p'$. Pour tout $x$, on a $mx + p = m'x + p'$. En prenant $x = 0$, on obtient $p = p'$, puis $mx = m'x$ pour tout $x$, donc $m = m'$. L'écriture est donc unique.
\item Si $b = 0$, l'équation devient $ax + c = 0$, soit $x = -\dfrac{c}{a}$ : tous les points de la droite ont la même abscisse, la droite est verticale, et on ne peut pas exprimer $y$ en fonction de $x$.
\end{enumerate}
\end{experience}

\begin{exemplebox}[De la forme cartésienne à la forme réduite]
Soit $(d)$ la droite d'équation $4x - 2y + 6 = 0$. Ici $b = -2 \neq 0$, donc $(d)$ admet une équation réduite :
\begin{align*}
4x - 2y + 6 &= 0 \\
-2y &= -4x - 6 \\
y &= \dfrac{-4x - 6}{-2} = 2x + 3.
\end{align*}
Le coefficient directeur est $m = 2$ et l'ordonnée à l'origine est $p = 3$ : la droite coupe l'axe des ordonnées au point $(0\,; 3)$.
\end{exemplebox}

\begin{attention}[Droite verticale : pas d'équation réduite]
Ne cherche jamais à mettre une droite verticale sous la forme $y = mx + p$ : c'est impossible ! La droite d'équation $x = 5$, par exemple, est verticale (tous ses points ont pour abscisse $5$) et n'a ni coefficient directeur ni ordonnée à l'origine. Avant de diviser par $b$, vérifie toujours que $b \neq 0$.
\end{attention}

\courssub{Calcul du coefficient directeur à partir de deux points}

\textbf{Intuition.} Le coefficient directeur mesure « l'inclinaison » de la droite, comme la pente d'une route entre deux villages. Pour la calculer, on compare la variation des ordonnées (le « dénivelé ») à la variation des abscisses (la « distance horizontale »).

\begin{propriete}[Coefficient directeur à partir de deux points]
Soient $A(x_A\,; y_A)$ et $B(x_B\,; y_B)$ deux points distincts d'une droite $(d)$ non verticale (donc $x_A \neq x_B$). Alors le coefficient directeur de $(d)$ est
\[ m = \dfrac{y_B - y_A}{x_B - x_A}. \]
\end{propriete}

\begin{experience}[Démonstration]
Comme $A$ et $B$ appartiennent à la droite $(d)$ d'équation réduite $y = mx + p$, leurs coordonnées vérifient cette équation :
\[ y_A = m x_A + p \qquad \text{et} \qquad y_B = m x_B + p. \]
En soustrayant membre à membre :
\[ y_B - y_A = m x_B - m x_A = m(x_B - x_A). \]
Comme la droite n'est pas verticale, $x_B - x_A \neq 0$, et on peut diviser :
\[ m = \dfrac{y_B - y_A}{x_B - x_A}. \]
\end{experience}

\begin{exemplebox}[Calcul d'un coefficient directeur]
Soient $A(1\,; 2)$ et $B(4\,; 8)$. Comme $x_A \neq x_B$, la droite $(AB)$ n'est pas verticale et
\[ m = \dfrac{y_B - y_A}{x_B - x_A} = \dfrac{8 - 2}{4 - 1} = \dfrac{6}{3} = 2. \]
La droite $(AB)$ a pour coefficient directeur $2$.
\end{exemplebox}

\begin{attention}[L'erreur classique : inverser les variations]
L'erreur la plus fréquente est de calculer $m = \dfrac{x_B - x_A}{y_B - y_A}$ (les abscisses en haut). C'est \textbf{faux} ! Retiens bien la formule : c'est toujours
\[ m = \dfrac{\text{variation des ordonnées}}{\text{variation des abscisses}} = \dfrac{y_B - y_A}{x_B - x_A}. \]
Autre piège : il faut prendre les points \emph{dans le même ordre} en haut et en bas. On peut écrire $\dfrac{y_B - y_A}{x_B - x_A}$ ou $\dfrac{y_A - y_B}{x_A - x_B}$, mais jamais $\dfrac{y_B - y_A}{x_A - x_B}$.
\end{attention}

\courssub{Interprétation graphique du coefficient directeur}

\begin{aretenir}[Ce que dit le coefficient directeur]
Sur une droite d'équation $y = mx + p$ :
\begin{itemize}
\item quand $x$ augmente de $1$, alors $y$ varie de $m$ (il augmente de $m$ si $m > 0$, il diminue de $-m$ si $m < 0$) ;
\item si $m > 0$, la droite « monte » de gauche à droite ;
\item si $m < 0$, la droite « descend » de gauche à droite ;
\item si $m = 0$, la droite est horizontale (équation $y = p$) ;
\item plus $|m|$ est grand, plus la droite est « raide ».
\end{itemize}
L'ordonnée à l'origine $p$ se lit directement : c'est l'ordonnée du point d'intersection de la droite avec l'axe des ordonnées.
\end{aretenir}

\begin{popfigure}
\begin{center}
\begin{tikzpicture}
\begin{axis}[width=0.85\linewidth, axis lines=middle, xlabel={$x$}, ylabel={$y$}, xmin=-1.5, xmax=3.5, ymin=-3, ymax=6, xtick={-1,0,1,2,3}, ytick={-2,-1,1,2,3,4,5}, grid=both, grid style={popline!40}, legend style={at={(0.02,0.98)}, anchor=north west}]
\addplot[popcurve, domain=-1:3.2, samples=100]{2*x - 1};
\addlegendentry{$y = 2x - 1$}
\addplot[mark=*, mark size=1.5pt, popdark] coordinates {(0,-1) (1,1) (2,3)};
\addplot[dashed, poporange, thick] coordinates {(1,1) (2,1) (2,3)};
\node[poporange, anchor=south] at (axis cs:1.5,1) {$+1$};
\node[poporange, anchor=west] at (axis cs:2,2) {$+2$};
\node[popdark, anchor=east] at (axis cs:0,-1) {$(0\,; p)$};
\end{axis}
\end{tikzpicture}
\end{center}
\legende{La droite $y = 2x - 1$ : le triangle de pente montre que quand $x$ augmente de $1$, $y$ augmente de $m = 2$. L'ordonnée à l'origine est $p = -1$.}
\end{popfigure}

\begin{popfigure}
\begin{center}
\begin{tikzpicture}
\begin{axis}[width=0.85\linewidth, axis lines=middle, xlabel={$x$}, ylabel={$y$}, xmin=-3.5, xmax=3.5, ymin=-3, ymax=5, xtick={-3,-2,-1,0,1,2,3}, ytick={-2,-1,1,2,3,4}, grid=both, grid style={popline!40}, legend style={at={(0.02,0.98)}, anchor=north west}]
\addplot[popcurve, domain=-2:2, samples=50]{2*x + 1};
\addlegendentry{$y = 2x + 1$}
\addplot[popcurve, poporange, domain=-2:3.2, samples=50]{-x + 1};
\addlegendentry{$y = -x + 1$}
\addplot[popcurve, popgreen, domain=-3:3.2, samples=50]{0.5*x + 1};
\addlegendentry{$y = 0{,}5x + 1$}
\addplot[mark=*, mark size=1.8pt, popdark] coordinates {(0,1)};
\node[popdark, anchor=west] at (axis cs:0.1,0.6) {$(0\,; 1)$};
\end{axis}
\end{tikzpicture}
\end{center}
\legende{Trois droites de même ordonnée à l'origine $p = 1$ mais de coefficients directeurs différents : $m = 2$ (monte vite), $m = -1$ (descend), $m = 0{,}5$ (monte doucement).}
\end{popfigure}

\begin{methode}[Lire graphiquement $m$ et $p$ sur une droite tracée]
\begin{enumerate}
\item \textbf{Lire $p$} : repérer le point où la droite coupe l'axe des ordonnées ; son ordonnée est $p$.
\item \textbf{Choisir deux points} de la droite dont les coordonnées sont entières (des « nœuds » du quadrillage), par exemple $A$ et $B$.
\item \textbf{Compter les variations} : en allant de $A$ vers $B$, compter de combien on avance horizontalement ($x_B - x_A$) et de combien on monte ou descend verticalement ($y_B - y_A$).
\item \textbf{Calculer} $m = \dfrac{y_B - y_A}{x_B - x_A}$, en comptant négativement une descente.
\item \textbf{Conclure} : l'équation réduite est $y = mx + p$.
\end{enumerate}
\end{methode}

\courssub{Droite de coefficient directeur donné passant par un point donné}

\textbf{Intuition.} Une droite est entièrement déterminée par sa pente et un point : c'est comme une route dont on connaît l'inclinaison et un point de passage. Il existe une formule directe pour obtenir son équation.

\begin{propriete}[Équation de la droite de coefficient directeur $m$ passant par $A$]
La droite de coefficient directeur $m$ passant par le point $A(x_A\,; y_A)$ a pour équation
\[ y - y_A = m(x - x_A). \]
\end{propriete}

\begin{experience}[Démonstration]
La droite cherchée a une équation réduite de la forme $y = mx + p$, puisque son coefficient directeur est $m$. Comme $A$ appartient à cette droite, ses coordonnées vérifient l'équation : $y_A = m x_A + p$, donc $p = y_A - m x_A$. En remplaçant dans $y = mx + p$ :
\[ y = mx + y_A - m x_A \quad \Longleftrightarrow \quad y - y_A = m(x - x_A). \]
\end{experience}

\begin{methode}[Déterminer l'équation réduite d'une droite de coefficient directeur $m$ passant par $A$]
\begin{enumerate}
\item Écrire la formule $y - y_A = m(x - x_A)$.
\item Remplacer $m$, $x_A$ et $y_A$ par leurs valeurs.
\item Développer et réduire le membre de droite.
\item Isoler $y$ pour obtenir la forme $y = mx + p$.
\item \textbf{Vérifier} : remplacer $x$ par $x_A$ dans l'équation obtenue ; on doit retrouver $y_A$.
\end{enumerate}
\end{methode}

\begin{exemplebox}[Coefficient directeur $-2$, passage par $A(3\,; 1)$]
Cherchons l'équation réduite de la droite $(d)$ de coefficient directeur $m = -2$ passant par $A(3\,; 1)$.
\begin{align*}
y - y_A &= m(x - x_A) \\
y - 1 &= -2(x - 3) \\
y - 1 &= -2x + 6 \\
y &= -2x + 7.
\end{align*}
L'équation réduite de $(d)$ est $y = -2x + 7$. Vérification : pour $x = 3$, $y = -2 \times 3 + 7 = 1$ : le point $A$ est bien sur la droite.
\end{exemplebox}

\begin{exoresolu}[D'un vecteur directeur à l'équation réduite]
Énoncé. Dans un repère orthonormé, on considère la droite $(d)$ de vecteur directeur $\vec{u}(-1\,; 3)$ passant par le point $C(2\,; -1)$. Déterminer l'équation réduite de $(d)$.
\tcblower
Solution détaillée. Un vecteur directeur $\vec{u}(-1\,; 3)$ donne la variation : quand $x$ varie de $-1$, $y$ varie de $3$. Le coefficient directeur est donc
\[ m = \dfrac{3}{-1} = -3. \]
(On retrouve la règle : pour un vecteur directeur $\vec{u}(\alpha\,; \beta)$ avec $\alpha \neq 0$, $m = \dfrac{\beta}{\alpha}$.) La droite passe par $C(2\,; -1)$, donc :
\begin{align*}
y - (-1) &= -3(x - 2) \\
y + 1 &= -3x + 6 \\
y &= -3x + 5.
\end{align*}
L'équation réduite de $(d)$ est $y = -3x + 5$. Vérification : pour $x = 2$, $y = -3 \times 2 + 5 = -1$ : le point $C$ appartient bien à $(d)$.
\end{exoresolu}

\begin{exoresolu}[Lecture graphique et modélisation]
Énoncé. Un opérateur de transport interurbain affiche un prix de course composé d'une prise en charge fixe et d'un montant proportionnel à la distance. Pour $10$ km, on paie $\num{2500}$ FCFA ; pour $30$ km, on paie $\num{6500}$ FCFA. On note $x$ la distance en km et $y$ le prix en FCFA, et on admet que $y = mx + p$. Déterminer $m$ et $p$, puis le prix d'une course de $50$ km.
\tcblower
Solution détaillée. Les points $A(10\,; \num{2500})$ et $B(30\,; \num{6500})$ appartiennent à la droite.
\[ m = \dfrac{y_B - y_A}{x_B - x_A} = \dfrac{\num{6500} - \num{2500}}{30 - 10} = \dfrac{\num{4000}}{20} = 200. \]
La droite passe par $A(10\,; \num{2500})$ :
\[ y - \num{2500} = 200(x - 10) \quad \Longleftrightarrow \quad y = 200x + 500. \]
Donc $m = 200$ et $p = 500$. Pour une course de $50$ km : $y = 200 \times 50 + 500 = \num{10500}$ FCFA.
\end{exoresolu}

\begin{savaistu}[Le coefficient directeur en économie]
En économie, le coefficient directeur d'une droite de coût représente le \emph{coût unitaire}, et l'ordonnée à l'origine le \emph{coût fixe}. Par exemple, si une agence de voyage facture $y = 500x + \num{2000}$ (en FCFA) pour un trajet de $x$ kilomètres, alors $500$ FCFA est le prix par kilomètre et $\num{2000}$ FCFA est la prise en charge, payée même si l'on parcourt $0$ km. Les économistes, les transporteurs de Douala à Bafoussam et même les opérateurs de téléphonie mobile utilisent en permanence ce modèle $y = mx + p$ pour prévoir leurs recettes !
\end{savaistu}

\begin{exercice}[Pour s'entraîner]
\begin{enumerate}
\item Déterminer l'équation réduite de la droite d'équation cartésienne $3x + 6y - 12 = 0$.
\item Calculer le coefficient directeur de la droite passant par $M(-2\,; 5)$ et $N(4\,; -1)$.
\item Déterminer l'équation réduite de la droite de coefficient directeur $\dfrac{3}{4}$ passant par $P(8\,; 2)$.
\item La droite $(\Delta)$ a pour équation $y = -x + 4$. Le point $Q(7\,; -3)$ appartient-il à $(\Delta)$ ?
\end{enumerate}
\end{exercice}

\begin{corrige}[Exercice 1]
\begin{enumerate}
\item $3x + 6y - 12 = 0 \Leftrightarrow 6y = -3x + 12 \Leftrightarrow y = -\dfrac{1}{2}x + 2$. Donc $m = -\dfrac{1}{2}$ et $p = 2$.
\item $m = \dfrac{-1 - 5}{4 - (-2)} = \dfrac{-6}{6} = -1$.
\item $y - 2 = \dfrac{3}{4}(x - 8) \Leftrightarrow y = \dfrac{3}{4}x - 6 + 2 \Leftrightarrow y = \dfrac{3}{4}x - 4$.
\item Pour $x = 7$ : $y = -7 + 4 = -3$. C'est bien l'ordonnée de $Q$, donc $Q \in (\Delta)$.
\end{enumerate}
\end{corrige}

\begin{aretenir}[L'essentiel de la section]
\begin{itemize}
\item Toute droite \textbf{non verticale} admet une unique équation réduite $y = mx + p$ : $m$ est le coefficient directeur, $p$ l'ordonnée à l'origine.
\item Passage de la forme cartésienne : si $b \neq 0$, $ax + by + c = 0 \Leftrightarrow y = -\dfrac{a}{b}x - \dfrac{c}{b}$. Si $b = 0$, la droite est verticale : $x = k$, pas d'équation réduite.
\item Coefficient directeur à partir de deux points : $m = \dfrac{y_B - y_A}{x_B - x_A}$ (variation des ordonnées sur variation des abscisses).
\item Quand $x$ augmente de $1$, $y$ varie de $m$ ; $m > 0$ : la droite monte, $m < 0$ : elle descend, $m = 0$ : elle est horizontale.
\item Droite de coefficient directeur $m$ passant par $A(x_A\,; y_A)$ : $y - y_A = m(x - x_A)$.
\end{itemize}
\end{aretenir}

\coursec{Positions relatives de deux droites}

Dans la section précédente, nous avons appris à écrire l'équation réduite $y = mx + p$ d'une droite et à lire son \cle{coefficient directeur} $m$, qui mesure son « inclinaison ». Une question naturelle se pose alors : si l'on trace \emph{deux} droites dans le même plan, comment savoir, \emph{sans même les dessiner}, si elles se croisent, si elles sont parallèles, ou si elles sont superposées ? C'est exactement le problème que se pose un ingénieur de CAMTEL qui vérifie si deux lignes de câbles se croisent, ou un conducteur qui veut savoir si deux routes vont se rencontrer. Nous allons voir que la réponse tient en quelques calculs très simples.

\courssub{Les trois positions possibles}

\begin{experience}[Activité de découverte]
Prends une feuille et une règle. Trace deux droites quelconques $(d)$ et $(d')$. Recommence plusieurs fois. Tu constateras qu'il n'existe que \textbf{trois situations possibles} :
\begin{enumerate}
\item les deux droites se coupent en \textbf{un seul point} : on dit qu'elles sont \cle{sécantes} ;
\item les deux droites ne se coupent \textbf{jamais} et gardent la même direction : on dit qu'elles sont \cle{strictement parallèles} ;
\item les deux droites sont \textbf{superposées} (tous leurs points sont communs) : on dit qu'elles sont \cle{confondues}.
\end{enumerate}
\end{experience}

\begin{aretenir}[Les trois positions relatives de deux droites]
Dans le plan, deux droites $(d)$ et $(d')$ sont dans \textbf{l'une et une seule} des trois situations suivantes :
\begin{center}
\begin{tabularx}{\linewidth}{|l|X|X|}
\hline
\rowcolor{popblueL} \textbf{Position} & \textbf{Nombre de points communs} & \textbf{Vocabulaire} \\
\hline
Sécantes & exactement $1$ point & elles se coupent \\
\hline
Strictement parallèles & $0$ point & même direction, distinctes \\
\hline
Confondues & une infinité & ce sont la même droite \\
\hline
\end{tabularx}
\end{center}
\end{aretenir}

\begin{popfigure}
\begin{center}
\begin{tikzpicture}[scale=0.85]
% Panneau (a) : sécantes
\begin{scope}
\draw[popline,->] (-2.2,0)--(2.2,0) node[below left]{\footnotesize $x$};
\draw[popline,->] (0,-1.6)--(0,2.6) node[below left]{\footnotesize $y$};
\draw[popblue,thick,domain=-1.8:1.8] plot(\x,{0.8*\x+0.4});
\draw[poporange,thick,domain=-1.8:1.8] plot(\x,{-0.9*\x+1.4});
\fill[popdark] (0.588,0.87) circle (2.2pt) node[above right]{\footnotesize $I$};
\node[popblue] at (1.7,2.0) {\footnotesize $(d)$};
\node[poporange] at (1.7,-0.4) {\footnotesize $(d')$};
\node at (0,-2.1) {\footnotesize \textbf{(a)} Sécantes en $I$};
\end{scope}
% Panneau (b) : strictement parallèles
\begin{scope}[xshift=6.2cm]
\draw[popline,->] (-2.2,0)--(2.2,0) node[below left]{\footnotesize $x$};
\draw[popline,->] (0,-1.6)--(0,2.6) node[below left]{\footnotesize $y$};
\draw[popblue,thick,domain=-1.8:1.8] plot(\x,{0.7*\x+1.2});
\draw[poporange,thick,domain=-1.8:1.8] plot(\x,{0.7*\x-0.6});
\node[popblue] at (1.6,2.6) {\footnotesize $(d)$};
\node[poporange] at (1.6,0.9) {\footnotesize $(d')$};
\node at (0,-2.1) {\footnotesize \textbf{(b)} Strictement parallèles};
\end{scope}
% Panneau (c) : confondues
\begin{scope}[xshift=12.4cm]
\draw[popline,->] (-2.2,0)--(2.2,0) node[below left]{\footnotesize $x$};
\draw[popline,->] (0,-1.6)--(0,2.6) node[below left]{\footnotesize $y$};
\draw[popblue,thick,domain=-1.8:1.8] plot(\x,{0.7*\x+0.3});
\draw[poporange,thick,dashed,domain=-1.75:1.75] plot(\x,{0.7*\x+0.35});
\node[popblue] at (1.6,1.9) {\footnotesize $(d)=(d')$};
\node at (0,-2.1) {\footnotesize \textbf{(c)} Confondues};
\end{scope}
\end{tikzpicture}
\end{center}
\legende{Les trois positions relatives de deux droites du plan : (a) un point commun, (b) aucun point commun, (c) tous les points communs.}
\end{popfigure}

\courssub{Critère de parallélisme par les vecteurs directeurs}

L'idée intuitive est simple : deux droites sont parallèles (strictement ou confondues) exactement lorsqu'elles ont \textbf{la même direction}. Or la direction d'une droite est donnée par ses vecteurs directeurs. Deux directions sont identiques lorsque les vecteurs sont \cle{colinéaires}.

\begin{propriete}[Théorème : critère de parallélisme (forme cartésienne)]
Soient $(d) : ax + by + c = 0$ et $(d') : a'x + b'y + c' = 0$ deux droites du plan.
\[ (d) \text{ et } (d') \text{ sont parallèles (ou confondues)} \iff ab' - a'b = 0. \]
Si $ab' - a'b \neq 0$, alors les deux droites sont \textbf{sécantes en un unique point}.
\end{propriete}

\begin{experience}[Démonstration guidée du théorème]
Rappelons d'abord qu'un vecteur directeur de $(d) : ax + by + c = 0$ est $\vec{u}(-b\,;\,a)$, et un vecteur directeur de $(d')$ est $\vec{u'}(-b'\,;\,a')$.
\begin{enumerate}
\item \textbf{Traduction géométrique.} $(d)$ et $(d')$ sont parallèles si et seulement si leurs vecteurs directeurs $\vec{u}$ et $\vec{u'}$ sont \cle{colinéaires}.
\item \textbf{Critère de colinéarité.} Deux vecteurs $\vec{u}(x\,;\,y)$ et $\vec{u'}(x'\,;\,y')$ sont colinéaires si et seulement si $xy' - x'y = 0$.
\item \textbf{Application.} Ici $x = -b$, $y = a$, $x' = -b'$, $y' = a'$. Donc :
\[ xy' - x'y = (-b)\times a' - (-b')\times a = -a'b + ab' = ab' - a'b. \]
\item \textbf{Conclusion.} $\vec{u}$ et $\vec{u'}$ sont colinéaires si et seulement si $ab' - a'b = 0$, ce qui prouve le théorème. \hfill$\blacksquare$
\end{enumerate}
\end{experience}

\begin{exemplebox}[Application directe du critère $ab' - a'b$]
Soient $(d) : 2x - 3y + 5 = 0$ et $(d') : 4x - 6y + 1 = 0$. Sont-elles parallèles ?

On identifie $a=2,\; b=-3$ et $a'=4,\; b'=-6$.
On calcule : $ab' - a'b = 2\times(-6) - 4\times(-3) = -12 + 12 = 0$.

Le critère est nul : les droites sont parallèles (reste à savoir si elles sont confondues ou strictement parallèles : voir plus bas).
\end{exemplebox}

\courssub{Critère par les coefficients directeurs}

Lorsque les deux droites ne sont \textbf{pas verticales}, on peut écrire leurs équations réduites $(d) : y = mx + p$ et $(d') : y = m'x + p'$. Le critère devient alors d'une simplicité remarquable.

\begin{propriete}[Positions relatives et équations réduites]
Soient $(d) : y = mx + p$ et $(d') : y = m'x + p'$ deux droites non verticales.
\begin{itemize}
\item $(d)$ et $(d')$ sont \textbf{sécantes} $\iff m \neq m'$ ;
\item $(d)$ et $(d')$ sont \textbf{strictement parallèles} $\iff m = m'$ \textbf{et} $p \neq p'$ ;
\item $(d)$ et $(d')$ sont \textbf{confondues} $\iff m = m'$ \textbf{et} $p = p'$.
\end{itemize}
\end{propriete}

\begin{experience}[Démonstration guidée : pourquoi « parallèles $\iff m = m'$ » ?]
Une droite d'équation réduite $y = mx + p$ admet pour vecteur directeur $\vec{u}(1\,;\,m)$ (quand $x$ augmente de $1$, $y$ augmente de $m$). De même, $(d')$ admet $\vec{u'}(1\,;\,m')$.
\begin{enumerate}
\item $(d)$ et $(d')$ parallèles $\iff \vec{u}(1\,;\,m)$ et $\vec{u'}(1\,;\,m')$ colinéaires.
\item Critère de colinéarité : $1\times m' - 1\times m = 0$, c'est-à-dire $m' - m = 0$.
\item Donc $(d) \parallel (d') \iff m = m'$. \hfill$\blacksquare$
\end{enumerate}
De plus, si $m = m'$ et $p = p'$, les deux équations réduites sont \emph{identiques} : c'est la même droite, elles sont confondues. Si $m = m'$ mais $p \neq p'$, les droites ont même direction mais passent par des points différents de l'axe des ordonnées : elles sont strictement parallèles.
\end{experience}

\courssub{Distinguer « strictement parallèles » et « confondues »}

Quand $m = m'$ (ou $ab' - a'b = 0$), le travail n'est \textbf{pas terminé} : il faut départager les deux cas. Deux stratégies équivalentes :

\begin{methode}[Départager parallèles strictes et droites confondues]
\begin{enumerate}
\item \textbf{Méthode 1 (ordonnées à l'origine).} Écrire les deux équations réduites. Si $p = p'$, les droites sont confondues ; si $p \neq p'$, elles sont strictement parallèles.
\item \textbf{Méthode 2 (test d'un point).} Choisir un point de $(d)$ (par exemple le point d'abscisse $0$), et tester si ses coordonnées vérifient l'équation de $(d')$. Si oui, les droites sont confondues ; si non, elles sont strictement parallèles.
\end{enumerate}
\end{methode}

\begin{exemplebox}[Exemple chiffré]
Reprenons $(d) : 2x - 3y + 5 = 0$ et $(d') : 4x - 6y + 1 = 0$ (on a vu que $ab' - a'b = 0$).

Équations réduites : $(d) : y = \dfrac{2}{3}x + \dfrac{5}{3}$ et $(d') : y = \dfrac{2}{3}x + \dfrac{1}{6}$.

On a $m = m' = \dfrac{2}{3}$ mais $p = \dfrac{5}{3} \neq \dfrac{1}{6} = p'$ : les droites sont \textbf{strictement parallèles}.

\emph{Vérification par un point :} $A(0\,;\,\tfrac{5}{3}) \in (d)$. Testons dans $(d')$ : $4\times 0 - 6\times\tfrac{5}{3} + 1 = -10 + 1 = -9 \neq 0$. Donc $A \notin (d')$ : confirmé, les droites sont distinctes.
\end{exemplebox}

\begin{attention}[Erreur fréquente : conclure trop vite]
L'erreur la plus classique au BEPC est de conclure « les droites sont parallèles » dès que $m = m'$, \textbf{sans comparer $p$ et $p'$}. Or si $p = p'$, les droites sont \emph{confondues}, pas strictement parallèles ! Rédige toujours la comparaison complète : « $m = m' = \ldots$ et $p = \ldots \neq \ldots = p'$, donc les droites sont strictement parallèles ».
\end{attention}

\courssub{Point d'intersection de deux droites sécantes}

Lorsque $m \neq m'$, les droites se coupent en un unique point $I$. Ses coordonnées vérifient \textbf{les deux équations à la fois} : on résout donc le système formé par les deux équations.

\begin{methode}[Trouver le point d'intersection de deux droites sécantes (forme réduite)]
\begin{enumerate}
\item Vérifier que les droites sont sécantes : $m \neq m'$ (ou $ab' - a'b \neq 0$).
\item Écrire le système formé par les deux équations réduites : $\begin{cases} y = mx + p \\ y = m'x + p' \end{cases}$ (si une droite est verticale, utiliser son équation cartésienne $x = k$).
\item Résoudre par \textbf{substitution} (égaliser $mx + p = m'x + p'$) ou par \textbf{combinaison}.
\item Reporter la valeur trouvée dans l'une des équations pour obtenir la deuxième coordonnée.
\item Conclure : $I(x_I\,;\,y_I)$, et éventuellement vérifier dans l'autre équation.
\end{enumerate}
\end{methode}

\begin{exoresolu}[Intersection de $(d) : y = 2x + 1$ et $(d') : y = -x + 4$]
Déterminer la position relative des droites $(d) : y = 2x + 1$ et $(d') : y = -x + 4$, puis leurs éventuels points communs.
\tcblower
\textbf{Étape 1 : position relative.} Les coefficients directeurs sont $m = 2$ et $m' = -1$. Comme $m \neq m'$, les droites sont \textbf{sécantes en un unique point} $I$.

\textbf{Étape 2 : résolution du système.} Les coordonnées de $I$ vérifient les deux équations :
\[ \begin{cases} y = 2x + 1 \\ y = -x + 4 \end{cases} \]
Par substitution, on égalise les deux expressions de $y$ :
\begin{align*}
2x + 1 &= -x + 4 \\
2x + x &= 4 - 1 \\
3x &= 3 \\
x &= 1
\end{align*}

\textbf{Étape 3 : calcul de $y$.} On reporte $x = 1$ dans la première équation : $y = 2\times 1 + 1 = 3$.

\textbf{Étape 4 : vérification.} Dans $(d')$ : $y = -1 + 4 = 3$. \checkmark

\textbf{Conclusion.} Les droites $(d)$ et $(d')$ sont sécantes au point $I(1\,;\,3)$.
\end{exoresolu}

\courssub{Lien avec le nombre de solutions d'un système}

Résoudre un système de deux équations à deux inconnues, c'est exactement chercher les points communs à deux droites. Les trois positions relatives correspondent donc aux trois nombres de solutions possibles :

\begin{center}
\begin{tabularx}{\linewidth}{|l|X|X|X|}
\hline
\rowcolor{popblueL} \textbf{Condition} & \textbf{Position des droites} & \textbf{Points communs} & \textbf{Solutions du système} \\
\hline
$m \neq m'$ & sécantes & $1$ point & une unique solution \\
\hline
$m = m'$ et $p \neq p'$ & strictement parallèles & aucun & aucune solution \\
\hline
$m = m'$ et $p = p'$ & confondues & une infinité & une infinité de solutions \\
\hline
\end{tabularx}
\end{center}

\begin{popfigure}
\begin{center}
\begin{tikzpicture}[
  node distance=0.9cm and 1.6cm,
  decision/.style={diamond, draw=popdark, fill=popgoldL, text width=2.6cm, align=center, inner sep=1pt, font=\footnotesize},
  result/.style={rectangle, rounded corners, draw=popdark, fill=popblueL, text width=3.4cm, align=center, font=\footnotesize, inner sep=4pt},
  lab/.style={font=\footnotesize\itshape, popdark}
]
\node[decision] (q1) {Comparer $m$ et $m'$};
\node[result, below left=1.1cm and 0.4cm of q1] (sec) {$(d)$ et $(d')$ \textbf{sécantes} : un unique point d'intersection};
\node[decision, below right=1.1cm and 0.4cm of q1] (q2) {Comparer $p$ et $p'$};
\node[result, below left=1.1cm and -0.6cm of q2] (conf) {$(d)$ et $(d')$ \textbf{confondues}};
\node[result, below right=1.1cm and -0.6cm of q2] (par) {$(d)$ et $(d')$ \textbf{strictement parallèles}};
\draw[-{Stealth}, thick, popdark] (q1) -- node[lab, above left]{$m \neq m'$} (sec);
\draw[-{Stealth}, thick, popdark] (q1) -- node[lab, above right]{$m = m'$} (q2);
\draw[-{Stealth}, thick, popdark] (q2) -- node[lab, above left]{$p = p'$} (conf);
\draw[-{Stealth}, thick, popdark] (q2) -- node[lab, above right]{$p \neq p'$} (par);
\end{tikzpicture}
\end{center}
\legende{Arbre de décision pour déterminer la position relative de deux droites non verticales $y = mx + p$ et $y = m'x + p'$.}
\end{popfigure}

\begin{savaistu}[Pourquoi un système a 0, 1 ou une infinité de solutions]
En résolvant des systèmes d'équations, tu as peut-être déjà rencontré des systèmes « impossibles » (on aboutit à $0 = 5$) ou « indéterminés » (on aboutit à $0 = 0$). La géométrie explique tout : chaque équation représente une droite, et résoudre le système revient à chercher les points communs aux deux droites. Deux droites du plan ne peuvent avoir que $0$, $1$ ou une infinité de points communs : il est donc \emph{impossible} qu'un tel système ait exactement 2 ou 3 solutions ! Ce beau pont entre l'algèbre et la géométrie est au cœur de ce qu'on appelle la \emph{géométrie analytique}, inventée par René Descartes au XVII\textsuperscript{e} siècle.
\end{savaistu}

\begin{exercice}[À toi de jouer]
Dans chaque cas, déterminer la position relative des droites $(d)$ et $(d')$, puis les coordonnées du point d'intersection s'il existe.
\begin{enumerate}
\item $(d) : y = 3x - 2$ et $(d') : y = x + 4$.
\item $(d) : y = -2x + 5$ et $(d') : y = -2x - 1$.
\item $(d) : 2x + y - 3 = 0$ et $(d') : y = -2x + 3$.
\end{enumerate}
\end{exercice}

\begin{corrige}[Exercice : positions relatives]
\begin{enumerate}
\item $m = 3 \neq m' = 1$ : droites sécantes. On résout $3x - 2 = x + 4$, soit $2x = 6$, $x = 3$, puis $y = 3 + 4 = 7$. Intersection : $I(3\,;\,7)$.
\item $m = m' = -2$ et $p = 5 \neq -1 = p'$ : droites \textbf{strictement parallèles}, pas de point d'intersection.
\item Équation réduite de $(d)$ : $y = -2x + 3$. Donc $m = m' = -2$ et $p = p' = 3$ : les droites sont \textbf{confondues} (une infinité de points communs).
\end{enumerate}
\end{corrige}

\coursec{Applications et modélisation}

Après avoir étudié les positions relatives de deux droites (parallèles, confondues, sécantes), nous disposons maintenant de tous les outils pour \textbf{modéliser des situations concrètes} par des équations de droites, les \textbf{exploiter} pour prendre des décisions (choisir un forfait, prévoir un coût, trouver un seuil de rentabilité) et \textbf{rédiger une démarche complète} répondant aux exigences du savoir-faire « Rédiger et justifier ». C'est l'aboutissement naturel de cette leçon : les mathématiques deviennent un outil d'aide à la décision.

\courssub{Modéliser une situation de proportionnalité ou d'évolution affine}

De nombreuses situations de la vie courante font intervenir deux grandeurs liées par une relation de la forme $y = mx + p$ (fonction affine) ou $y = mx$ (proportionnalité, cas particulier où $p=0$).
\begin{itemize}
    \item \textbf{Transport :} coût total en fonction de la distance (prise en charge + prix au km).
    \item \textbf{Télécommunications (MTN, Orange) :} coût d'un forfait en fonction de la durée d'appels ou du volume de data au-delà du forfait de base.
    \item \textbf{Électricité / Eau (ENEO, CAMWATER) :} facture = abonnement fixe + prix au kWh ou au m$^3$ consommé.
    \item \textbf{Commerce :} salaire d'un vendeur = fixe + commission sur ventes.
\end{itemize}

La modélisation consiste à \textbf{traduire le texte en langage mathématique} :
\begin{enumerate}
    \item Identifier la \textbf{variable indépendante} $x$ (ce qui varie librement : distance, consommation, heures...) et la \textbf{variable dépendante} $y$ (ce qui en dépend : prix total, salaire, facture...).
    \item Repérer le \textbf{taux de variation} $m$ (prix unitaire, commission, coût au km...). C'est le \textbf{coefficient directeur} de la droite.
    \item Repérer la \textbf{valeur initiale} $p$ (abonnement, prise en charge, salaire de base...) quand $x=0$. C'est l'\textbf{ordonnée à l'origine}.
    \item Écrire l'équation : $\mathbf{y = mx + p}$.
\end{enumerate}

\begin{methode}[Modéliser par une équation de droite]
\textbf{Étape 1 :} Choisir les inconnues. $x$ = grandeur qui varie (distance, quantité, temps...), $y$ = grandeur à prévoir (coût total, salaire...). Préciser les unités.
\textbf{Étape 2 :} Identifier $m$ (taux, prix unitaire, coefficient directeur) et $p$ (valeur fixe, abonnement, ordonnée à l'origine).
\textbf{Étape 3 :} Écrire l'équation de la droite : $y = mx + p$.
\textbf{Étape 4 :} Exploiter l'équation (calculer $y$ pour un $x$ donné, ou $x$ pour un $y$ donné, ou résoudre un système de deux équations pour comparer deux offres).
\textbf{Étape 5 :} \textbf{Conclure par une phrase rédigée} qui répond exactement à la question posée, avec les unités.
\end{methode}

\begin{exemplebox}[Comparaison de deux compagnies de taxi — Contexte camerounais]
\textbf{Énoncé :} Pour un trajet en ville, la \textbf{Compagnie A} facture une prise en charge de \SI{500}{\text{FCFA}} puis \SI{300}{\text{FCFA}} par kilomètre. La \textbf{Compagnie B} facture \SI{200}{\text{FCFA}} de prise en charge puis \SI{400}{\text{FCFA}} par kilomètre.
\begin{enumerate}
    \item Modéliser le prix $y$ (en FCFA) en fonction de la distance $x$ (en km) pour chaque compagnie.
    \item Pour quelle distance les deux tarifs sont-ils égaux ?
    \item Laquelle choisir pour un trajet de \SI{2}{km} ? De \SI{5}{km} ?
\end{enumerate}

\textbf{Solution :}
\begin{enumerate}
    \item \textbf{Choix des variables :} $x$ = distance en km ; $y_A$ = prix Compagnie A (FCFA) ; $y_B$ = prix Compagnie B (FCFA).
    \item \textbf{Identification des paramètres :}
    \begin{itemize}
        \item Cie A : $m_A = 300$ (FCFA/km), $p_A = 500$ (FCFA). \textbf{Équation :} $y_A = 300x + 500$.
        \item Cie B : $m_B = 400$ (FCFA/km), $p_B = 200$ (FCFA). \textbf{Équation :} $y_B = 400x + 200$.
    \end{itemize}
    \item \textbf{Égalité des tarifs :} On résout $y_A = y_B$.
    \[
    300x + 500 = 400x + 200 \iff 100x = 300 \iff x = 3.
    \]
    Les tarifs sont égaux pour \textbf{\SI{3}{km}} (prix commun : $300 \times 3 + 500 = \SI{1400}{\text{FCFA}}$).
    \item \textbf{Comparaison :}
    \begin{itemize}
        \item Pour $x = 2$ km : $y_A = 1100$, $y_B = 1000$. \textbf{La Compagnie B est moins chère.}
        \item Pour $x = 5$ km : $y_A = 2000$, $y_B = 2200$. \textbf{La Compagnie A est moins chère.}
    \end{itemize}
    \item \textbf{Conclusion :} Pour les courts trajets ($< 3$ km), la Compagnie B est avantageuse grâce à sa faible prise en charge. Au-delà de \SI{3}{km}, la Compagnie A devient plus intéressante car son prix au km est plus bas.
\end{enumerate}
\end{exemplebox}

\begin{popfigure}
\begin{tikzpicture}
\begin{axis}[
    width=\linewidth, height=0.6\linewidth,
    axis lines=middle,
    xlabel={Distance $x$ (km)}, ylabel={Prix $y$ (FCFA)},
    xmin=0, xmax=6, ymin=0, ymax=2500,
    xtick={0,1,2,3,4,5,6}, ytick={0,500,1000,1400,1500,2000,2500},
    xticklabel style={font=\small}, yticklabel style={font=\small},
    grid=major, grid style={popline},
    legend style={at={(0.02,0.98)}, anchor=north west, font=\small, fill=white, draw=popmuted},
    clip=false
]
% Compagnie A : y = 300x + 500
\addplot[popblue, thick, domain=0:6, samples=2, popcurve] {300*x + 500};
\addlegendentry{Compagnie A : $y_A = 300x + 500$}
% Compagnie B : y = 400x + 200
\addplot[poporange, thick, domain=0:6, samples=2, popcurve] {400*x + 200};
\addlegendentry{Compagnie B : $y_B = 400x + 200$}
% Point d'intersection (3, 1400)
\addplot[only marks, mark=*, mark size=2.5pt, popgreen] coordinates {(3,1400)};
\node[above right, font=\small, popgreen] at (axis cs:3,1400) {$I(3\,;\,1400)$};
% Lignes pointillées vers les axes
\draw[popgreen, dashed, thin] (axis cs:3,0) -- (axis cs:3,1400);
\draw[popgreen, dashed, thin] (axis cs:0,1400) -- (axis cs:3,1400);
\node[below, font=\tiny, popgreen] at (axis cs:3,0) {3 km};
\node[left, font=\tiny, popgreen] at (axis cs:0,1400) {1400 FCFA};
% Zones
\node[align=center, font=\footnotesize, popblue!80!black] at (axis cs:5,2300) {Au-delà de 3 km :\\A moins chère};
\node[align=center, font=\footnotesize, poporange!80!black] at (axis cs:1.5,800) {Avant 3 km :\\B moins chère};
\end{axis}
\end{tikzpicture}
\legende{Comparaison graphique des tarifs des deux compagnies. Le point d'intersection $I(3\,;\,1400)$ correspond au kilométrage d'égalité des prix (seuil de rentabilité).}
\end{popfigure}

\courssub{Exploiter le modèle : prévoir, comparer, déterminer un seuil}

L'exploitation du modèle repose sur la résolution d'équations du premier degré (déjà maîtrisées) ou de systèmes de deux équations (intersection de deux droites).

\begin{exoresolu}[Forfaits Internet — Modélisation et décision]
\textbf{Énoncé :} Un élève compare deux forfaits Internet pour son modem 4G.
\begin{itemize}
    \item \textbf{Forfait « Étudiant » (MTN) :} \SI{2000}{\text{FCFA}}/mois pour \SI{5}{Go}, puis \SI{500}{\text{FCFA}}/Go supplémentaire.
    \item \textbf{Forfait « Illimité Soir » (Orange) :} \SI{5000}{\text{FCFA}}/mois pour \SI{20}{Go} (utilisables de 18h à 6h), au-delà \SI{300}{\text{FCFA}}/Go.
\end{itemize}
On suppose que l'élève consomme $x$ Go supplémentaires au-delà du volume de base inclus (donc $x \ge 0$).
\begin{enumerate}
    \item Exprimer le coût mensuel $C_M(x)$ (MTN) et $C_O(x)$ (Orange) en fonction de $x$.
    \item Déterminer le volume supplémentaire $x$ pour lequel les deux forfaits coûtent le même prix.
    \item Conseiller l'élève s'il prévoit \SI{3}{Go} supplémentaires, puis \SI{15}{Go} supplémentaires.
\end{enumerate}

\textbf{Solution détaillée :}
\begin{enumerate}
    \item \textbf{Modélisation :}
    \begin{align*}
        C_M(x) &= 500x + 2000 \quad \text{(prix de base 2000 + 500 par Go suppl.)} \\
        C_O(x) &= 300x + 5000 \quad \text{(prix de base 5000 + 300 par Go suppl.)}
    \end{align*}
    \item \textbf{Seuil d'égalité (intersection) :} On résout $C_M(x) = C_O(x)$.
    \[
    500x + 2000 = 300x + 5000 \iff 200x = 3000 \iff x = 15.
    \]
    Pour \textbf{\SI{15}{Go} supplémentaires} (soit 20 Go total pour MTN, 35 Go total pour Orange), les coûts sont égaux : $500 \times 15 + 2000 = \SI{9500}{\text{FCFA}}$.
    \item \textbf{Analyse et conseil :}
    \begin{itemize}
        \item Pour $x = 3$ Go suppl. : $C_M(3) = 3500$ FCFA ; $C_O(3) = 5900$ FCFA. \textbf{Le forfait MTN est bien moins cher.}
        \item Pour $x = 15$ Go suppl. : Égalité (\SI{9500}{\text{FCFA}}).
        \item Pour $x > 15$ (ex. $x=20$) : $C_M(20)=12000$, $C_O(20)=11000$. \textbf{Le forfait Orange devient moins cher.}
    \end{itemize}
    \item \textbf{Conclusion rédigée :} Si l'élève consomme peu de data supplémentaire ($< 15$ Go), le forfait « Étudiant » MTN est plus économique. S'il est gros consommateur ($> 15$ Go supplémentaires), le forfait « Illimité Soir » Orange devient avantageux grâce à son coût marginal plus faible (\SI{300}{\text{FCFA}}/Go contre \SI{500}{\text{FCFA}}/Go).
\end{enumerate}
\end{exoresolu}

\begin{attention}[Exigence de rédaction — Piège fréquent au BEPC]
\textbf{Ne jamais terminer un exercice par un simple nombre ou une inégalité.} Le savoir-faire « Rédiger et justifier » impose une \textbf{phrase de conclusion complète}, en français, qui répond à la question concrète.
\begin{itemize}
    \item \textbf{Mauvais :} $x = 3$ km. / $x > 15$.
    \item \textbf{Bon :} « Les deux tarifs sont égaux pour une distance de \SI{3}{km}. Pour un trajet de \SI{5}{km}, il faut choisir la Compagnie A qui coûte \SI{2000}{\text{FCFA}} contre \SI{2200}{\text{FCFA}} pour la Compagnie B. »
\end{itemize}
Vérifie toujours que ta conclusion contient : la réponse chiffrée \textbf{ET} l'unité \textbf{ET} la réponse à la question (choix, prévision, interprétation).
\end{attention}

\courssub{Ouverture (Terminale) : Droites perpendiculaires et distance d'un point à une droite}

\emph{Les résultats ci-dessous ne sont \textbf{pas exigibles} en classe de 3ème ni au BEPC. Ils sont donnés en ouverture pour préparer la classe de Seconde/Terminale où ils seront démontrés et utilisés systématiquement.}

\begin{savaistu}[Produit des coefficients directeurs = $-1$ — Ouverture Terminale]
Dans un repère orthonormé, deux droites \textbf{non verticales} d'équations $y = mx + p$ et $y = m'x + p'$ sont \textbf{perpendiculaires} (orthogonales) \textbf{si et seulement si} :
\[
\boxed{m \times m' = -1}
\]
\textbf{Exemple :} La droite $d_1 : y = 2x + 1$ est perpendiculaire à la droite $d_2 : y = -\frac{1}{2}x + 3$ car $2 \times \left(-\frac{1}{2}\right) = -1$.
\textbf{Cas particuliers :} Une droite horizontale ($m=0$) est perpendiculaire à une droite verticale (équation $x = a$, pas de coefficient directeur). La règle $m \times m' = -1$ ne s'applique pas directement car la verticale n'a pas de coefficient directeur.
\end{savaistu}

\begin{savaistu}[Distance d'un point à une droite — Ouverture Terminale]
La distance d'un point $A(x_A\,;\,y_A)$ à une droite $d$ d'équation cartésienne $ax + by + c = 0$ (avec $a^2+b^2 \ne 0$) est donnée par la formule :
\[
\boxed{d(A,d) = \frac{|ax_A + by_A + c|}{\sqrt{a^2+b^2}}}
\]
Cette formule permet de calculer la plus courte distance (le segment perpendiculaire) sans tracer la figure. Elle sera essentielle en géométrie analytique au lycée.
\end{savaistu}

\courssub{Synthèse de la leçon : Les quatre façons de déterminer une équation de droite}

Toute la leçon « Équations de droites » se résume à ce tableau de correspondance entre la \textbf{donnée géométrique de départ} et la \textbf{méthode algébrique} pour obtenir l'équation cartésienne.

\begin{popfigure}
\begin{tikzpicture}[
    node distance=1.2cm and 1.5cm,
    start/.style={rectangle, rounded corners, draw=popblue, thick, fill=popblueL, text width=3.2cm, align=center, font=\small\bfseries},
    meth/.style={rectangle, rounded corners, draw=poporange, thick, fill=poporangeL, text width=3.5cm, align=center, font=\small},
    eq/.style={rectangle, rounded corners, draw=popgreen, thick, fill=popgreenL, text width=3.5cm, align=center, font=\small},
    arrow/.style={-Stealth, thick, popdark}
]
% Ligne 1 : Deux points
\node[start, align=center] (s1){Donnée : \textbf{Deux points} $A(x_A,y_A)$ et $B(x_B,y_B)$ \\ ($x_A \ne x_B$)};
\node[meth, right=of s1, align=center] (m1){Calculer le coefficient directeur : \\ $m = \dfrac{y_B - y_A}{x_B - x_A}$ \\ Puis utiliser point $A$ : $y - y_A = m(x - x_A)$};
\node[eq, right=of m1, align=center] (e1){Forme obtenue : \\ $y = mx + p$ \\ (équation réduite)};

% Ligne 2 : Point + Vecteur directeur
\node[start, below=of s1, align=center] (s2){Donnée : \textbf{Un point} $A(x_A,y_A)$ \\ et un \textbf{vecteur directeur} $\vec{u}(u_1,u_2)$ \\ ($u_1 \ne 0$)};
\node[meth, right=of s2, align=center] (m2){Coefficient directeur : $m = \dfrac{u_2}{u_1}$ \\ Puis équation point-pente : $y - y_A = m(x - x_A)$};
\node[eq, right=of m2, align=center] (e2){Forme obtenue : \\ $y = mx + p$ \\ ou forme cartésienne : $u_2 x - u_1 y + c = 0$};

% Ligne 3 : Point + Coefficient directeur
\node[start, below=of s2, align=center] (s3){Donnée : \textbf{Un point} $A(x_A,y_A)$ \\ et le \textbf{coefficient directeur} $m$};
\node[meth, right=of s3, align=center] (m3){Équation point-pente directe : \\ $y - y_A = m(x - x_A)$};
\node[eq, right=of m3, align=center] (e3){Forme obtenue : \\ $y = mx + p$ \\ (équation réduite)};

% Ligne 4 : Parallèle / Perpendiculaire (Ouverture)
\node[start, below=of s3, align=center] (s4){Donnée (Ouverture) : \\ Droite $\Delta$ connue et \\ \textbf{parallèle} ou \textbf{perpendiculaire} à $\Delta$ passant par un point};
\node[meth, right=of s4, align=center] (m4){Parallèle : même $m$. \\ Perpendiculaire : $m' = -\frac{1}{m}$ (si $m \ne 0$). \\ Puis point-pente.};
\node[eq, right=of m4, align=center] (e4){Forme obtenue : \\ $y = mx + p$ ou $y = m'x + p'$};

% Flèches
\foreach \i in {1,2,3,4} {
    \draw[arrow] (s\i) -- (m\i);
    \draw[arrow] (m\i) -- (e\i);
}

% Critère position relative (bas)
\node[rectangle, draw=poppurple, thick, fill=poppurpleL, rounded corners, text width=15cm, align=center, font=\small, below=1.5cm of s4] (crit) {
    \textbf{Positions relatives de deux droites $d_1(y=m_1x+p_1)$ et $d_2(y=m_2x+p_2)$ (non verticales) :} \\
    \textbf{Parallèles distinctes} $\iff m_1 = m_2$ et $p_1 \ne p_2$ \quad | \quad
    \textbf{Confondues} $\iff m_1 = m_2$ et $p_1 = p_2$ \quad | \quad
    \textbf{Sécantes} $\iff m_1 \ne m_2$ (intersection : résoudre $m_1x+p_1 = m_2x+p_2$)
};
\end{tikzpicture}
\legende{Tableau de synthèse : de la donnée géométrique à l'équation cartésienne, et critère de position relative de deux droites.}
\end{popfigure}

\begin{aretenir}[Bilan — Ce qu'il faut maîtriser pour le BEPC]
\begin{enumerate}
    \item \textbf{Modéliser :} Lire un énoncé $\rightarrow$ identifier $x$ (variable), $y$ (réponse), $m$ (taux/unitaire), $p$ (fixe/abonné) $\rightarrow$ écrire $y=mx+p$.
    \item \textbf{Exploiter :} Calculer une image ($y$ pour $x$ donné), un antécédent ($x$ pour $y$ donné), comparer deux offres (résoudre $f(x)=g(x)$ $\rightarrow$ abscisse du point d'intersection = seuil de rentabilité).
    \item \textbf{Rédiger :} Phrase de conclusion \textbf{obligatoire}, réponse à la question, unités respectées.
    \item \textbf{Quatre méthodes pour l'équation :} (1) Deux points, (2) Point + vecteur directeur, (3) Point + coefficient directeur, (4) Parallèle/perpendiculaire à une droite connue (ouverture).
    \item \textbf{Positions relatives :} Parallèles ($m$ égaux, $p$ différents), Confondues ($m$ et $p$ égaux), Sécantes ($m$ différents).
\end{enumerate}
\end{aretenir}

\begin{exercice}[Entraînement — Modélisation complète]
\textbf{Énoncé :} Pour l'irrigation d'un champ, un agriculteur hésite entre deux systèmes de pompage.
\begin{itemize}
    \item \textbf{Système S1 (thermique) :} Achat de la pompe \SI{150000}{\text{FCFA}}, consommation carburant \SI{2000}{\text{FCFA}} par heure de fonctionnement.
    \item \textbf{Système S2 (solaire) :} Achat de l'installation \SI{400000}{\text{FCFA}}, coût de fonctionnement quasi nul (\SI{100}{\text{FCFA}}/h pour maintenance).
\end{itemize}
On note $x$ le nombre d'heures de fonctionnement total et $C_1(x)$, $C_2(x)$ les coûts globaux en FCFA.
\begin{enumerate}
    \item Établir les expressions de $C_1(x)$ et $C_2(x)$.
    \item Déterminer le nombre d'heures $x$ pour lequel les deux systèmes reviennent au même coût global (seuil de rentabilité).
    \item L'agriculteur prévoit \SI{150}{h} de pompage par an sur 5 ans. Quel système lui conseillez-vous ? Justifiez par un calcul et une phrase de conclusion.
\end{enumerate}
\end{exercice}

\begin{corrige}[Exercice — Modélisation complète]
\begin{enumerate}
    \item $C_1(x) = 2000x + 150000$ \quad ; \quad $C_2(x) = 100x + 400000$.
    \item Seuil : $2000x + 150000 = 100x + 400000 \iff 1900x = 250000 \iff x \approx 131,58$ h.
    \item Prévision : $150 \times 5 = 750$ h au total. $750 > 131,58$, donc au-delà du seuil.
    $C_1(750) = 2000 \times 750 + 150000 = 1\,650\,000$ FCFA.
    $C_2(750) = 100 \times 750 + 400000 = 475\,000$ FCFA.
    \textbf{Conclusion :} Sur la durée prévue (750 h), le système solaire S2 est beaucoup plus économique (coût global \SI{475000}{\text{FCFA}} contre \SI{1650000}{\text{FCFA}} pour le thermique). L'investissement initial plus élevé est rentabilisé bien avant la fin de la période.
\end{enumerate}
\end{corrige}

\newpage

\coursec{Activité d'intégration}

\coursec{Équations de droites : Activité d'intégration}

\courssub{Objectifs de l'activité}
Cette activité d'intégration vise à mobiliser l'ensemble des savoir-faire travaillés dans la leçon « Équations de droites » pour résoudre un problème authentique de choix économique. Elle permet à l'élève de :
\begin{itemize}
    \item Modéliser une situation concrète par deux fonctions affines (équations réduites de droites).
    \item Représenter graphiquement ces droites dans un repère orthonormé adapté.
    \item Déterminer algébriquement le point d'intersection de deux droites et l'interpréter dans le contexte.
    \item Vérifier la cohérence du modèle en testant des points sur les droites.
    \item Argumenter un choix optimal en fonction d'un paramètre (la charge à transporter).
    \item Utiliser la condition de parallélisme ($m = m'$) pour écarter une hypothèse.
    \item Rédiger une démarche complète, justifiée et structurée, conforme aux attendus du BEPC.
\end{itemize}

\courssub{Situation}
La coopérative \emph{« Bananes de Penja »}, située dans le Moungo, récolte une production de bananes plantains destinée au grand marché de Douala (marché de Bepanda ou de Nkololoun). Pour acheminer sa marchandise, le président de la coopérative contacte deux sociétés de transport routier bien connues sur l'axe Douala--Ouest :

\begin{itemize}
    \item \textbf{Kamga Trans} : facture \textbf{10\,000 FCFA} de frais fixes (chargement, assurance, papier) plus \textbf{800 FCFA par quintal} transporté.
    \item \textbf{Étoile Logistique} : facture \textbf{4\,000 FCFA} de frais fixes plus \textbf{1\,000 FCFA par quintal} transporté.
\end{itemize}

Le président doit expédier un lot dont le poids (en quintaux) peut varier selon la commande : \textbf{5 qx}, \textbf{10 qx} ou \textbf{30 qx}. Il souhaite choisir le transporteur le moins cher pour chaque cas, mais aussi comprendre à partir de quelle charge l'un devient plus avantageux que l'autre.

Un troisième transporteur, \textbf{Bamileke Express}, propose un tarif dont la courbe représentative est une droite \textbf{parallèle} à celle de Kamga Trans, mais passant par le point d'abscisse 0 et d'ordonnée 12\,000 (frais fixes plus élevés, même coût au quintal).

\courssub{Consignes}
\begin{enumerate}
    \item \textbf{Modélisation.} Soit $x$ la masse en quintaux ($x \ge 0$) et $y$ le coût total en FCFA.
    \begin{enumerate}
        \item Écrire l'équation réduite de la droite $(d_K)$ modélisant le tarif de Kamga Trans.
        \item Écrire l'équation réduite de la droite $(d_E)$ modélisant le tarif d'Étoile Logistique.
        \item Écrire l'équation réduite de la droite $(d_B)$ modélisant le tarif de Bamileke Express.
    \end{enumerate}
    \item \textbf{Représentation graphique.} Dans un repère orthonormé $(O; \vec{\imath}, \vec{\jmath})$ (1 cm pour 2 qx en abscisses, 1 cm pour 2\,000 FCFA en ordonnées), tracer les trois droites $(d_K)$, $(d_E)$ et $(d_B)$ pour $x \in [0 ; 35]$. On placera les points caractéristiques (ordonnées à l'origine, intersection).
    \item \textbf{Intersection et interprétation.}
    \begin{enumerate}
        \item Déterminer par le calcul les coordonnées du point d'intersection $I$ des droites $(d_K)$ et $(d_E)$.
        \item Interpréter concrètement l'abscisse et l'ordonnée de $I$ pour le président de la coopérative.
    \end{enumerate}
    \item \textbf{Vérification du modèle.} Pour chacune des droites $(d_K)$ et $(d_E)$, choisir deux valeurs de $x$ distinctes (différentes de l'abscisse de $I$), calculer les $y$ correspondants et vérifier que les points obtenus appartiennent bien à la droite (vérification graphique et algébrique).
    \item \textbf{Recommandation argumentée.} Pour chacune des charges prévues (5 qx, 10 qx, 30 qx), indiquer le transporteur le moins cher et justifier votre réponse par le calcul du coût ou par la position relative des droites.
    \item \textbf{Prolongement : le cas Bamileke Express.} Justifier rigoureusement, en utilisant la condition de parallélisme, pourquoi Bamileke Express ne pourra \textbf{jamais} être le moins cher des trois, quelle que soit la charge $x > 0$.
\end{enumerate}

\courssub{Ressources à mobiliser}
\begin{itemize}
    \item \textbf{Équation réduite d'une droite :} Une droite non verticale d'équation $y = mx + p$ a pour coefficient directeur $m$ et pour ordonnée à l'origine $p$.
    \item \textbf{Droite passant par deux points :} Si la droite passe par $A(x_A; y_A)$ et $B(x_B; y_B)$ ($x_A \ne x_B$), son coefficient directeur est $m = \dfrac{y_B - y_A}{x_B - x_A}$.
    \item \textbf{Droite de coefficient directeur donné passant par un point :} $y - y_0 = m(x - x_0)$.
    \item \textbf{Positions relatives :} Deux droites $y = mx + p$ et $y = m'x + p'$ sont \textbf{parallèles} si et seulement si $m = m'$ (et $p \ne p'$). Elles sont \textbf{confondues} si $m = m'$ et $p = p'$. Elles sont \textbf{sécantes} si $m \ne m'$ ; leur intersection s'obtient en résolvant $mx + p = m'x + p'$.
    \item \textbf{Appartenance :} Un point $M(x_0; y_0)$ appartient à la droite $y = mx + p$ si et seulement si $y_0 = mx_0 + p$.
\end{itemize}

\courssub{Démarche et corrigé commenté}
\begin{exoresolu}[Corrigé détaillé : Modélisation, Graphique, Intersection]
\paragraph{1. Modélisation.}
On note $x$ la masse en quintaux ($x \ge 0$) et $y$ le coût total en FCFA.
\begin{itemize}
    \item \textbf{Kamga Trans :} Frais fixes = 10\,000 FCFA $\rightarrow$ ordonnée à l'origine $p_K = 10\,000$. Coût au quintal = 800 FCFA $\rightarrow$ coefficient directeur $m_K = 800$.
    \begin{center}
        \boxed{(d_K) : y = 800x + 10\,000}
    \end{center}
    \item \textbf{Étoile Logistique :} Frais fixes = 4\,000 FCFA $\rightarrow$ $p_E = 4\,000$. Coût au quintal = 1\,000 FCFA $\rightarrow$ $m_E = 1\,000$.
    \begin{center}
        \boxed{(d_E) : y = 1\,000x + 4\,000}
    \end{center}
    \item \textbf{Bamileke Express :} Droite parallèle à $(d_K)$ $\rightarrow$ même coefficient directeur $m_B = m_K = 800$. Passe par $(0; 12\,000)$ $\rightarrow$ ordonnée à l'origine $p_B = 12\,000$.
    \begin{center}
        \boxed{(d_B) : y = 800x + 12\,000}
    \end{center}
\end{itemize}

\paragraph{2. Représentation graphique.}
\begin{popfigure}
\begin{tikzpicture}[scale=0.7]
    % Axes : x max 35 -> 17.5cm, y max 36000 -> 18cm. On prend un peu de marge.
    \draw[->, thick] (0,0) -- (19,0) node[right] {$x$ (qx)};
    \draw[->, thick] (0,0) -- (0,20) node[above] {$y$ (kFCFA)};
    % Graduations abscisses : 1 cm = 2 qx -> coord = x/2
    \foreach \x in {0,5,10,15,20,25,30,35} {
        \draw (\x/2, -0.3) -- (\x/2, 0.3) node[below] {\small $\x$};
    }
    % Graduations ordonnées : 1 cm = 2000 FCFA -> coord = y/2000. Label en kFCFA (y/1000)
    \foreach \k in {0,2,4,6,8,10,12,14,16,18,20,22,24,26,28,30,32,34,36} {
        \draw (-0.3, \k/2) -- (0.3, \k/2) node[left] {\small $\k$};
    }
    % Droite d_K : y = 800x + 10000. En coord graphique : Y = 0.8*X + 5 (X=x/2, Y=y/2000)
    % Domaine x in [0,35] -> X in [0, 17.5]
    \draw[popblue, thick, domain=0:17.5, samples=100] plot (\x, {0.8*\x + 5}) node[right] {$(d_K)$};
    % Droite d_E : y = 1000x + 4000 -> Y = 1.0*X + 2
    \draw[poporange, thick, domain=0:17.5, samples=100] plot (\x, {1.0*\x + 2}) node[right] {$(d_E)$};
    % Droite d_B : y = 800x + 12000 -> Y = 0.8*X + 6
    \draw[popgreen, thick, domain=0:17.5, samples=100] plot (\x, {0.8*\x + 6}) node[right] {$(d_B)$};
    % Point I intersection : x=30 -> X=15. y=34000 -> Y=17.
    \fill[popdark] (15, 17) circle (3pt) node[above left] {$I(30; 34\,000)$};
    % Points ordonnées origine
    \fill[popblue] (0,5) circle (2pt) node[left] {$K_0(0; 10\,000)$};
    \fill[poporange] (0,2) circle (2pt) node[left] {$E_0(0; 4\,000)$};
    \fill[popgreen] (0,6) circle (2pt) node[left] {$B_0(0; 12\,000)$};
\end{tikzpicture}
\legende{Représentation graphique des trois tarifs. Échelle : 1 cm = 2 qx (abscisses), 1 cm = 2 000 FCFA (ordonnées).}
\end{popfigure}

\paragraph{3. Intersection et interprétation.}
Les droites $(d_K)$ et $(d_E)$ sont sécantes car $m_K = 800 \ne 1000 = m_E$.
On résout $800x + 10\,000 = 1\,000x + 4\,000$ :
\begin{align*}
    10\,000 - 4\,000 &= 1\,000x - 800x \\
    6\,000 &= 200x \\
    x &= \frac{6\,000}{200} = 30
\end{align*}
En remplaçant dans $(d_K)$ : $y = 800 \times 30 + 10\,000 = 24\,000 + 10\,000 = 34\,000$.
Le point d'intersection est $I(30; 34\,000)$.
\end{exoresolu}

\begin{aretenir}[Interprétation du point d'intersection]
Pour une charge de \textbf{30 quintaux}, les deux transporteurs facturent exactement le même montant : \textbf{34\,000 FCFA}.
\begin{itemize}
    \item Si $x < 30$ (charge légère), la droite $(d_E)$ est \textbf{sous} $(d_K)$ : Étoile Logistique est moins cher.
    \item Si $x > 30$ (charge lourde), la droite $(d_K)$ est \textbf{sous} $(d_E)$ : Kamga Trans est moins cher.
\end{itemize}
L'abscisse 30 est le \textbf{seuil de rentabilité} (ou point d'équilibre) à partir duquel Kamga Trans devient plus avantageux.
\end{aretenir}

\begin{exoresolu}[Corrigé détaillé : Vérification et Recommandation]
\paragraph{4. Vérification du modèle.}
\begin{itemize}
    \item \textbf{Pour $(d_K) : y = 800x + 10\,000$}
    \begin{itemize}
        \item $x = 5$ : $y = 800 \times 5 + 10\,000 = 14\,000$. Point $A_K(5; 14\,000)$. Vérification : $14\,000 = 800 \times 5 + 10\,000 = 14\,000$ \checkmark
        \item $x = 20$ : $y = 800 \times 20 + 10\,000 = 26\,000$. Point $B_K(20; 26\,000)$. Vérification : $26\,000 = 16\,000 + 10\,000 = 26\,000$ \checkmark
    \end{itemize}
    \item \textbf{Pour $(d_E) : y = 1\,000x + 4\,000$}
    \begin{itemize}
        \item $x = 5$ : $y = 1\,000 \times 5 + 4\,000 = 9\,000$. Point $A_E(5; 9\,000)$. Vérification : $9\,000 = 5\,000 + 4\,000 = 9\,000$ \checkmark
        \item $x = 20$ : $y = 1\,000 \times 20 + 4\,000 = 24\,000$. Point $B_E(20; 24\,000)$. Vérification : $24\,000 = 20\,000 + 4\,000 = 24\,000$ \checkmark
    \end{itemize}
\end{itemize}
Sur le graphique, ces points sont bien alignés sur leurs droites respectives.

\paragraph{5. Recommandation argumentée.}
On compare les coûts pour les trois charges demandées.

\begin{center}
\begin{tabular}{|c|c|c|c|c|}\hline
Charge ($x$) & Coût Kamga Trans $(d_K)$ & Coût Étoile Logistique $(d_E)$ & Coût Bamileke $(d_B)$ & \textbf{Choix optimal} \\ \hline
5 qx & $800\times5+10\,000 = 14\,000$ & $1\,000\times5+4\,000 = \mathbf{9\,000}$ & $800\times5+12\,000 = 16\,000$ & \textbf{Étoile Logistique} \\ \hline
10 qx & $800\times10+10\,000 = 18\,000$ & $1\,000\times10+4\,000 = \mathbf{14\,000}$ & $800\times10+12\,000 = 20\,000$ & \textbf{Étoile Logistique} \\ \hline
30 qx & $800\times30+10\,000 = \mathbf{34\,000}$ & $1\,000\times30+4\,000 = \mathbf{34\,000}$ & $800\times30+12\,000 = 36\,000$ & \textbf{Indifférent (Kamga ou Étoile)} \\ \hline
\end{tabular}
\end{center}
\end{exoresolu}

\begin{attention}[Rédaction attendue au BEPC]
Pour 5 qx et 10 qx : « Comme $x < 30$, la droite $(d_E)$ est située sous $(d_K)$ et sous $(d_B)$, donc Étoile Logistique est le moins cher. »
Pour 30 qx : « Les deux droites $(d_K)$ et $(d_E)$ se coupent en $I(30; 34\,000)$, les coûts sont égaux. Bamileke reste plus cher. »
\end{attention}

\begin{exoresolu}[Corrigé détaillé : Prolongement Bamileke Express]
\paragraph{6. Prolongement : Bamileke Express.}
La droite $(d_B)$ a pour équation $y = 800x + 12\,000$. La droite $(d_K)$ a pour équation $y = 800x + 10\,000$.
Elles ont le \textbf{même coefficient directeur} ($m_B = m_K = 800$) et des \textbf{ordonnées à l'origine différentes} ($12\,000 \ne 10\,000$).
\cle{Propriété} : Deux droites de même coefficient directeur et d'ordonnées à l'origine distinctes sont \textbf{parallèles strictement} (elles n'ont aucun point commun).

De plus, pour tout $x \ge 0$ :
$$y_B - y_K = (800x + 12\,000) - (800x + 10\,000) = 2\,000 > 0.$$
La courbe de Bamileke Express est \textbf{toujours au-dessus} de celle de Kamga Trans de 2\,000 FCFA constants.
Puisque Kamga Trans est déjà moins cher qu'Étoile Logistique pour $x > 30$, et qu'Étoile Logistique est moins cher que Kamga Trans pour $x < 30$, Bamileke Express, qui est \textbf{toujours plus cher que Kamga Trans}, ne peut \textbf{jamais} être le moins cher des trois, quelle que soit la charge $x > 0$.
\end{exoresolu}

\courssub{Bilan de l'activité}
Cette activité a permis de mettre en évidence la puissance de la modélisation par des fonctions affines (droites) pour comparer des offres commerciales réelles. Les compétences clés validées sont :
\begin{enumerate}
    \item \textbf{Traduire un énoncé concret en langage mathématique} (identifier $p$ comme coût fixe et $m$ comme coût variable unitaire).
    \item \textbf{Manipuler l'équation réduite $y = mx + p$} dans les trois cas de figure : deux points (données chiffrées), coefficient directeur et point (parallélisme), coefficient directeur et ordonnée à l'origine.
    \item \textbf{Résoudre graphiquement et algébriquement} un système de deux équations du premier degré (recherche d'intersection).
    \item \textbf{Interpréter le résultat mathématique} (coordonnées du point $I$) en termes de décision économique (seuil de rentabilité / point d'équilibre).
    \item \textbf{Justifier rigoureusement} l'ordre des coûts par la position relative des droites (sécantes, parallèles) et le signe de la différence des fonctions.
    \item \textbf{Rédiger une réponse complète} : calculs, phrases explicatives, conclusion contextuelle.
\end{enumerate}
L'erreur fréquente à éviter est de confondre « coefficient directeur » et « ordonnée à l'origine » lors de la lecture du tarif, ou d'oublier de vérifier la condition $p \ne p'$ pour affirmer le parallélisme strict. L'exemple de Bamileke Express illustre parfaitement qu'un transporteur aligné sur le coût variable d'un concurrent mais avec des frais fixes plus élevés est \emph{stricto sensu} dominé : il ne sera jamais optimal.

\newpage

\coursec{Méthodes et exercices résolus}

\courssub{Méthode 1 : Déterminer une équation cartésienne d'une droite passant par deux points}

Dans tout ce chapitre, le plan est muni d'un repère $(O\,;\vec{i},\vec{j})$. On sait depuis la classe de seconde qu'une droite est entièrement déterminée par deux points distincts, ou par un point et un vecteur directeur. L'objectif est de traduire ces données géométriques par une \cle{équation} : une relation que doivent vérifier les coordonnées $(x\,;y)$ de tout point $M$ de la droite. L'outil central est la \textbf{condition de colinéarité de deux vecteurs} : $\vec{u}\begin{pmatrix} a \\ b \end{pmatrix}$ et $\vec{v}\begin{pmatrix} a' \\ b' \end{pmatrix}$ sont colinéaires si et seulement si $ab' - a'b = 0$.

\begin{methode}[Équation d'une droite passant par deux points]
Pour déterminer une équation cartésienne de la droite $(AB)$ avec $A(x_A\,;y_A)$ et $B(x_B\,;y_B)$ :
\begin{enumerate}
\item \textbf{Vérifier} que $A \neq B$ (sinon il n'y a pas une unique droite passant par $A$ et $B$).
\item \textbf{Calculer} les coordonnées du vecteur directeur $\overrightarrow{AB}\begin{pmatrix} x_B - x_A \\ y_B - y_A \end{pmatrix}$.
\item \textbf{Écrire} qu'un point $M(x\,;y)$ appartient à $(AB)$ si et seulement si $\overrightarrow{AM}$ et $\overrightarrow{AB}$ sont colinéaires, c'est-à-dire :
\[ (x - x_A)(y_B - y_A) - (y - y_A)(x_B - x_A) = 0 \]
\item \textbf{Développer et réduire} pour obtenir une équation de la forme $ax + by + c = 0$.
\item \textbf{Contrôler} en vérifiant que les coordonnées de $A$ et de $B$ satisfont l'équation trouvée.
\end{enumerate}
\textbf{Formule à retenir :} si $\vec{u}\begin{pmatrix} \alpha \\ \beta \end{pmatrix}$ est un vecteur directeur de $(D)$, alors $(D)$ admet une équation de la forme $-\beta x + \alpha y + c = 0$ : les coefficients de $x$ et de $y$ se lisent directement sur le vecteur directeur (avec échange des coordonnées et changement d'un signe).
\end{methode}

\begin{exoresolu}[Droite passant par deux points]
Le plan est muni d'un repère $(O\,;\vec{i},\vec{j})$. On donne $A(2\,;-1)$ et $B(5\,;3)$. Déterminer une équation cartésienne de la droite $(AB)$.
\tcblower
\textbf{Étape 1 :} $A \neq B$, donc la droite $(AB)$ existe et est unique.

\textbf{Étape 2 :} Calculons $\overrightarrow{AB}$ :
\[ \overrightarrow{AB}\begin{pmatrix} 5 - 2 \\ 3 - (-1) \end{pmatrix} = \begin{pmatrix} 3 \\ 4 \end{pmatrix}. \]

\textbf{Étape 3 :} Soit $M(x\,;y)$ un point du plan. Alors $\overrightarrow{AM}\begin{pmatrix} x - 2 \\ y - (-1) \end{pmatrix} = \begin{pmatrix} x - 2 \\ y + 1 \end{pmatrix}$. On a :
\[ M \in (AB) \iff \overrightarrow{AM} \text{ et } \overrightarrow{AB} \text{ colinéaires} \iff 4(x-2) - 3(y+1) = 0. \]

\textbf{Étape 4 :} Développons :
\begin{align*}
4(x-2) - 3(y+1) &= 0 \\
4x - 8 - 3y - 3 &= 0 \\
4x - 3y - 11 &= 0.
\end{align*}

\textbf{Étape 5 :} Vérification : pour $A(2\,;-1)$ : $4(2) - 3(-1) - 11 = 8 + 3 - 11 = 0$ \checkmark ; pour $B(5\,;3)$ : $4(5) - 3(3) - 11 = 20 - 9 - 11 = 0$ \checkmark.

\textbf{Conclusion :} une équation cartésienne de $(AB)$ est $\boxed{4x - 3y - 11 = 0}$.
\end{exoresolu}

\courssub{Méthode 2 : Déterminer une équation d'une droite de vecteur directeur donné}

\begin{methode}[Droite définie par un point et un vecteur directeur]
Pour trouver une équation de la droite $(D)$ passant par $A(x_A\,;y_A)$ et de vecteur directeur $\vec{u}\begin{pmatrix} \alpha \\ \beta \end{pmatrix}$ (avec $\vec{u} \neq \vec{0}$) :
\begin{enumerate}
\item \textbf{Écrire} la condition de colinéarité : $M(x\,;y) \in (D) \iff \overrightarrow{AM}$ et $\vec{u}$ sont colinéaires, soit $\beta(x - x_A) - \alpha(y - y_A) = 0$.
\item \textbf{Développer} pour obtenir une équation de la forme $ax + by + c = 0$.
\item \textbf{Réflexe rapide :} l'équation est directement de la forme $-\beta x + \alpha y + c = 0$ ; on calcule $c$ en remplaçant $x$ et $y$ par les coordonnées de $A$.
\item \textbf{Vérifier} que $(a\,;b) \neq (0\,;0)$, ce qui est garanti par l'hypothèse $\vec{u} \neq \vec{0}$.
\end{enumerate}
\textbf{Attention :} un vecteur directeur de $(D) : ax + by + c = 0$ est $\vec{u}\begin{pmatrix} -b \\ a \end{pmatrix}$, et non $\begin{pmatrix} a \\ b \end{pmatrix}$.
\end{methode}

\begin{exoresolu}[Droite de vecteur directeur donné]
Dans un repère $(O\,;\vec{i},\vec{j})$, déterminer une équation cartésienne de la droite $(D)$ passant par $A(-1\,;4)$ et de vecteur directeur $\vec{u}\begin{pmatrix} 2 \\ -5 \end{pmatrix}$.
\tcblower
\textbf{Méthode détaillée :} Soit $M(x\,;y)$. On a $\overrightarrow{AM}\begin{pmatrix} x - (-1) \\ y - 4 \end{pmatrix} = \begin{pmatrix} x + 1 \\ y - 4 \end{pmatrix}$.
\[ M \in (D) \iff (-5)(x+1) - 2(y - 4) = 0 \]
Développons :
\begin{align*}
-5x - 5 - 2y + 8 &= 0 \\
-5x - 2y + 3 &= 0.
\end{align*}
En multipliant les deux membres par $-1$ (ce qui ne change pas l'ensemble des solutions, donc pas la droite) : $5x + 2y - 3 = 0$.

\textbf{Vérification par le réflexe rapide :} avec $\vec{u}\begin{pmatrix} 2 \\ -5 \end{pmatrix}$, l'équation est de la forme $-(-5)x + 2y + c = 0$, c'est-à-dire $5x + 2y + c = 0$. Le point $A(-1\,;4)$ appartient à $(D)$, donc $5(-1) + 2(4) + c = 0$, soit $-5 + 8 + c = 0$, d'où $c = -3$.

\textbf{Conclusion :} une équation de $(D)$ est $\boxed{5x + 2y - 3 = 0}$.
\end{exoresolu}

\courssub{Méthode 3 : Passer de l'équation cartésienne à l'équation réduite}

\begin{methode}[Équation réduite et coefficient directeur]
Soit $(D) : ax + by + c = 0$ une droite.
\begin{enumerate}
\item \textbf{Tester le coefficient $b$ :}
\begin{itemize}
\item Si $b \neq 0$, on isole $y$ : $(D) : y = -\dfrac{a}{b}x - \dfrac{c}{b}$. C'est l'\cle{équation réduite}, de la forme $y = mx + p$ avec $m = -\dfrac{a}{b}$ le \cle{coefficient directeur} et $p$ l'\cle{ordonnée à l'origine}.
\item Si $b = 0$, la droite est \cle{parallèle à l'axe des ordonnées} : $(D) : x = -\dfrac{c}{a}$ ; elle n'a pas de coefficient directeur.
\end{itemize}
\item \textbf{Lire} $m$ et $p$ directement dans l'écriture $y = mx + p$.
\item \textbf{Interpréter} : $p$ est l'ordonnée du point d'intersection de $(D)$ avec l'axe des ordonnées, c'est-à-dire le point $(0\,;p)$.
\end{enumerate}
\textbf{Formule à retenir :} $(D) : ax + by + c = 0$ avec $b \neq 0 \Rightarrow m = -\dfrac{a}{b}$.
\end{methode}

\begin{exoresolu}[Équation réduite et coefficient directeur]
On donne les droites $(D_1) : 2x - 4y + 8 = 0$ et $(D_2) : 3x - 6 = 0$.
\begin{enumerate}
\item Déterminer l'équation réduite de $(D_1)$, son coefficient directeur et son ordonnée à l'origine.
\item La droite $(D_2)$ admet-elle un coefficient directeur ? Justifier.
\end{enumerate}
\tcblower
\textbf{1.} Pour $(D_1)$, le coefficient de $y$ est $-4 \neq 0$, donc on peut isoler $y$ :
\begin{align*}
2x - 4y + 8 &= 0 \\
-4y &= -2x - 8 \\
y &= \frac{-2x - 8}{-4} = \frac{1}{2}x + 2.
\end{align*}
L'équation réduite de $(D_1)$ est $y = \dfrac{1}{2}x + 2$. Son coefficient directeur est $m = \dfrac{1}{2}$ et son ordonnée à l'origine est $p = 2$ : la droite coupe l'axe des ordonnées au point $(0\,;2)$.

\textbf{2.} Pour $(D_2) : 3x - 6 = 0$, le coefficient de $y$ est nul ($b = 0$). On obtient $x = 2$ : c'est une droite parallèle à l'axe des ordonnées. Par définition, une telle droite \textbf{n'admet pas de coefficient directeur}.

\textbf{Conclusion :} $(D_1) : y = \dfrac{1}{2}x + 2$, $m = \dfrac{1}{2}$, $p = 2$ ; $(D_2) : x = 2$ n'a pas de coefficient directeur.
\end{exoresolu}

\courssub{Méthode 4 : Déterminer une équation d'une droite de coefficient directeur donné}

\begin{methode}[Droite passant par un point, de coefficient directeur donné]
Pour déterminer l'équation réduite de la droite $(D)$ passant par $A(x_A\,;y_A)$ et de coefficient directeur $m$ :
\begin{enumerate}
\item \textbf{Écrire} la forme générale : $(D) : y = mx + p$.
\item \textbf{Calculer $p$} en utilisant le fait que $A \in (D)$ : $y_A = m x_A + p$, donc $p = y_A - m x_A$.
\item \textbf{Conclure} en écrivant l'équation complète, puis vérifier en remplaçant $x$ par $x_A$.
\end{enumerate}
\textbf{Formule à retenir :} $(D) : y = m(x - x_A) + y_A$.
\end{methode}

\begin{exoresolu}[Droite de coefficient directeur donné]
Déterminer l'équation réduite de la droite $(\Delta)$ passant par $A(3\,;-2)$ et de coefficient directeur $m = -\dfrac{2}{3}$.
\tcblower
\textbf{Étape 1 :} $(\Delta)$ admet une équation de la forme $y = -\dfrac{2}{3}x + p$.

\textbf{Étape 2 :} Comme $A(3\,;-2) \in (\Delta)$, ses coordonnées vérifient l'équation :
\[ -2 = -\frac{2}{3} \times 3 + p \iff -2 = -2 + p \iff p = 0. \]

\textbf{Étape 3 :} L'équation réduite de $(\Delta)$ est donc $y = -\dfrac{2}{3}x$.

\textbf{Vérification :} pour $x = 3$ : $y = -\dfrac{2}{3} \times 3 = -2$ \checkmark. De plus $p = 0$ signifie que $(\Delta)$ passe par l'origine $O$ du repère.

\textbf{Conclusion :} $(\Delta) : y = -\dfrac{2}{3}x$.
\end{exoresolu}

\courssub{Méthode 5 : Étudier la position relative de deux droites}

\begin{methode}[Positions relatives de deux droites]
Soient $(D) : y = mx + p$ et $(D') : y = m'x + p'$ deux droites non parallèles à l'axe des ordonnées.
\begin{enumerate}
\item \textbf{Comparer les coefficients directeurs :}
\begin{itemize}
\item Si $m \neq m'$ : $(D)$ et $(D')$ sont \cle{sécantes}. On trouve le point d'intersection en résolvant l'équation $mx + p = m'x + p'$.
\item Si $m = m'$ et $p = p'$ : les droites sont \cle{confondues}.
\item Si $m = m'$ et $p \neq p'$ : les droites sont \cle{strictement parallèles}.
\end{itemize}
\item \textbf{Pour le point d'intersection}, résoudre le système formé par les deux équations, puis vérifier le couple solution dans les deux équations.
\end{enumerate}
\textbf{Réflexe :} si les équations sont cartésiennes $ax + by + c = 0$ et $a'x + b'y + c' = 0$, les droites sont parallèles (ou confondues) si et seulement si $ab' - a'b = 0$ (vecteurs directeurs colinéaires).
\end{methode}

\begin{exoresolu}[Positions relatives et point d'intersection]
Dans un repère, on donne $(D) : y = 2x - 1$, $(D') : y = -x + 5$ et $(D'') : y = 2x + 3$.
\begin{enumerate}
\item Étudier la position relative de $(D)$ et $(D'')$.
\item Étudier la position relative de $(D)$ et $(D')$ ; déterminer, s'il existe, leur point d'intersection.
\end{enumerate}
\tcblower
\textbf{1.} $(D)$ a pour coefficient directeur $m = 2$ et $(D'')$ a pour coefficient directeur $m'' = 2$. On a $m = m''$ mais $p = -1 \neq 3 = p''$.

Donc $(D)$ et $(D'')$ sont \textbf{strictement parallèles}.

\textbf{2.} $(D')$ a pour coefficient directeur $m' = -1 \neq 2 = m$. Donc $(D)$ et $(D')$ sont \textbf{sécantes} en un unique point.

Cherchons ce point $I(x\,;y)$. Ses coordonnées vérifient les deux équations, donc :
\begin{align*}
2x - 1 &= -x + 5 \\
2x + x &= 5 + 1 \\
3x &= 6 \\
x &= 2.
\end{align*}
En reportant dans l'équation de $(D)$ : $y = 2 \times 2 - 1 = 3$.

\textbf{Vérification} dans $(D')$ : $y = -2 + 5 = 3$ \checkmark.

\textbf{Conclusion :} $(D)$ et $(D'')$ sont strictement parallèles ; $(D)$ et $(D')$ sont sécantes au point $I(2\,;3)$.
\end{exoresolu}

\begin{popfigure}
\begin{center}
\begin{tikzpicture}[scale=0.85]
\draw[->,popmuted] (-1,0) -- (5,0) node[below right] {$x$};
\draw[->,popmuted] (0,-2) -- (0,6) node[above left] {$y$};
\draw[popblue,thick,domain=-0.5:3.2] plot (\x,{2*\x-1}) node[right] {$(D)$};
\draw[poporange,thick,domain=-0.5:5] plot (\x,{-\x+5}) node[right] {$(D')$};
\draw[popgreen,thick,domain=-0.5:1.6] plot (\x,{2*\x+3}) node[right] {$(D'')$};
\fill[popink] (2,3) circle (2pt) node[below right] {$I(2\,;3)$};
\end{tikzpicture}
\end{center}
\legende{$(D)$ et $(D')$ sont sécantes en $I(2\,;3)$ ; $(D)$ et $(D'')$ sont strictement parallèles.}
\end{popfigure}

\courssub{Méthode 6 : Modéliser une situation concrète par des équations de droites}

\begin{methode}[Modélisation par des droites]
Face à un problème concret (tarifs, trajets, abonnements) :
\begin{enumerate}
\item \textbf{Choisir un repère} et identifier ce que représentent $x$ et $y$ (préciser les unités).
\item \textbf{Traduire} chaque donnée : un point connu, un taux de variation (coefficient directeur), un tarif fixe (ordonnée à l'origine).
\item \textbf{Déterminer} les équations des droites en jeu avec les méthodes précédentes.
\item \textbf{Répondre à la question} : comparer les ordonnées, ou chercher l'intersection (valeur pour laquelle deux offres sont équivalentes).
\item \textbf{Conclure par une phrase} dans le langage de la situation, avec les unités.
\end{enumerate}
\end{methode}

\begin{exoresolu}[Comparaison de deux offres de transport]
À Douala, deux compagnies de transport proposent les tarifs suivants pour un trajet interurbain :
\begin{itemize}
\item \textbf{Compagnie A :} une carte d'abonnement de \SI{2000}{F} puis \SI{500}{F} par voyage ;
\item \textbf{Compagnie B :} \SI{700}{F} par voyage, sans abonnement.
\end{itemize}
On note $x$ le nombre de voyages effectués dans l'année et $y$ la dépense annuelle totale en francs.
\begin{enumerate}
\item Déterminer les équations réduites des droites $(D_A)$ et $(D_B)$ traduisant les dépenses avec chaque compagnie.
\item Pour quel nombre de voyages les deux offres sont-elles équivalentes ?
\item Un client prévoit 15 voyages dans l'année. Quelle compagnie doit-il choisir ?
\end{enumerate}
\tcblower
\textbf{1.} Avec la compagnie A, la dépense est $y = 500x + 2000$ : $(D_A) : y = 500x + 2000$ (coefficient directeur $500$, ordonnée à l'origine $2000$ qui représente l'abonnement fixe).

Avec la compagnie B : $(D_B) : y = 700x$.

\textbf{2.} Les deux offres sont équivalentes lorsque les dépenses sont égales :
\begin{align*}
500x + 2000 &= 700x \\
2000 &= 200x \\
x &= 10.
\end{align*}
Les deux offres coûtent le même prix pour \textbf{10 voyages} (dépense commune : $700 \times 10 = \SI{7000}{F}$).

\textbf{3.} Pour $x = 15$ :
\begin{itemize}
\item Compagnie A : $y = 500 \times 15 + 2000 = 7500 + 2000 = \SI{9500}{F}$ ;
\item Compagnie B : $y = 700 \times 15 = \SI{10500}{F}$.
\end{itemize}
Comme $9500 < 10500$, le client a intérêt à choisir la \textbf{compagnie A}, avec une économie de $\SI{10500}{F} - \SI{9500}{F} = \SI{1000}{F}$.

\textbf{Conclusion :} les offres sont équivalentes pour 10 voyages ; au-delà de 10 voyages, la compagnie A est plus avantageuse.
\end{exoresolu}

\begin{attention}[Les 5 erreurs qui coûtent des points au BEPC]
\begin{enumerate}
\item \textbf{Confondre vecteur directeur et coefficients de l'équation.} Un vecteur directeur de $(D) : ax + by + c = 0$ est $\vec{u}\begin{pmatrix} -b \\ a \end{pmatrix}$, et non $\begin{pmatrix} a \\ b \end{pmatrix}$. Écris toujours la condition de colinéarité si tu as un doute.
\item \textbf{Oublier le cas des droites verticales.} La droite $x = 3$ n'a \emph{pas} d'équation réduite de la forme $y = mx + p$ ni de coefficient directeur. Avant d'isoler $y$, vérifie que le coefficient de $y$ est non nul.
\item \textbf{Se tromper de signe dans le coefficient directeur.} Pour $(D) : ax + by + c = 0$, on a $m = -\dfrac{a}{b}$. Exemple : pour $2x - 4y + 8 = 0$, $m = -\dfrac{2}{-4} = \dfrac{1}{2}$ et non $-\dfrac{1}{2}$.
\item \textbf{Conclure « parallèles » dès que $m = m'$ sans comparer les ordonnées à l'origine.} Si $m = m'$ et $p = p'$, les droites sont \emph{confondues}, pas strictement parallèles. Compare toujours $p$ et $p'$ avant de conclure.
\item \textbf{Oublier de vérifier le point d'intersection.} Après avoir résolu $mx + p = m'x + p'$, reporte la valeur de $x$ dans une équation pour obtenir $y$, puis vérifie le couple $(x\,;y)$ dans l'\emph{autre} équation. Une intersection est un \emph{point} : donne toujours ses deux coordonnées.
\end{enumerate}
\end{attention}

\newpage

\coursec{Exercices}

\coursec{Équations de droites}

\courssub{Équation cartésienne d'une droite}

\begin{definition}[Équation cartésienne]
Dans un repère du plan, une \cle{droite} $(d)$ admet une \cle{équation cartésienne} de la forme :
\[
ax + by + c = 0 \qquad \text{avec } (a\,;\,b) \neq (0\,;\,0).
\]
Les coefficients $a$ et $b$ ne sont pas tous les deux nuls.
\end{definition}

\begin{propriete}[Vecteur normal et vecteur directeur]
Soit $(d)$ une droite d'équation cartésienne $ax + by + c = 0$.
\begin{itemize}
\item Un \cle{vecteur normal} à $(d)$ est $\vec{n}(a\,;\,b)$.
\item Un \cle{vecteur directeur} de $(d)$ est $\vec{u}(-b\,;\,a)$ ou $\vec{u}(b\,;\,-a)$ (tout vecteur orthogonal à $\vec{n}$).
\end{itemize}
\end{propriete}

\begin{propriete}[Appartenance d'un point à une droite]
Un point $M(x_M\,;\,y_M)$ appartient à la droite $(d) : ax + by + c = 0$ \textbf{si et seulement si} ses coordonnées vérifient l'équation :
\[
ax_M + by_M + c = 0.
\]
\end{propriete}

\begin{exemplebox}[Vérifier l'appartenance]
La droite $(d)$ a pour équation $2x - 3y + 1 = 0$.
Le point $A(1\,;\,1)$ appartient-il à $(d)$ ?
On calcule $2(1) - 3(1) + 1 = 0$. L'égalité est vraie, donc $A \in (d)$.
Le point $B(0\,;\,0)$ : $2(0) - 3(0) + 1 = 1 \neq 0$, donc $B \notin (d)$.
\end{exemplebox}

\courssub{Droite passant par deux points}

\begin{methode}[Équation d'une droite passant par deux points $A$ et $B$ ($A \neq B$)]
\begin{enumerate}
\item Calculer les coordonnées du vecteur $\overrightarrow{AB}(x_B-x_A\,;\,y_B-y_A)$.
\item En déduire un vecteur normal $\vec{n}$ (par exemple $\vec{n}(y_A-y_B\,;\,x_B-x_A)$ car $\overrightarrow{AB} \cdot \vec{n} = 0$).
\item Une équation cartésienne est de la forme $a x + b y + c = 0$ avec $\vec{n}(a\,;\,b)$.
\item Déterminer $c$ en utilisant le fait que $A$ (ou $B$) appartient à la droite : $c = -ax_A - by_A$.
\end{enumerate}
\end{methode}

\begin{exemplebox}[Droite $(AB)$ avec $A(2\,;\,-1)$ et $B(5\,;\,3)$]
$\overrightarrow{AB}(3\,;\,4)$. Un vecteur normal est $\vec{n}(4\,;\,-3)$.
Équation : $4x - 3y + c = 0$.
$A \in (d) \Rightarrow 4(2) - 3(-1) + c = 0 \Rightarrow 8 + 3 + c = 0 \Rightarrow c = -11$.
Équation cartésienne : $\boxed{4x - 3y - 11 = 0}$.
\end{exemplebox}

\courssub{Droite de vecteur directeur donné}

\begin{methode}[Équation d'une droite passant par $A$ et de vecteur directeur $\vec{u}(u_1\,;\,u_2)$]
\begin{enumerate}
\item Un vecteur normal $\vec{n}$ est orthogonal à $\vec{u}$ : on peut prendre $\vec{n}(u_2\,;\,-u_1)$ ou $\vec{n}(-u_2\,;\,u_1)$.
\item L'équation est $u_2 x - u_1 y + c = 0$ (ou $-u_2 x + u_1 y + c = 0$).
\item Déterminer $c$ par les coordonnées de $A$ : $c = -u_2 x_A + u_1 y_A$.
\end{enumerate}
\end{methode}

\begin{exemplebox}[Droite par $A(1\,;\,2)$, $\vec{u}(-3\,;\,4)$]
Vecteur normal $\vec{n}(4\,;\,3)$ (car $4 \times (-3) + 3 \times 4 = 0$).
Équation : $4x + 3y + c = 0$.
$A \in (d) \Rightarrow 4(1) + 3(2) + c = 0 \Rightarrow 10 + c = 0 \Rightarrow c = -10$.
Équation cartésienne : $\boxed{4x + 3y - 10 = 0}$.
\end{exemplebox}

\courssub{Équation réduite et coefficient directeur}

\begin{definition}[Équation réduite]
Une droite non parallèle à l'axe des ordonnées (c.-à-d. $b \neq 0$ dans $ax+by+c=0$) admet une \cle{équation réduite} de la forme :
\[
y = mx + p
\]
où $m = -\dfrac{a}{b}$ est le \cle{coefficient directeur} et $p = -\dfrac{c}{b}$ est l'\cle{ordonnée à l'origine} (point $I(0\,;\,p)$).
\end{definition}

\begin{attention}[Droite verticale]
Une droite parallèle à l'axe des ordonnées a pour équation $x = k$ ($k \in \mathbb{R}$).
\begin{itemize}
\item Elle n'a \textbf{pas} d'équation réduite.
\item Elle n'a \textbf{pas} de coefficient directeur.
\item Un vecteur directeur est $\vec{j}(0\,;\,1)$ (ou tout vecteur colinéaire).
\end{itemize}
\end{attention}

\begin{propriete}[Coefficient directeur et vecteur directeur]
Si la droite admet une équation réduite $y = mx + p$, alors :
\begin{itemize}
\item Un vecteur directeur est $\vec{u}(1\,;\,m)$.
\item Le coefficient directeur $m$ est le quotient $\dfrac{y_B - y_A}{x_B - x_A}$ pour deux points $A,B$ distincts de la droite.
\end{itemize}
\end{propriete}

\courssub{Positions relatives de deux droites}

\begin{propriete}[Parallélisme]
Soient $(d)$ et $(d')$ deux droites.
\begin{itemize}
\item \textbf{Par vecteurs directeurs :} $(d) \parallel (d') \iff$ leurs vecteurs directeurs sont colinéaires.
\item \textbf{Par équations cartésiennes $ax+by+c=0$ et $a'x+b'y+c'=0$ :} $(d) \parallel (d') \iff ab' - a'b = 0$ (determinant nul).
\item \textbf{Par coefficients directeurs (si existants) :} $(d) \parallel (d') \iff m = m'$.
\end{itemize}
Deux droites parallèles sont \textbf{confondues} si elles ont un point commun (même équation à un facteur près), \textbf{strictement parallèles} sinon.
\end{propriete}

\begin{propriete}[Perpendiculaire]
$(d) \perp (d') \iff$ leurs vecteurs directeurs sont orthogonaux (produit scalaire nul).
Si les coefficients directeurs $m$ et $m'$ existent : $(d) \perp (d') \iff m \times m' = -1$.
\end{propriete}

\begin{propriete}[Droites sécantes]
Deux droites non parallèles sont \cle{sécantes}. Leur point d'intersection $I$ s'obtient en résolvant le système de leurs deux équations.
\end{propriete}

\begin{aretenir}[Tableau récapitulatif des positions relatives]
\begin{center}
\begin{tabularx}{\linewidth}{|l|X|X|}\hline
\textbf{Position} & \textbf{Condition (vecteurs directeurs $\vec{u},\vec{u}'$)} & \textbf{Condition (coeff. directeurs $m,m'$)} \\ \hline
Parallèles strictes & $\vec{u} \parallel \vec{u}'$ et droites distinctes & $m = m'$ et $p \neq p'$ \\ \hline
Confondues & $\vec{u} \parallel \vec{u}'$ et un point commun & $m = m'$ et $p = p'$ \\ \hline
Sécantes & $\vec{u}$ non colinéaire à $\vec{u}'$ & $m \neq m'$ \\ \hline
Perpendiculaires & $\vec{u} \cdot \vec{u}' = 0$ & $m \times m' = -1$ \\ \hline
\end{tabularx}
\end{center}
\end{aretenir}

\courssub{Applications et modélisation}

\begin{methode}[Modélisation par une fonction affine]
Une situation où une grandeur $y$ varie proportionnellement à $x$ avec un coût initial se modélise par $y = mx + p$ :
\begin{itemize}
\item $m$ : coût unitaire / taux de variation (coefficient directeur).
\item $p$ : coût fixe / valeur initiale (ordonnée à l'origine).
\item L'intersection de deux droites $y = m_1x+p_1$ et $y = m_2x+p_2$ donne le seuil de rentabilité / point d'égalité.
\end{itemize}
\end{methode}

\begin{savaistu}[Contexte camerounais]
Au Cameroun, les tarifs de l'électricité (ENEO), de l'eau (CAMWATER) ou des forfaits internet (MTN, Orange) suivent souvent ce modèle : une partie fixe (abonnement) + une partie proportionnelle à la consommation. Comparer deux offres revient à étudier l'intersection de deux droites.
\end{savaistu}

% ============================================================
% EXERCICES
% ============================================================

\courssub{Je vérifie mes connaissances}

\begin{exercice}[QCM : une seule bonne réponse]
Pour chaque question, entoure la bonne réponse. Aucune justification n'est demandée.
\begin{enumerate}
\item Une équation cartésienne d'une droite est une équation de la forme :
\begin{enumerate}
\item $ax + by + c = 0$ avec $(a\,;\,b) \neq (0\,;\,0)$
\item $ax^2 + bx + c = 0$
\item $y = ax^2 + b$
\end{enumerate}
\item La droite d'équation $2x - 3y + 5 = 0$ admet pour vecteur directeur :
\begin{enumerate}
\item $\vec{u}(2\,;\,-3)$
\item $\vec{u}(3\,;\,2)$
\item $\vec{u}(-3\,;\,2)$
\end{enumerate}
\item Le coefficient directeur de la droite d'équation $y = -4x + 7$ est :
\begin{enumerate}
\item $7$
\item $-4$
\item $4$
\end{enumerate}
\item Deux droites de coefficients directeurs $m$ et $m'$ sont parallèles si et seulement si :
\begin{enumerate}
\item $m \times m' = -1$
\item $m = m'$
\item $m + m' = 0$
\end{enumerate}
\end{enumerate}
\end{exercice}

\begin{exercice}[Vrai ou faux, en justifiant]
Réponds par \textbf{Vrai} ou \textbf{Faux} et justifie ta réponse par un calcul ou une propriété du cours.
\begin{enumerate}
\item Le point $M(3\,;\,1)$ appartient à la droite d'équation $x - 2y - 1 = 0$.
\item La droite d'équation $y = 5$ est parallèle à l'axe des abscisses.
\item Les droites $(d) : 2x - y + 3 = 0$ et $(d') : 4x - 2y - 5 = 0$ sont sécantes.
\item Toute droite du plan admet une équation réduite de la forme $y = mx + p$.
\end{enumerate}
\end{exercice}

\begin{exercice}[Lecture d'équations]
Pour chacune des droites suivantes, donne un vecteur directeur, le coefficient directeur (lorsqu'il existe) et l'ordonnée à l'origine :
\begin{enumerate}
\item $(d_1) : 3x + y - 2 = 0$ ;
\item $(d_2) : y = -\dfrac{1}{2}x + 4$ ;
\item $(d_3) : x - 5 = 0$.
\end{enumerate}
\end{exercice}

\begin{exercice}[Calculs directs]
Dans un repère orthonormé $(O\,;\,\vec{i}\,,\,\vec{j})$, on donne les points $A(2\,;\,-1)$ et $B(5\,;\,3)$.
\begin{enumerate}
\item Calcule les coordonnées du vecteur $\overrightarrow{AB}$.
\item Détermine une équation cartésienne de la droite $(AB)$.
\item Le point $C(-1\,;\,-5)$ appartient-il à la droite $(AB)$ ? Justifie.
\end{enumerate}
\end{exercice}

\courssub{Je m'entraîne}

\begin{exercice}[Droite de vecteur directeur donné]
Dans un repère orthonormé, on considère la droite $(d)$ passant par le point $A(1\,;\,2)$ et de vecteur directeur $\vec{u}(-3\,;\,4)$.
\begin{enumerate}
\item Détermine une équation cartésienne de $(d)$.
\item Déduis-en l'équation réduite de $(d)$.
\item Précise le coefficient directeur et l'ordonnée à l'origine de $(d)$.
\item Trace $(d)$ dans un repère bien choisi.
\end{enumerate}
\end{exercice}

\begin{exercice}[Parallélisme et équation à tiroirs]
Dans un repère orthonormé, on donne la droite $(d') : 4x - 2y + 1 = 0$ et le point $A(-2\,;\,3)$.
\begin{enumerate}
\item Donne un vecteur directeur de $(d')$.
\item Justifie que la droite $(d)$ passant par $A$ et parallèle à $(d')$ admet une équation de la forme $4x - 2y + c = 0$.
\item Détermine la valeur de $c$ en utilisant les coordonnées de $A$, puis écris une équation cartésienne de $(d)$.
\item Donne l'équation réduite de $(d)$ et celle de $(d')$, puis vérifie le parallélisme à l'aide des coefficients directeurs.
\end{enumerate}
\end{exercice}

\begin{exercice}[Positions relatives par deux méthodes]
On donne les droites $(d) : 2x - 3y + 5 = 0$ et $(d') : y = \dfrac{2}{3}x - 1$.
\begin{enumerate}
\item \textbf{Première méthode (vecteurs directeurs).} Donne un vecteur directeur de chaque droite, puis montre, à l'aide du critère de colinéarité, que $(d)$ et $(d')$ sont parallèles.
\item \textbf{Deuxième méthode (coefficients directeurs).} Écris l'équation réduite de $(d)$, compare les coefficients directeurs et les ordonnées à l'origine, puis précise si $(d)$ et $(d')$ sont strictement parallèles ou confondues.
\item Conclus sur la position relative de $(d)$ et $(d')$.
\end{enumerate}
\end{exercice}

\begin{exercice}[Intersection de deux droites]
Dans un repère orthonormé, on donne les droites $(d) : x + 2y - 7 = 0$ et $(d') : 3x - y + 4 = 0$.
\begin{enumerate}
\item Justifie que $(d)$ et $(d')$ sont sécantes.
\item Détermine par le calcul les coordonnées de leur point d'intersection $I$.
\item Trace les deux droites dans un repère (unité : \SI{1}{cm}) et vérifie graphiquement le résultat obtenu.
\end{enumerate}
\end{exercice}

\courssub{Je me prépare au Probatoire}

\begin{exercice}[Location de voitures à Bafoussam]
Une société de location de voitures de Bafoussam propose à ses clients deux formules pour la location d'une voiture pendant une journée :
\begin{itemize}
\item \textbf{Formule A :} un forfait fixe de \num{10000} FCFA, plus $150$ FCFA par kilomètre parcouru ;
\item \textbf{Formule B :} un forfait fixe de \num{16000} FCFA, plus $100$ FCFA par kilomètre parcouru.
\end{itemize}
On note $x$ le nombre de kilomètres parcourus dans la journée, et on se place dans un repère orthogonal (unités : \SI{1}{cm} pour $50$ km en abscisse et \SI{1}{cm} pour \num{5000} FCFA en ordonnée).
\begin{enumerate}
\item Exprime, en fonction de $x$, le prix $y_A$ payé avec la formule A et le prix $y_B$ payé avec la formule B.
\item Justifie que chaque formule est représentée par une droite, et donne le coefficient directeur et l'ordonnée à l'origine de chacune.
\item Trace les deux droites dans le repère pour $x$ compris entre $0$ et $200$ km.
\item Détermine par le calcul le kilométrage pour lequel les deux formules coûtent le même prix. Quel est ce prix ?
\item Un commerçant de Bafoussam prévoit de parcourir $180$ km dans la journée. Rédige un conseil argumenté lui indiquant la formule la plus avantageuse, en justifiant ta démarche par un calcul.
\end{enumerate}
\end{exercice}

\begin{exercice}[Étude complète d'une configuration]
Dans un repère orthonormé $(O\,;\,\vec{i}\,,\,\vec{j})$, on donne les points $A(-2\,;\,1)$, $B(4\,;\,3)$ et $C(1\,;\,-3)$.
\begin{enumerate}
\item Place les points $A$, $B$ et $C$ dans un repère soigné.
\item Détermine une équation cartésienne de la droite $(AB)$.
\item Vérifie que le point $C$ n'appartient pas à la droite $(AB)$.
\item Détermine une équation de la droite $(\Delta)$ passant par $C$ et parallèle à $(AB)$.
\item Détermine une équation de la droite $(AC)$, puis calcule les coordonnées du point d'intersection $K$ de $(AC)$ avec la droite d'équation $y = -x + 2$.
\item Les droites $(AB)$ et $(AC)$ sont-elles perpendiculaires ? Justifie en comparant leurs coefficients directeurs.
\end{enumerate}
\end{exercice}

\begin{exercice}[Lecture graphique et synthèse]
Sur la figure ci-dessous, tracée dans un repère orthonormé, on a représenté trois droites $(d_1)$, $(d_2)$ et $(d_3)$.

\begin{popfigure}
\begin{center}
\begin{tikzpicture}[scale=0.85]
\draw[popline,->] (-1.5,0) -- (6.5,0) node[below left] {$x$};
\draw[popline,->] (0,-2.5) -- (0,5.5) node[below left] {$y$};
\foreach \x in {1,2,3,4,5,6} \draw[popline] (\x,0.06) -- (\x,-0.06) node[below] {\scriptsize $\x$};
\foreach \y in {-2,-1,1,2,3,4,5} \draw[popline] (0.06,\y) -- (-0.06,\y) node[left] {\scriptsize $\y$};
\node[below left] at (0,0) {\scriptsize $O$};
\draw[popblue,thick,domain=-1.2:6.2] plot (\x,{0.5*\x+1}) node[above left] {\scriptsize $(d_1)$};
\draw[poporange,thick,domain=-1.2:5] plot (\x,{0.5*\x-1.5}) node[below right] {\scriptsize $(d_2)$};
\draw[poppurple,thick,domain=-1.2:4.5] plot (\x,{-1.5*\x+4.5}) node[above right] {\scriptsize $(d_3)$};
\end{tikzpicture}
\end{center}
\legende{Les trois droites $(d_1)$, $(d_2)$ et $(d_3)$ dans un repère orthonormé.}
\end{popfigure}

\begin{enumerate}
\item Par lecture graphique, détermine l'équation réduite de chacune des droites $(d_1)$ et $(d_2)$ (on précisera le coefficient directeur et l'ordonnée à l'origine, lus sur la figure).
\item Justifie que $(d_1)$ et $(d_2)$ sont parallèles.
\item La droite $(d_3)$ passe par les points $E(1\,;\,3)$ et $F(3\,;\,0)$. Détermine par le calcul une équation cartésienne de $(d_3)$, puis son équation réduite.
\item Calcule les coordonnées du point d'intersection $M$ des droites $(d_1)$ et $(d_3)$, puis vérifie le résultat sur la figure.
\item Le point $P(7\,;\,2)$ appartient-il à la droite $(d_2)$ ? Justifie par un calcul.
\end{enumerate}
\end{exercice}

\newpage

\coursec{Corrigés des exercices}

\coursec{Exercices corrigés — Équations de droites}

\begin{corrige}[Exercice 1 — QCM : une seule bonne réponse]
\begin{enumerate}
\item \textbf{Réponse 1.} Une équation cartésienne d'une droite dans le plan est de la forme $ax + by + c = 0$ avec $(a\,;\,b) \neq (0\,;\,0)$. Les autres formes correspondent à une équation du second degré ou à une parabole.
\item \textbf{Réponse 2.} La droite $2x - 3y + 5 = 0$ a pour vecteur normal $\vec{n}(2\,;\,-3)$. Un vecteur directeur $\vec{u}$ est orthogonal à $\vec{n}$ : on peut prendre $\vec{u}(-b\,;\,a) = (3\,;\,2)$ ou $\vec{u}(b\,;\,-a) = (-3\,;\,-2)$. Parmi les propositions, $\vec{u}(3\,;\,2)$ convient.
\item \textbf{Réponse 3.} L'équation réduite $y = mx + p$ donne directement le coefficient directeur $m$. Ici $m = -4$.
\item \textbf{Réponse 4.} Deux droites non verticales sont parallèles si et seulement si elles ont même coefficient directeur : $m = m'$. La condition $m \times m' = -1$ caractérise le perpendiculaire.
\end{enumerate}
\end{corrige}

\begin{corrige}[Exercice 2 — Vrai ou faux, en justifiant]
\begin{enumerate}
\item \textbf{Vrai.} On substitue les coordonnées de $M(3\,;\,1)$ dans l'équation $x - 2y - 1 = 0$ : $3 - 2(1) - 1 = 3 - 2 - 1 = 0$. L'égalité est vérifiée, donc $M \in (d)$.
\item \textbf{Vrai.} L'équation $y = 5$ s'écrit $0\cdot x + 1\cdot y - 5 = 0$. Le coefficient directeur est $m = -\frac{0}{1} = 0$. Une droite de coefficient directeur nul est horizontale, donc parallèle à l'axe des abscisses.
\item \textbf{Faux.} On met les équations sous forme réduite :
$(d) : 2x - y + 3 = 0 \iff y = 2x + 3$ (coefficient directeur $m = 2$).
$(d') : 4x - 2y - 5 = 0 \iff 2y = 4x - 5 \iff y = 2x - 2,5$ (coefficient directeur $m' = 2$).
Les deux droites ont même coefficient directeur ($m=m'=2$) mais des ordonnées à l'origine différentes ($3 \neq -2,5$). Elles sont donc \textbf{parallèles strictes}, et non sécantes.
\item \textbf{Faux.} Une droite parallèle à l'axe des ordonnées (verticale) a pour équation $x = k$. Elle ne peut pas s'écrire sous la forme $y = mx + p$ car elle n'a pas de coefficient directeur (la pente est infinie). Seules les droites non verticales admettent une équation réduite.
\end{enumerate}
\end{corrige}

\begin{corrige}[Exercice 3 — Lecture d'équations]
\begin{enumerate}
\item $(d_1) : 3x + y - 2 = 0$.
\begin{itemize}
\item Vecteur normal $\vec{n}(3\,;\,1)$. Un vecteur directeur : $\vec{u}(-1\,;\,3)$ (ou $(1\,;\,-3)$).
\item Forme réduite : $y = -3x + 2$. Coefficient directeur $m = -3$. Ordonnée à l'origine $p = 2$ (point $I(0\,;\,2)$).
\end{itemize}
\item $(d_2) : y = -\dfrac{1}{2}x + 4$.
\begin{itemize}
\item Déjà sous forme réduite. Coefficient directeur $m = -\dfrac{1}{2}$. Ordonnée à l'origine $p = 4$ (point $I(0\,;\,4)$).
\item Un vecteur directeur : $\vec{u}(1\,;\,m) = \left(1\,;\,-\frac{1}{2}\right)$ ou, plus simplement, $\vec{u}(2\,;\,-1)$.
\end{itemize}
\item $(d_3) : x - 5 = 0 \iff x = 5$.
\begin{itemize}
\item Droite verticale (parallèle à l'axe des ordonnées).
\item Vecteur directeur : $\vec{u}(0\,;\,1)$ (ou tout vecteur colinéaire à $\vec{j}$).
\item \textbf{Pas de coefficient directeur} (pente infinie), \textbf{pas d'ordonnée à l'origine} (ne coupe pas l'axe des ordonnées).
\end{itemize}
\end{enumerate}
\end{corrige}

\begin{corrige}[Exercice 4 — Calculs directs]
\begin{enumerate}
\item $\overrightarrow{AB}(x_B-x_A\,;\,y_B-y_A) = (5-2\,;\,3-(-1)) = (3\,;\,4)$.
\item Un vecteur normal à $(AB)$ est orthogonal à $\overrightarrow{AB}(3\,;\,4)$. On prend $\vec{n}(4\,;\,-3)$ (car $3\times4 + 4\times(-3)=0$).
Une équation cartésienne est $4x - 3y + c = 0$.
Le point $A(2\,;\,-1)$ appartient à la droite : $4(2) - 3(-1) + c = 0 \iff 8 + 3 + c = 0 \iff c = -11$.
Équation cartésienne de $(AB)$ : $\boxed{4x - 3y - 11 = 0}$.
\item On teste l'appartenance de $C(-1\,;\,-5)$ : $4(-1) - 3(-5) - 11 = -4 + 15 - 11 = 0$.
L'égalité est vérifiée, donc \textbf{$C$ appartient à la droite $(AB)$}. (Les points $A,B,C$ sont alignés).
\end{enumerate}
\end{corrige}

\begin{corrige}[Exercice 5 — Droite de vecteur directeur donné]
\begin{enumerate}
\item La droite $(d)$ passe par $A(1\,;\,2)$ et a pour vecteur directeur $\vec{u}(-3\,;\,4)$.
Un vecteur normal $\vec{n}$ est orthogonal à $\vec{u}$ : $\vec{n}(4\,;\,3)$ (car $4\times(-3)+3\times4=0$).
Équation cartésienne : $4x + 3y + c = 0$.
$A \in (d) \Rightarrow 4(1) + 3(2) + c = 0 \Rightarrow 10 + c = 0 \Rightarrow c = -10$.
Équation cartésienne : $\boxed{4x + 3y - 10 = 0}$.
\item On isole $y$ : $3y = -4x + 10 \iff y = -\dfrac{4}{3}x + \dfrac{10}{3}$.
Équation réduite : $\boxed{y = -\frac{4}{3}x + \frac{10}{3}}$.
\item Coefficient directeur $m = -\dfrac{4}{3}$. Ordonnée à l'origine $p = \dfrac{10}{3}$ (point $I\left(0\,;\,\frac{10}{3}\right)$).
\item \textbf{Trace (indications) :} Repère orthonormé. Unité \SI{1}{cm}.
\begin{itemize}
\item Point $A(1\,;\,2)$.
\item Ordonnée à l'origine $I(0\,;\,3,33\ldots)$.
\item Abscisse à l'origine : $4x - 10 = 0 \Rightarrow x = 2,5$, point $J(2,5\,;\,0)$.
\item Tracer la droite $(IJ)$ qui passe par $A$. Vecteur directeur $\vec{u}(-3\,;\,4)$ : depuis $A$, reculer de 3 unités, monter de 4 unités pour trouver un autre point.
\end{itemize}
\end{enumerate}
\end{corrige}

\begin{corrige}[Exercice 6 — Parallélisme et équation à tiroirs]
\begin{enumerate}
\item $(d') : 4x - 2y + 1 = 0$. Vecteur normal $\vec{n}(4\,;\,-2)$. Un vecteur directeur : $\vec{u}(2\,;\,4)$ ou $\vec{u}(1\,;\,2)$.
\item Deux droites parallèles ont des vecteurs normaux colinéaires (donc mêmes coefficients $a,b$ à un facteur près). Comme $(d) \parallel (d')$, elles ont même vecteur normal $\vec{n}(4\,;\,-2)$. L'équation de $(d)$ est donc de la forme $4x - 2y + c = 0$.
\item $A(-2\,;\,3) \in (d) \Rightarrow 4(-2) - 2(3) + c = 0 \Rightarrow -8 - 6 + c = 0 \Rightarrow c = 14$.
Équation cartésienne de $(d)$ : $\boxed{4x - 2y + 14 = 0}$.
\item Forme réduite de $(d)$ : $2y = 4x + 14 \iff y = 2x + 7$. ($m=2,\; p=7$).
Forme réduite de $(d')$ : $2y = 4x + 1 \iff y = 2x + 0,5$. ($m'=2,\; p'=0,5$).
Les coefficients directeurs sont égaux ($m=m'=2$) et les ordonnées à l'origine sont différentes ($7 \neq 0,5$).
Donc $(d)$ et $(d')$ sont \textbf{parallèles strictes}.
\end{enumerate}
\end{corrige}

\begin{corrige}[Exercice 7 — Positions relatives par deux méthodes]
\begin{enumerate}
\item \textbf{Méthode 1 (Vecteurs directeurs).}
$(d) : 2x - 3y + 5 = 0$. Vecteur normal $\vec{n}(2\,;\,-3)$. Vecteur directeur $\vec{u}(3\,;\,2)$.
$(d') : y = \frac{2}{3}x - 1$. Coefficient directeur $m' = \frac{2}{3}$. Vecteur directeur $\vec{u}'(3\,;\,2)$ (ou $(1\,;\,\frac{2}{3})$).
$\vec{u}(3\,;\,2)$ et $\vec{u}'(3\,;\,2)$ sont colinéaires (ils sont même égaux).
\Rightarrow $(d) \parallel (d')$.
\item \textbf{Méthode 2 (Coefficients directeurs).}
$(d) : 2x - 3y + 5 = 0 \iff 3y = 2x + 5 \iff y = \frac{2}{3}x + \frac{5}{3}$.
Coefficient directeur $m = \frac{2}{3}$, ordonnée à l'origine $p = \frac{5}{3}$.
$(d') : y = \frac{2}{3}x - 1$. Coefficient directeur $m' = \frac{2}{3}$, ordonnée à l'origine $p' = -1$.
On a $m = m' = \frac{2}{3}$ et $p \neq p'$ ($\frac{5}{3} \neq -1$).
\Rightarrow $(d)$ et $(d')$ sont \textbf{parallèles strictes}.
\item \textbf{Conclusion :} Les droites $(d)$ et $(d')$ sont \textbf{parallèles strictes} (pas de point d'intersection).
\end{enumerate}
\end{corrige}

\begin{corrige}[Exercice 8 — Intersection de deux droites]
\begin{enumerate}
\item $(d) : x + 2y - 7 = 0$. Vecteur normal $\vec{n}(1\,;\,2)$, vecteur directeur $\vec{u}(-2\,;\,1)$.
$(d') : 3x - y + 4 = 0$. Vecteur normal $\vec{n}'(3\,;\,-1)$, vecteur directeur $\vec{u}'(1\,;\,3)$.
Déterminant (critère de colinéarité) : $\det(\vec{u},\vec{u}') = (-2)\times 3 - 1\times 1 = -6 - 1 = -7 \neq 0$.
Les vecteurs directeurs ne sont pas colinéaires, donc les droites ne sont pas parallèles.
\Rightarrow $(d)$ et $(d')$ sont \textbf{sécantes}.
\item Résolution du système :
\[
\begin{cases}
x + 2y = 7 & (1)\\
3x - y = -4 & (2)
\end{cases}
\]
De (2) : $y = 3x + 4$.
Substitution dans (1) : $x + 2(3x + 4) = 7 \iff x + 6x + 8 = 7 \iff 7x = -1 \iff x = -\dfrac{1}{7}$.
$y = 3\left(-\dfrac{1}{7}\right) + 4 = -\dfrac{3}{7} + \dfrac{28}{7} = \dfrac{25}{7}$.
Point d'intersection : $\boxed{I\left(-\dfrac{1}{7}\,;\,\dfrac{25}{7}\right)} \approx I(-0,14\,;\,3,57)$.
\item \textbf{Vérification graphique :} Dans un repère orthonormé (unité \SI{1}{cm}), tracer :
$(d)$ : passe par $(0\,;\,3,5)$ et $(7\,;\,0)$.
$(d')$ : passe par $(0\,;\,4)$ et $(-1,33\,;\,0)$.
Les droites se coupent bien au voisinage de $(-0,14\,;\,3,57)$, confirmant le calcul.
\end{enumerate}
\end{corrige}

\begin{corrige}[Exercice 9 — Location de voitures à Bafoussam]
\begin{enumerate}
\item $y_A(x) = 10000 + 150x$ \quad (prix en FCFA, $x$ en km).
$y_B(x) = 16000 + 100x$.
\item Ce sont des fonctions affines ($y = mx + p$), donc leurs représentations graphiques sont des droites.
Formule A : coefficient directeur $m_A = 150$ (coût au km), ordonnée à l'origine $p_A = 10000$ (forfait fixe).
Formule B : coefficient directeur $m_B = 100$, ordonnée à l'origine $p_B = 16000$.
\item \textbf{Trace (repère orthogonal) :}
Unités : \SI{1}{cm} pour 50 km (abscisses) $\rightarrow$ 200 km = 4 cm.
\SI{1}{cm} pour \num{5000} FCFA (ordonnées) $\rightarrow$ max $\approx$ \num{40000} FCFA = 8 cm.
Droite A : passe par $(0\,;\,10000)$ et $(200\,;\,40000)$.
Droite B : passe par $(0\,;\,16000)$ et $(200\,;\,36000)$.
(Les droites se coupent entre 100 et 150 km).
\item \textbf{Seuil d'égalité (intersection) :} $y_A = y_B$
$10000 + 150x = 16000 + 100x \iff 50x = 6000 \iff x = 120$.
Pour $x = 120$ km, le prix commun est $y = 10000 + 150 \times 120 = 10000 + 18000 = \num{28000}$ FCFA.
\item \textbf{Conseil pour 180 km :}
$y_A(180) = 10000 + 150 \times 180 = 10000 + 27000 = \num{37000}$ FCFA.
$y_B(180) = 16000 + 100 \times 180 = 16000 + 18000 = \num{34000}$ FCFA.
Comme $34000 < 37000$, la \textbf{Formule B} est plus avantageuse pour 180 km.
\textit{Argument :} Au-delà du seuil de 120 km, la formule B (forfait plus cher mais km moins cher) devient la moins coûteuse.
\end{enumerate}
\end{corrige}
\begin{aretenir}[Barème indicatif — Probatoire]
\begin{itemize}
\item Questions 1, 2 : 2 pts chacune (expression, identification $m,p$).
\item Question 3 : 3 pts (repère correct, unités respectées, droites tracées avec 2 points chacune).
\item Question 4 : 3 pts (mise en équation, résolution, calcul du prix, conclusion en km).
\item Question 5 : 3 pts (calculs des deux prix, comparaison, conseil rédigé et justifié).
\textbf{Total : 13 pts.}
\end{itemize}
\end{aretenir}
\begin{attention}[Points de vigilance]
\begin{itemize}
\item Ne pas oublier les unités (FCFA, km) dans les réponses finales.
\item Pour le tracé : respecter l'échelle demandée (50 km = 1 cm, 5000 FCFA = 1 cm), ne pas utiliser un repère orthonormé standard.
\item Justifier le choix de la formule B par l'inégalité numérique $y_B(180) < y_A(180)$, pas seulement par "c'est moins cher".
\end{itemize}
\end{attention}

\begin{corrige}[Exercice 10 — Étude complète d'une configuration]
\begin{enumerate}
\item \textbf{Placement :} Repère orthonormé. $A(-2\,;\,1)$ (quadrant II), $B(4\,;\,3)$ (quadrant I), $C(1\,;\,-3)$ (quadrant IV).
\item \textbf{Droite $(AB)$ :} $\overrightarrow{AB}(4-(-2)\,;\,3-1) = (6\,;\,2)$. Vecteur directeur $\vec{u}(3\,;\,1)$.
Vecteur normal $\vec{n}(1\,;\,-3)$ (car $3\times1+1\times(-3)=0$).
Équation : $x - 3y + c = 0$. $A(-2\,;\,1) \in (AB) \Rightarrow -2 - 3(1) + c = 0 \Rightarrow c = 5$.
$\boxed{(AB) : x - 3y + 5 = 0}$.
\item \textbf{Appartenance de $C$ :} $1 - 3(-3) + 5 = 1 + 9 + 5 = 15 \neq 0$.
$C$ n'appartient pas à $(AB)$. (Triangle $ABC$ non dégénéré).
\item \textbf{Droite $(\Delta) \parallel (AB)$ passant par $C$ :} Même vecteur normal $\vec{n}(1\,;\,-3)$.
Équation : $x - 3y + c = 0$. $C(1\,;\,-3) \in (\Delta) \Rightarrow 1 - 3(-3) + c = 0 \Rightarrow 10 + c = 0 \Rightarrow c = -10$.
$\boxed{(\Delta) : x - 3y - 10 = 0}$.
\item \textbf{Droite $(AC)$ :} $\overrightarrow{AC}(1-(-2)\,;\,-3-1) = (3\,;\,-4)$.
Vecteur normal $\vec{n}(4\,;\,3)$ (car $3\times4 + (-4)\times3 = 0$).
Équation : $4x + 3y + c = 0$. $A(-2\,;\,1) \in (AC) \Rightarrow 4(-2) + 3(1) + c = 0 \Rightarrow -5 + c = 0 \Rightarrow c = 5$.
$(AC) : 4x + 3y + 5 = 0 \iff y = -\frac{4}{3}x - \frac{5}{3}$.
\textbf{Intersection $K$ avec $y = -x + 2$ :}
Système : $\begin{cases} 4x + 3y = -5 \\ y = -x + 2 \end{cases}$
$4x + 3(-x+2) = -5 \iff 4x - 3x + 6 = -5 \iff x = -11$.
$y = -(-11) + 2 = 13$.
$\boxed{K(-11\,;\,13)}$.
\item \textbf{Perpendiculaire ?}
$(AB) : y = \frac{1}{3}x + \frac{5}{3} \Rightarrow m_{AB} = \frac{1}{3}$.
$(AC) : y = -\frac{4}{3}x - \frac{5}{3} \Rightarrow m_{AC} = -\frac{4}{3}$.
$m_{AB} \times m_{AC} = \frac{1}{3} \times \left(-\frac{4}{3}\right) = -\frac{4}{9} \neq -1$.
Les droites $(AB)$ et $(AC)$ ne sont \textbf{pas perpendiculaires}.
\end{enumerate}
\end{corrige}
\begin{aretenir}[Barème indicatif — Probatoire]
\begin{itemize}
\item Q1 : 2 pts (repère, points). Q2 : 3 pts (vecteur, normale, équation, substitution). Q3 : 2 pts (calcul, conclusion).
\item Q4 : 3 pts (parallélisme $\rightarrow$ même normale, passage par C). Q5 : 4 pts (équation AC, système, résolution, coordonnées K).
\item Q6 : 3 pts (coefficients directeurs, produit, conclusion).
\textbf{Total : 17 pts.}
\end{itemize}
\end{aretenir}
\begin{attention}[Points de vigilance]
\begin{itemize}
\item Vérifier que le vecteur normal choisi est bien orthogonal au vecteur directeur (produit scalaire nul).
\item Pour l'intersection $K$, bien résoudre le système par substitution (ou combinaison) et donner les coordonnées exactes (fractions ou entiers).
\item Pour le perpendiculaire : citer explicitement la condition $m \times m' = -1$ et montrer le calcul.
\end{itemize}
\end{attention}

\begin{corrige}[Exercice 11 — Lecture graphique et synthèse]
\begin{enumerate}
\item \textbf{Lecture graphique sur la figure fournie :}
\begin{itemize}
\item $(d_1)$ : Passe par $(0\,;\,1)$ (ordonnée à l'origine $p_1=1$) et $(2\,;\,2)$. Coefficient directeur $m_1 = \frac{2-1}{2-0} = 0,5$.
Équation réduite : $\boxed{y = 0,5x + 1}$.
\item $(d_2)$ : Passe par $(0\,;\,-1,5)$ (ordonnée à l'origine $p_2=-1,5$) et $(3\,;\,0)$. Coefficient directeur $m_2 = \frac{0-(-1,5)}{3-0} = 0,5$.
Équation réduite : $\boxed{y = 0,5x - 1,5}$.
\end{itemize}
\item $(d_1)$ et $(d_2)$ ont même coefficient directeur $m_1 = m_2 = 0,5$ et ordonnées à l'origine différentes ($1 \neq -1,5$).
\Rightarrow Elles sont \textbf{parallèles strictes}.
\item $(d_3)$ passe par $E(1\,;\,3)$ et $F(3\,;\,0)$.
$\overrightarrow{EF}(2\,;\,-3)$. Vecteur normal $\vec{n}(3\,;\,2)$.
Équation cartésienne : $3x + 2y + c = 0$. $E(1\,;\,3) \in (d_3) \Rightarrow 3(1) + 2(3) + c = 0 \Rightarrow 9 + c = 0 \Rightarrow c = -9$.
$\boxed{(d_3) : 3x + 2y - 9 = 0}$.
Forme réduite : $2y = -3x + 9 \iff \boxed{y = -1,5x + 4,5}$.
(Cela correspond à la lecture graphique : $p=4,5$, $m=-1,5$).
\item \textbf{Intersection $M$ de $(d_1)$ et $(d_3)$ :}
$\begin{cases} y = 0,5x + 1 \\ y = -1,5x + 4,5 \end{cases}$
$0,5x + 1 = -1,5x + 4,5 \iff 2x = 3,5 \iff x = 1,75 = \frac{7}{4}$.
$y = 0,5 \times 1,75 + 1 = 0,875 + 1 = 1,875 = \frac{15}{8}$.
$\boxed{M\left(\frac{7}{4}\,;\,\frac{15}{8}\right)} \approx M(1,75\,;\,1,875)$.
\textbf{Vérification graphique :} Sur la figure, l'intersection se situe bien entre $x=1$ et $x=2$, et $y$ entre $1$ et $2$, proche de $(1,75\,;\,1,875)$. Le tracé confirme le calcul.
\item \textbf{Appartenance de $P(7\,;\,2)$ à $(d_2)$ :}
Équation $(d_2) : y = 0,5x - 1,5$.
Pour $x=7$ : $y = 0,5 \times 7 - 1,5 = 3,5 - 1,5 = 2$.
Les coordonnées de $P$ vérifient l'équation.
\Rightarrow \textbf{Oui, $P$ appartient à $(d_2)$}.
\end{enumerate}
\end{corrige}
\begin{aretenir}[Barème indicatif — Probatoire]
\begin{itemize}
\item Q1 : 4 pts (2 par droite : lecture $p$ et $m$, écriture équation).
\item Q2 : 2 pts (comparaison $m$ et $p$, vocabulaire "parallèles strictes").
\item Q3 : 4 pts (vecteur $\overrightarrow{EF}$, normale, équation cartésienne, passage forme réduite).
\item Q4 : 3 pts (système, résolution, coordonnées exactes, vérification graphique mentionnée).
\item Q5 : 2 pts (substitution, conclusion).
\textbf{Total : 15 pts.}
\end{itemize}
\end{aretenir}
\begin{attention}[Points de vigilance]
\begin{itemize}
\item En lecture graphique, bien repérer les graduations (ici 1 unité = 1 cm). Préciser les points utilisés pour calculer $m$.
\item Pour l'équation cartésienne de $(d_3)$, partir des coordonnées exactes des points $E$ et $F$ donnés dans l'énoncé, pas de la lecture graphique approximative.
\item Donner les coordonnées de $M$ sous forme exacte (fractions $\frac{7}{4}, \frac{15}{8}$) de préférence, ou décimales exactes ($1,75; 1,875$).
\item Pour l'appartenance, toujours montrer le calcul de substitution.
\end{itemize}
\end{attention}

\newpage

\coursec{Fiche bilan}

\begin{popfigure}
\centering
\begin{tikzpicture}[mindmap, grow cyclic,
  every node/.style={concept, font=\small},
  root concept/.append style={concept color=popblue, font=\bfseries},
  level 1/.append style={level distance=3.4cm, sibling angle=60},
  level 1 concept/.append style={font=\footnotesize}]
  \node[root concept]{Équations de droites}
    child[concept color=poporange] { node {Équation cartésienne $ax+by+c=0$} }
    child[concept color=popgreen] { node {Droite par deux points $A$, $B$} }
    child[concept color=poppurple] { node {Vecteur directeur $\vec{u}(\alpha;\beta)$} }
    child[concept color=popgold] { node {Équation réduite $y=mx+p$} }
    child[concept color=poppink] { node {Positions relatives : sécantes, parallèles, confondues} }
    child[concept color=popblue!70] { node {Modélisation de situations concrètes} };
\end{tikzpicture}
\legende{Carte mentale de la leçon : les six idées clés autour des équations de droites.}
\end{popfigure}

\begin{bilan}
\textbf{1. Définitions clés}
\begin{itemize}
  \item Une \cle{équation cartésienne} d'une droite est une relation de la forme $ax+by+c=0$ (avec $(a;b)\neq(0;0)$) vérifiée par les coordonnées $(x;y)$ de tout point de la droite.
  \item Un \cle{vecteur directeur} d'une droite $(d)$ est tout vecteur non nul $\vec{u}$ qui a la même direction que $(d)$.
  \item L'\cle{équation réduite} d'une droite non verticale est $y=mx+p$ : $m$ est le \cle{coefficient directeur}, $p$ l'\cle{ordonnée à l'origine}.
\end{itemize}

\textbf{2. Formules à connaître par cœur}
\begin{center}
\begin{tabularx}{\linewidth}{|l|X|}
\hline
\rowcolor{popblueL} \textbf{Situation} & \textbf{Résultat} \\
\hline
Droite de vecteur directeur $\vec{u}(\alpha;\beta)$ & Équation : $-\beta x+\alpha y+c=0$ (on trouve $c$ avec un point de la droite) \\
\hline
Coefficient directeur par $A(x_A;y_A)$ et $B(x_B;y_B)$, $x_A\neq x_B$ & $m=\dfrac{y_B-y_A}{x_B-x_A}$ \\
\hline
Droite verticale ($x_A=x_B$) & Équation : $x=x_A$ (pas de coefficient directeur) \\
\hline
Droites $y=mx+p$ et $y=m'x+p'$ & Parallèles (ou confondues) $\Leftrightarrow m=m'$ ; confondues si de plus $p=p'$ ; sécantes $\Leftrightarrow m\neq m'$ \\
\hline
Point d'intersection de deux droites sécantes & On résout le système formé par leurs deux équations \\
\hline
\end{tabularx}
\end{center}

\textbf{3. Méthodes express}
\begin{itemize}
  \item \textbf{Équation par deux points $A$ et $B$ :} calculer $m=\dfrac{y_B-y_A}{x_B-x_A}$, écrire $y=mx+p$, remplacer par les coordonnées de $A$ pour trouver $p$.
  \item \textbf{Vérifier qu'un point est sur une droite :} remplacer $x$ et $y$ par ses coordonnées dans l'équation ; l'égalité doit être vraie.
  \item \textbf{Tracer une droite :} choisir deux valeurs de $x$, calculer les $y$ correspondants, placer les deux points.
  \item \textbf{Intersection de deux droites :} résoudre le système $\begin{cases} y=mx+p \\ y=m'x+p' \end{cases}$ en égalant les deux expressions de $y$.
\end{itemize}
\end{bilan}

\begin{aretenir}[Checklist avant le BEPC]
\begin{itemize}
  \item $\square$ Je sais qu'une équation cartésienne de droite s'écrit $ax+by+c=0$.
  \item $\square$ Je sais trouver l'équation réduite $y=mx+p$ d'une droite passant par deux points donnés.
  \item $\square$ Je sais calculer un coefficient directeur : $m=\dfrac{y_B-y_A}{x_B-x_A}$.
  \item $\square$ Je n'oublie pas le cas particulier de la droite verticale $x=k$ (pas de coefficient directeur).
  \item $\square$ Je sais écrire une équation de droite connaissant un vecteur directeur et un point.
  \item $\square$ Je sais vérifier si un point appartient à une droite en testant son équation.
  \item $\square$ Je sais comparer les coefficients directeurs pour décider si deux droites sont parallèles ou sécantes.
  \item $\square$ Je sais distinguer droites confondues ($m=m'$ et $p=p'$) et droites strictement parallèles ($m=m'$, $p\neq p'$).
  \item $\square$ Je sais résoudre un système de deux équations pour trouver le point d'intersection de deux droites.
  \item $\square$ Je sais tracer une droite à partir de son équation réduite en plaçant deux points.
  \item $\square$ Je sais modéliser une situation concrète (tarif, consommation) par une droite et interpréter $m$ et $p$.
  \item $\square$ Je rédige mes justifications avec des phrases complètes et je vérifie mes calculs.
\end{itemize}
\end{aretenir}

\findecours

\end{document}
